
%%FILE blueprint/src/content.tex
% In this file you should put the actual content of the blueprint.
% It will be used both by the web and the print version.
% It should *not* include the \begin{document}
%
% If you want to split the blueprint content into several files then
% the current file can be a simple sequence of \input. Otherwise It
% can start with a \section or \chapter for instance.

\input{../../.lake/build/blueprint/library/BlrPcp}

% \inputleanmodule{BlrPcp.Basic}

% \input{../../.lake/build/blueprint/library/BlrPcp}
\chapter{Linearity test over finite fields}
\input{subsections/blr-revised.tex}
%\input{subsections/blr-lower-bounds.tex}
\chapter{Gowers Hatami theorem}
\input{subsections/gowers-hatami}
% \input{subsections/ghZ2.tex}  
%the Z_q version is formalized proved in `GHFqPrime', but it does not give a tight bound beyond Z2
\input{subsections/gh_blr.tex}
\chapter{Exponential-length PCP}
\input{subsections/pcp.tex}
\input{subsections/blr-F2.tex}
\section{Linear PCP for linear equations}
\input{subsections/lpcp-lineq.tex}
\section{Linear PCP for tensor checks}
\input{subsections/lpcp-tensor.tex}
\section{Linear PCP to PCP}
\input{subsections/lpcp-to-pcp.tex}
\section{PCP for QESAT}
\input{subsections/qesat-exp-pcp.tex}

%%FILE blueprint/src/subsections/blr-F2.tex
\section{Simpler BLR verifier for \(\F_{2}\)}
\begin{definition}\label{BLR.basicVerifier}
\uses{OracleComp,randOracleSpec}
\lean{BLR.basicVerifier}\leanok 
Let \(V_{\mathrm{BLR}(\F_{2}, n)}^{f}\) be the verifier with oracle access to \(f: \F_{2}^n \to \F_{2}\) defined as follows:
\begin{enumerate}
  \item Sample \(x, y \leftarrow \F_2^n\).
  \item Accept if and only if \(f(x) + f(y) = f(x+y)\).
\end{enumerate}
\end{definition}

\begin{theorem}\label{BLR_basic_completeness_ZMod2}
\lean{BLR_basic_completeness_ZMod2}
\uses{BLR.basicVerifier, def:linear-functions} 
% TODO: it is not really def-linear-functions used in the theorem statement but BlrPcp.linearSet
\leanok
Let \(n > 0\). For every linear function \(f: \F_{2}^n \to \F_{2}\), we have
\[\Pr \left[V_{\mathrm{BLR}(\F_{2}, n)}^{f} = 1\right] = 1.\]
\end{theorem}
\begin{proof}\leanok
% TODO: uses the ZMod 2 acceptance-probability transfer and `BLR_completeness`
\end{proof}

\begin{theorem}\label{BLR_basic_soundness_ZMod2}\lean{BLR_basic_soundness_ZMod2}
  \uses{BLR.basicVerifier, def:distance-to-linearity}\leanok
% TODO: The statement actually uses `distanceToLin`. Unify `distanceToLin` and `BlrPcp.distanceToLinear`
Let \(n > 0\) and \(f: \F_{2}^n \to \F_{2}\).
\[\Pr \left[V_{\mathrm{BLR}(\F_{2}, n)}^{f} = 1\right] \leq 1 - \distanceToLin{f}.\]
\end{theorem}
\begin{proof}\leanok
% TODO: uses the ZMod 2 acceptance-probability transfer and `BLR_soundness`
\uses{thm:finite-field-blr-soundness} % TODO: check that this is actually the correct version that is used
\end{proof}


%%FILE blueprint/src/subsections/blr-lower-bounds.tex
\begin{theorem}[Counterexample to the proposed cubic lower bound]
\label{thm:blr-cubic-lower-bound-counterexample}
The inequality
\[
    \left(1-\delta(f,\lin)\right)^3
    \le
    \Pr[\BLRVerifier=1]
\]
does not hold for all functions $f:\bbF_q^n\to\bbF_q$ with the finite-field BLR verifier
from \cref{def:finite-field-blr-test}.
\end{theorem}

\begin{proof}
It suffices to give a concrete counterexample. Take $q=2$ and $n=1$, so the domain and
codomain are both $\bbF_2$. Let
\[
    f:\bbF_2\to\bbF_2,
    \qquad
    f(0)=1,\quad f(1)=1.
\]
The homogeneous linear maps $\bbF_2\to\bbF_2$ are exactly the zero map and the identity map.
The function $f$ agrees with the zero map at no point, and it agrees with the identity map
only at the point $1$. Hence
\[
    \delta(f,\lin)=\frac{1}{2},
    \qquad
    1-\delta(f,\lin)=\frac{1}{2}.
\]

For the verifier, the only possible nonzero field elements are $a=b=1$. Thus the verifier
accepts on a pair $(x,y)$ exactly when
\[
    f(x)+f(y)=f(x+y).
\]
Since $f$ is constantly $1$, the left-hand side is $1+1=0$ in $\bbF_2$, while the
right-hand side is $1$. Therefore the verifier rejects for every choice of $x,y$, and
\[
    \Pr[\BLRVerifier=1]=0.
\]
But
\[
    \left(1-\delta(f,\lin)\right)^3
    =
    \left(\frac{1}{2}\right)^3
    =
    \frac{1}{8}
    >
    0
    =
    \Pr[\BLRVerifier=1].
\]
This contradicts the proposed lower bound.
\end{proof}

\begin{theorem}[Lower bound from agreement with a linear function]
\label{thm:blr-lower-bound-from-agreement}
Let $f:\bbF_q^n\to\bbF_q$. Suppose there exists $\ell\in\lin$ such that
\[
    \Pr_{x\in\bbF_q^n}[f(x)=\ell(x)]\ge 1-\varepsilon.
\]
Then
\[
    \Pr[\BLRVerifier=1]\ge 1-3\varepsilon+2\varepsilon^2.
\]
In particular,
\[
    \Pr[\BLRVerifier=1]\ge 1-3\delta(f,\lin)+2\delta(f,\lin)^2.
\]
\end{theorem}

\begin{proof}
Let
\[
    G:=\{z\in\bbF_q^n : f(z)=\ell(z)\}
\]
be the agreement set, and let $B:=\bbF_q^n\setminus G$ be the bad set. Write
\[
    \eta:=\Pr_z[z\in B].
\]
By assumption, $\eta\le\varepsilon$.

The verifier samples $a,b\in\bbF_q^\times$ and $x,y\in\bbF_q^n$, and queries
\[
    x,\qquad y,\qquad ax+by.
\]
For fixed $a,b\in\bbF_q^\times$, the random variables $x$, $y$, and $ax+by$ are each
uniform on $\bbF_q^n$. Indeed, $x$ and $y$ are uniform by definition, and $ax+by$ is
uniform because, after fixing $a$, $b$, and $x$, the map
\[
    y\mapsto ax+by
\]
is a bijection of $\bbF_q^n$.

Therefore
\[
    \Pr[x\in B]=\Pr[y\in B]=\Pr[ax+by\in B]=\eta.
\]
Moreover the three queried points are pairwise independent. For fixed
$a,b\in\bbF_q^\times$, each of the maps
\[
    (x,y)\mapsto (x,y),\qquad
    (x,y)\mapsto (x,ax+by),\qquad
    (x,y)\mapsto (y,ax+by)
\]
is a bijection from $\bbF_q^n\times\bbF_q^n$ to itself. Hence
\[
    \Pr[x\in B,\ y\in B]
    =
    \Pr[x\in B,\ ax+by\in B]
    =
    \Pr[y\in B,\ ax+by\in B]
    =
    \eta^2.
\]
By inclusion--exclusion and the trivial bound
\[
    \Pr[x\in B,\ y\in B,\ ax+by\in B]\le \Pr[x\in B,\ y\in B]=\eta^2,
\]
we get
\[
    \Pr[x\in B\text{ or }y\in B\text{ or }ax+by\in B]
    \le
    3\eta-3\eta^2+\eta^2
    =
    3\eta-2\eta^2.
\]
On the complementary event, the verifier sees the same three values as it would see from
$\ell$:
\[
    f(x)=\ell(x),\qquad
    f(y)=\ell(y),\qquad
    f(ax+by)=\ell(ax+by).
\]
Since $\ell$ is linear,
\[
    \ell(ax+by)=a\ell(x)+b\ell(y).
\]
Thus
\[
    f(ax+by)=a f(x)+b f(y),
\]
so the verifier accepts. Hence the acceptance probability is at least the probability of
the complementary event:
\[
    \Pr[\BLRVerifier=1]\ge 1-3\eta+2\eta^2.
\]
If $\varepsilon\ge 1/2$, then $1-3\varepsilon+2\varepsilon^2\le 0$, so the desired bound is
trivial. Otherwise $\varepsilon\le 1/2$, and since $\eta\le\varepsilon$, the function
$1-3t+2t^2$ is decreasing on the interval containing both $\eta$ and $\varepsilon$. This gives
\[
    \Pr[\BLRVerifier=1]\ge 1-3\varepsilon+2\varepsilon^2.
\]

Taking $\ell$ to be a closest linear function gives the final statement with
$\varepsilon=\delta(f,\lin)$.
\end{proof}

\begin{theorem}[Boolean tightness for subspace error sets]
\label{thm:boolean-blr-lower-bound-tight-subspace}
The bound in \cref{thm:blr-lower-bound-from-agreement} is tight for the Boolean BLR test
for infinitely many values of $\varepsilon$. More precisely, let $q=2$ and let
$B\le \bbF_2^n$ be a linear subspace of density $\varepsilon=2^{-k}$. Define
\[
    f:\bbF_2^n\to\bbF_2,
    \qquad
    f(x)=
    \begin{cases}
        1 & \text{if } x\in B,\\
        0 & \text{if } x\notin B.
    \end{cases}
\]
Then $\delta(f,\lin)=\varepsilon$ and
\[
    \Pr[\BLRVerifier=1]=1-3\varepsilon+2\varepsilon^2.
\]
\end{theorem}

\begin{proof}
The zero linear function disagrees with $f$ exactly on $B$, so
$\delta(f,\lin)\le\varepsilon$. If $m:\bbF_2^n\to\bbF_2$ is any nonzero linear function,
then $m$ is balanced. Thus $m$ agrees with the indicator of $B$ on at most half of the
domain, whereas the zero linear function agrees on a $1-\varepsilon$ fraction of the domain.
For $\varepsilon\le 1/2$, this shows that the closest linear function is the zero function,
so $\delta(f,\lin)=\varepsilon$.

Since $q=2$, the verifier has $a=b=1$ and accepts iff
\[
    f(x)+f(y)=f(x+y).
\]
Because $B$ is a subspace, among the three points $x$, $y$, and $x+y$, the number lying in
$B$ is never exactly two. Therefore the verifier rejects exactly when at least one of the
three queried points lies in $B$.

The three events
\[
    x\in B,\qquad y\in B,\qquad x+y\in B
\]
each have probability $\varepsilon$, and each pair has probability $\varepsilon^2$.
Since $B$ is a subspace, the triple event also has probability $\varepsilon^2$: if
$x,y\in B$, then $x+y\in B$. Hence
\[
    \Pr[\BLRVerifier=0]
    =
    3\varepsilon-3\varepsilon^2+\varepsilon^2
    =
    3\varepsilon-2\varepsilon^2.
\]
Taking complements gives
\[
    \Pr[\BLRVerifier=1]=1-3\varepsilon+2\varepsilon^2.
\]
\end{proof}


%%FILE blueprint/src/subsections/blr-revised.tex
\section{Basic tools for the analysis of functions over finite fields}

Throughout this section, let $q=p^s$ for a prime $p$ and $s\in\bbN$ with $s\ge 1$.
All vectors are in $\bbF_q^n$, and all expectations are uniform over the indicated finite sets.
Let $\omega_p = e^{2\pi i/p}$.

\begin{definition}[Field trace]
\label{def:field-trace}
\lean{Algebra.trace}
\leanok
The field trace $\Tr: \bbF_q \to \bbF_p$ is
\[
    \Tr(z) \coloneqq \sum_{i=0}^{s-1} z^{p^i}.
\]
\end{definition}

\begin{definition}[Base additive character]
\label{def:base-character}
\lean{BlrPcp.baseChar}
\uses{def:field-trace}
\leanok
The base additive character $\psi:\bbF_q\to\bbC$ is
\[
    \psi(z) \coloneqq \omega_p^{\Tr(z)}.
\]
\end{definition}

\begin{definition}[Additive characters]
\label{def:additive-character}
\lean{BlrPcp.charFn}
\uses{def:base-character}
\leanok
For $\alpha\in\bbF_q^n$, define the additive character
$\chi_\alpha: \bbF_q^n\to\bbC$ by
\[
    \chi_\alpha(x) \coloneqq \psi(\inner{\alpha}{x})
    = \omega_p^{\Tr(\inner{\alpha}{x})},
\]
where $\inner{\alpha}{x}=\sum_{i=1}^n \alpha_i x_i\in\bbF_q$.
\end{definition}

\begin{definition}[Expectation]
\label{def:expectation}
\lean{BlrPcp.expectation}
\leanok
For a function $h:\bbF_q^n\to\bbC$, define its uniform expectation by
\[
    \bbE_{x\in\bbF_q^n}[h(x)]
    \coloneqq \frac{1}{q^n}\sum_{x\in\bbF_q^n} h(x).
\]
\end{definition}

\begin{definition}[Inner product]
\label{def:function-inner-product}
\lean{BlrPcp.fnInner}
\uses{def:expectation}
\leanok
For functions $h_1,h_2:\bbF_q^n\to\bbC$, define
\[
    \inner{h_1}{h_2}
    \coloneqq \bbE_{x\in\bbF_q^n}\left[h_1(x)\overline{h_2(x)}\right].
\]
\end{definition}

\begin{definition}[$L_2$ norm]
\label{def:function-l2-norm}
\lean{BlrPcp.fnNorm}
\uses{def:expectation,def:function-inner-product}
\leanok
For a function $h:\bbF_q^n\to\bbC$, define its $L_2$ norm by
\[
    \|h\|_2
    \coloneqq
    \left(\bbE_{x\in\bbF_q^n}\left[\abs{h(x)}^2\right]\right)^{1/2}
    =
    \sqrt{\inner{h}{h}}.
\]
\end{definition}

\begin{definition}[Fourier coefficient]
\label{def:fourier-coefficient}
\lean{BlrPcp.fourierCoeff}
\uses{def:function-inner-product,def:additive-character}
\leanok
For $h:\bbF_q^n\to\bbC$ and $\alpha\in\bbF_q^n$, define
\[
    \widehat h(\alpha)
    \coloneqq \inner{h}{\chi_\alpha}
    = \bbE_{x\in\bbF_q^n}\left[h(x)\overline{\chi_\alpha(x)}\right].
\]
\end{definition}

\begin{lemma}[Characters are orthonormal]
\label{lem:character-orthonormal}
\lean{BlrPcp.characters_orthonormal}
\uses{def:additive-character,def:function-inner-product}
\leanok
For every $\alpha,\beta\in\bbF_q^n$,
\[
    \inner{\chi_\alpha}{\chi_\beta}
    =
    \begin{cases}
        1 & \text{if } \alpha=\beta,\\
        0 & \text{if } \alpha\ne\beta.
    \end{cases}
\]
\end{lemma}

\begin{proof}\leanok
First observe that for every $\alpha,\beta\in\bbF_q^n$ and every $x\in\bbF_q^n$,
\[
    \chi_\alpha(x)\overline{\chi_\beta(x)}
    =
    \omega_p^{\Tr(\inner{\alpha}{x})}
    \omega_p^{-\Tr(\inner{\beta}{x})}
    =
    \omega_p^{\Tr(\inner{\alpha-\beta}{x})}
    =
    \chi_{\alpha-\beta}(x),
\]
where we used the $\bbF_p$-linearity of the trace map.

If $\alpha=\beta$, then $\alpha-\beta=0$, so $\chi_{\alpha-\beta}(x)=1$ for every
$x\in\bbF_q^n$. Hence
\[
    \bbE_{x\in\bbF_q^n}\left[\chi_\alpha(x)\overline{\chi_\beta(x)}\right]
    =
    \bbE_{x\in\bbF_q^n}[1]
    =
    1.
\]

Now suppose $\alpha\ne\beta$. Let $\gamma \coloneqq \alpha-\beta$. Then $\gamma\ne 0$, so there exists
an index $j$ such that $\gamma_j\ne 0$. Fix all coordinates of $x$ except $x_j$. For the fixed
remaining coordinates, the map
\[
    x_j\mapsto \inner{\gamma}{x}
\]
has the form
\[
    x_j\mapsto \gamma_j x_j + c
\]
for some constant $c\in\bbF_q$. Since $\gamma_j\ne 0$, multiplication by $\gamma_j$ is a
bijection on $\bbF_q$, and therefore the affine map $x_j\mapsto \gamma_jx_j+c$ is also a
bijection on $\bbF_q$. Thus, as $x_j$ ranges uniformly over $\bbF_q$, the value
$\inner{\gamma}{x}$ also ranges uniformly over $\bbF_q$. Averaging first over $x_j$ and then
over the remaining coordinates gives
\[
    \bbE_{x\in\bbF_q^n}\left[\chi_\gamma(x)\right]
    =
    \bbE_{t\in\bbF_q}\left[\omega_p^{\Tr(t)}\right].
\]

It remains to show that this last average is zero. The trace map
$\Tr:\bbF_q\to\bbF_p$ is a nonzero $\bbF_p$-linear map, so every fiber of $\Tr$ has the same
cardinality. Since $\abs{\bbF_q}=q$ and $\abs{\bbF_p}=p$, each fiber has cardinality $q/p$.
Therefore
\begin{align*}
    \bbE_{t\in\bbF_q}\left[\omega_p^{\Tr(t)}\right]
    &=
    \frac{1}{q}\sum_{t\in\bbF_q}\omega_p^{\Tr(t)} \\
    &=
    \frac{1}{q}\sum_{r\in\bbF_p}
        \sum_{\substack{t\in\bbF_q\\ \Tr(t)=r}}
        \omega_p^r \\
    &=
    \frac{1}{q}\sum_{r\in\bbF_p}\frac{q}{p}\omega_p^r \\
    &=
    \frac{1}{p}\sum_{r\in\bbF_p}\omega_p^r \\
    &=0,
\end{align*}
because the sum of all $p$-th roots of unity is zero. Hence, if $\alpha\ne\beta$, then
\[
    \bbE_{x\in\bbF_q^n}\left[\chi_\alpha(x)\overline{\chi_\beta(x)}\right]
    =
    0.
\]

Combining the two cases proves the stated orthogonality relation. The equivalent formulation
follows by taking $\beta=0$.
\end{proof}

\begin{lemma}[Fourier inversion]
\label{lem:fourier-inversion}
\lean{BlrPcp.fourier_inversion}
\uses{def:fourier-coefficient,def:additive-character}
\leanok
For every function $h:\bbF_q^n\to\bbC$ we have
\[
    h(x)=\sum_{\alpha\in\bbF_q^n}\widehat h(\alpha)\chi_\alpha(x)
    \quad\text{for every }x\in\bbF_q^n.
\]
\end{lemma}

\begin{proof}\uses{lem:character-orthonormal, def:function-inner-product}\leanok
By \cref{lem:character-orthonormal}, the additive characters
\[
    \{\chi_\alpha : \alpha \in \mathbb F_q^n\}
\]
are orthonormal with respect to the inner product
\[
    \langle h_1,h_2\rangle
    =
    \mathbb E_{x\in \mathbb F_q^n}
    \left[h_1(x)\overline{h_2(x)}\right].
\]
Moreover, these characters span the vector space of all complex-valued
functions on \(\mathbb F_q^n\). Indeed, there are exactly \(q^n\) characters,
and the space of functions \(\mathbb F_q^n\to\mathbb C\) has dimension \(q^n\).
Since the characters are nonzero and pairwise orthogonal, they are linearly
independent. Hence they form an orthonormal basis.

Therefore every function \(h:\mathbb F_q^n\to\mathbb C\) has a unique expansion
in this basis:
\[
    h
    =
    \sum_{\alpha\in\mathbb F_q^n}
        c_\alpha \chi_\alpha
\]
for some coefficients \(c_\alpha\in\mathbb C\). Taking the inner product of
both sides with \(\chi_\beta\), we obtain
\[
    \langle h,\chi_\beta\rangle
    =
    \sum_{\alpha\in\mathbb F_q^n}
        c_\alpha \langle \chi_\alpha,\chi_\beta\rangle .
\]
By orthonormality, all terms in the sum vanish except the term
\(\alpha=\beta\), and \(\langle \chi_\beta,\chi_\beta\rangle=1\). Hence
\[
    \langle h,\chi_\beta\rangle = c_\beta .
\]
By definition of the Fourier coefficient,
\[
    \widehat h(\beta)=\langle h,\chi_\beta\rangle .
\]
Thus \(c_\beta=\widehat h(\beta)\) for every
\(\beta\in\mathbb F_q^n\). Therefore
\[
    h
    =
    \sum_{\alpha\in\mathbb F_q^n}
        \widehat h(\alpha)\chi_\alpha .
\]
Evaluating this identity at \(x\in\mathbb F_q^n\) gives the Fourier inversion
formula
\[
    h(x)
    =
    \sum_{\alpha\in\mathbb F_q^n}
        \widehat h(\alpha)\chi_\alpha(x).
\]
\end{proof}

\begin{lemma}[Plancherel identity]
\label{lem:fourier-plancherel}
\lean{BlrPcp.fourier_plancherel}
\uses{def:function-inner-product,def:fourier-coefficient}
\leanok
For every pair of functions $f,g:\bbF_q^n\to\bbC$,
\[
    \sum_{\alpha\in\bbF_q^n}
        \widehat f(\alpha)\overline{\widehat g(\alpha)}
    =
    \inner{f}{g}.
\]
\end{lemma}

\begin{proof}\uses{lem:fourier-inversion}\leanok
By \cref{lem:fourier-inversion},
\[
    f=\sum_{\alpha\in\bbF_q^n}\widehat f(\alpha)\chi_\alpha .
\]
Taking the inner product with \(g\) and using linearity in the first argument gives
\[
    \inner{f}{g}
    =
    \sum_{\alpha\in\bbF_q^n}
        \widehat f(\alpha)\inner{\chi_\alpha}{g}.
\]
Since
\[
    \inner{\chi_\alpha}{g}
    =
    \overline{\inner{g}{\chi_\alpha}}
    =
    \overline{\widehat g(\alpha)},
\]
we obtain
\[
    \inner{f}{g}
    =
    \sum_{\alpha\in\bbF_q^n}
        \widehat f(\alpha)\overline{\widehat g(\alpha)}.
\]
\end{proof}

\begin{lemma}[Parseval identity]
\label{lem:parseval-identity}
\lean{BlrPcp.parseval_identity}
\uses{def:fourier-coefficient,def:function-l2-norm}
\leanok
For every function $h:\bbF_q^n\to\bbC$,
\[
    \sum_{\alpha\in\bbF_q^n}\abs{\widehat h(\alpha)}^2
    =
    \|h\|_2^2.
\]
\end{lemma}

\begin{proof}\uses{lem:fourier-plancherel,def:function-inner-product}\leanok
Apply \cref{lem:fourier-plancherel} with \(f=g=h\):
\[
    \sum_{\alpha\in\bbF_q^n}
        \widehat h(\alpha)\overline{\widehat h(\alpha)}
    =
    \inner{h}{h}.
\]
The left-hand side is
\[
    \sum_{\alpha\in\bbF_q^n}\abs{\widehat h(\alpha)}^2,
\]
and by the definitions of the inner product and the $L_2$ norm, the right-hand side is
\[
    \inner{h}{h}
    =
    \bbE_{x\in\bbF_q^n}\left[\abs{h(x)}^2\right]
    =
    \|h\|_2^2.
\]
\end{proof}

\begin{definition}[Nonzero field elements]
\label{def:nonzero-field-elements}
\lean{BlrPcp.nonzeroF}
\leanok
Write
\[
    \bbF_q^\times \coloneqq \{c\in\bbF_q : c\ne 0\}.
\]
\end{definition}

\begin{definition}[Fourier expansion of a finite-field-valued function]
\label{def:fourier-expansion-field-valued}
\lean{BlrPcp.phaseLift}
\uses{def:base-character,def:nonzero-field-elements}
\leanok
Let $f:\bbF_q^n\to\bbF_q$. For each $c\in\bbF_q^\times$, define the $c$-phase lift
\[
    \fexp{\ff}{c}(x) \coloneqq \omega_p^{\Tr(c f(x))}.
\]
The Fourier expansion of $f$ is the set $\{\varphi_c\}_{c\in \bbF_q^\times}$.
\end{definition}

\begin{lemma}[Character-sum indicator]
\label{lem:character-sum-indicator}
\lean{BlrPcp.character_sum_indicator_eq}
\uses{def:base-character}
\leanok
For every $u,v\in\bbF_q$,
\[
    \ind[u=v]
    = \frac{1}{q}\sum_{c\in\bbF_q}\omega_p^{\Tr(c(u-v))}.
\]
\end{lemma}
\begin{proof}\leanok
For each $z \in \bbF_q$ we define
\[
    \psi(z) \coloneqq \omega_p^{\Tr(z)}
\]
be the standard additive character of \(\bbF_q\). We prove that, for every
\(t\in \bbF_q\),
\[
    \ind[t=0]
    =
    \frac{1}{q}\sum_{c\in\bbF_q}\psi(ct).
\]
The desired statement follows by taking \(t=u-v\).

If \(t=0\), then \(ct=0\) for every \(c\in\bbF_q\), so
\[
    \frac{1}{q}\sum_{c\in\bbF_q}\psi(ct)
    =
    \frac{1}{q}\sum_{c\in\bbF_q}1
    =
    1.
\]

Now suppose \(t\neq 0\). Then multiplication by \(t\) is a bijection
\(\bbF_q\to\bbF_q\), so
\[
    \sum_{c\in\bbF_q}\psi(ct)
    =
    \sum_{z\in\bbF_q}\psi(z).
\]
We claim that the latter sum is \(0\). Let
\[
    S \coloneqq \sum_{z\in\bbF_q}\psi(z).
\]
Since \(\Tr:\bbF_q\to\bbF_p\) is a nonzero \(\bbF_p\)-linear map, it is
surjective. Hence there exists \(a\in\bbF_q\) such that \(\Tr(a)=1\).
Using the bijection \(z\mapsto z+a\), we get
\[
    S
    =
    \sum_{z\in\bbF_q}\psi(z+a)
    =
    \sum_{z\in\bbF_q}\omega_p^{\Tr(z+a)}
    =
    \sum_{z\in\bbF_q}\omega_p^{\Tr(z)}\omega_p^{\Tr(a)}
    =
    \omega_p S.
\]
Since \(\omega_p\neq 1\), this implies \(S=0\). Therefore, when
\(t\neq 0\),
\[
    \frac{1}{q}\sum_{c\in\bbF_q}\psi(ct)=0.
\]

Combining the two cases gives
\[
    \ind[t=0]
    =
    \frac{1}{q}\sum_{c\in\bbF_q}\omega_p^{\Tr(ct)}.
\]
Substituting \(t=u-v\), we obtain
\[
    \ind[u=v]
    =
    \frac{1}{q}\sum_{c\in\bbF_q}\omega_p^{\Tr(c(u-v))}.
\]
\end{proof}

\begin{definition}[Relative Hamming distance]
\label{def:relative-hamming-distance}
\lean{BlrPcp.distance}
\leanok
For $f,g:\bbF_q^n\to\bbF_q$, define
\[
    \delta(f,g)
    \coloneqq \Pr_{x\in\bbF_q^n}\left[f(x)\ne g(x)\right].
\]
\end{definition}

\section{Linear functions over finite fields}

\begin{definition}[Linear functions]
\label{def:linear-functions}
\lean{BlrPcp.linearFn, BlrPcp.IsLinear, BlrPcp.LinearSet}
\leanok
For $\alpha\in\bbF_q^n$, define the linear scalar function
$\ell_\alpha:\bbF_q^n\to\bbF_q$ by
\[
    \ell_\alpha(x) \coloneqq \inner{\alpha}{x}.
\]
A function $h:\bbF_q^n\to\bbF_q$ is linear if $h=\ell_\alpha$ for some
$\alpha\in\bbF_q^n$. The finite set of all linear scalar functions is
\[
    \lin
    \coloneqq
    \left\{\ell_\alpha \mid \alpha\in\bbF_q^n\right\}.
\]
\end{definition}

\begin{definition}[Distance to linearity]
\label{def:distance-to-linearity}
\lean{BlrPcp.distanceToLinear}
\uses{def:relative-hamming-distance,def:linear-functions}
\leanok
The distance from $f:\bbF_q^n\to\bbF_q$ to linearity is
\[
    \delta(f,\lin) \coloneqq \min_{g\in\lin}\delta(f,g).
\]
\end{definition}



\begin{lemma}[Well-separatedness of $\lin$]
\label{lem:linear-set-well-separated}
\lean{BlrPcp.linearSet_wellSeparated}
\uses{def:linear-functions,def:relative-hamming-distance}
\leanok
For any distinct $f,g \in \lin$, one has $\delta(f,g)= 1-1/|\bbF|$.
\end{lemma}
\begin{proof}\uses{def:linear-functions,def:relative-hamming-distance}\leanok
Let $a,b \in \bbF^n$ be the unique vectors such that $f(x) = \langle a, x\rangle$ and $g(x) = \langle b, x\rangle$ for each $x$.
Let $c = a-b$, which must be a non-zero vector. In particular, there exists $i \in [n]$ such that $c_i \neq 0$. Take any permutation $\sigma: [n] \rightarrow [n]$ such that $\sigma(i)=n$. Then
\begin{align*}
    \delta(f,g) & = 1- \Pr_{x \sim \bbF^n}[\langle c, x\rangle =0] \\
    & = 1 - \frac{1}{|\bbF|^{n-1}}\sum_{x_1,\dots,x_{n-1} \in \bbF}\frac{1}{|\bbF|}\sum_{x_n \in \bbF} \ind[c_i \cdot x_n = -\sum_{j \in [n]\setminus {i}} c_j \cdot x_{\sigma(j)}] \\
    & = 1 - \frac{1}{|\bbF|^{n-1}}\sum_{x_1,\dots,x_{n-1} \in \bbF}\frac{1}{|\bbF|}\sum_{x_n \in \bbF} \ind[x_n = -c_i^{-1} \cdot \sum_{j \in [n]\setminus {i}} c_j \cdot x_{\sigma(j)}] 
\end{align*}
where the last equality follows because $c_i \neq 0$. Now note that  $\ind[x_n = -c_i^{-1} \cdot \sum_{j \in [n]\setminus {i}} c_j \cdot x_{\sigma(j)}]$ evaluates to $1$ for exactly one value of $x_n$. Thus we continue from above and get
\begin{align*}
    \delta(f,g) = 1-\frac{1}{|\bbF|^{n-1}}\sum_{x_1,\dots,x_{n-1} \in \bbF}\frac{1}{|\bbF|} = 1-\frac{1}{|\bbF|} \, .
\end{align*}
\end{proof}


\begin{lemma}[Distance formula]
\label{lem:distance-formula}
\lean{BlrPcp.distance_formula_via_phase_fourier_coefficients}
\uses{def:relative-hamming-distance,def:fourier-expansion-field-valued,def:function-inner-product}
\leanok
Let $f,g:\bbF_q^n\to\bbF_q$.
Let $\{\fexp{\ff}{c}\}_{c\in\bbF_q^\times}$ and
$\{\fexp{\fg}{c}\}_{c\in\bbF_q^\times}$ be their Fourier expansions.
Then
\[
    \delta(f,g)
    = 1-\frac{1}{q}\left(1+
        \sum_{c\in\bbF_q^\times}
        \inner{\fexp{\ff}{c}}{\fexp{\fg}{c}}
    \right).
\]
\end{lemma}

\begin{proof}\uses{lem:character-sum-indicator,lem:fourier-plancherel}\leanok
By \cref{lem:character-sum-indicator},
\[
    \Pr_x[f(x)=g(x)]
    = \bbE_x\left[\frac{1}{q}\sum_{c\in\bbF_q}
        \omega_p^{\Tr(c(f(x)-g(x)))}\right].
\]
The term $c=0$ contributes $1/q$. For $c\ne 0$,
\[
    \omega_p^{\Tr(c(f(x)-g(x)))}
    = \fexp{\ff}{c}(x)\overline{\fexp{\fg}{c}(x)}.
\]
Therefore
\[
    \Pr_x[f(x)=g(x)]
    =\frac{1}{q}\left(1+\sum_{c\in\bbF_q^\times}
        \inner{\fexp{\ff}{c}}{\fexp{\fg}{c}}
    \right),
\]
which gives the first formula after using
$\delta(f,g)=1-\Pr_x[f(x)=g(x)]$.
\end{proof}

\begin{lemma}[Fourier expansion of a linear function]
\label{lem:linear-function-fourier-expansion}
\lean{BlrPcp.phaseLift_linearFn}
\uses{def:linear-functions,def:fourier-expansion-field-valued,def:additive-character}
\leanok
Fix $\alpha\in\bbF_q^n$ and let $\{\lambda_c\}_{c \in \bbF_q^\times}$ be the Fourier expansion of $\ell_\alpha$.
Then
\[
    \fexp{\lambda}{c}=\chi_{c\alpha}.
\]
\end{lemma}

\begin{proof}\leanok
For every $x\in\bbF_q^n$,
\[
    \fexp{\lambda}{c}(x)
    = \omega_p^{\Tr(c\inner{\alpha}{x})}
    = \omega_p^{\Tr(\inner{c\alpha}{x})}
    = \chi_{c\alpha}(x).
\]
\end{proof}

\begin{definition}[Linear Fourier score]
\label{def:linear-fourier-score}
\lean{BlrPcp.linearFourierScore, BlrPcp.linearFourierScoreC}
\uses{def:nonzero-field-elements,def:fourier-expansion-field-valued,def:fourier-coefficient}
\leanok
For $f:\bbF_q^n\to\bbF_q$ and $\alpha\in\bbF_q^n$, define the linear Fourier score
of $f$ at $\alpha$ by
\[
    S_\alpha(f) \coloneqq \sum_{c\in\bbF_q^\times}
        \widehat{\fexp{\ff}{c}}(c\alpha),
\]
where $\{\fexp{\ff}{c}\}_{c \in \bbF_q^\times}$ is the Fourier expansion of $f$.
\end{definition}

\begin{lemma}[Distance to linearity in Fourier form]
\label{lem:distance-to-linearity-fourier}
\lean{BlrPcp.distance_to_linearity_fourier}
\uses{def:distance-to-linearity,def:linear-fourier-score}
\leanok
Let $f:\bbF_q^n\to\bbF_q$ have Fourier expansion
$\{\fexp{\ff}{c}\}_{c\in\bbF_q^\times}$. Then
\[
    \delta(f,\lin)
    = 1-\frac{1}{q}\left(1+
        \max_{\alpha\in\bbF_q^n}
        S_\alpha(f)
    \right).
\]
Moreover, for every $\alpha\in\bbF_q^n$, the score $S_\alpha(f)$ is real and
satisfies $-1\le S_\alpha(f)\le q-1$.
\end{lemma}

\begin{proof}\uses{lem:distance-formula,lem:linear-function-fourier-expansion}\leanok
Apply \cref{lem:distance-formula} with $g=\ell_\alpha$ and use
\cref{lem:linear-function-fourier-expansion}. This gives
\[
    \delta(f,\ell_\alpha)
    =1-\frac{1}{q}\left(1+
        S_\alpha(f)
    \right).
\]
Taking the minimum over $\alpha$ gives the claimed formula for
$\delta(f,\lin)$.
For fixed $\alpha$, the same formula expresses $S_\alpha(f)$ as
$q(1-\delta(f,\ell_\alpha))-1$. Hence $S_\alpha(f)$ is real, and since
$0\le \delta(f,\ell_\alpha)\le 1$, we get $-1\le S_\alpha(f)\le q-1$.
\end{proof}

\section{BLR linearity test}

\begin{definition}[Finite-field BLR test]
\label{def:finite-field-blr-test}
\lean{BlrPcp.BLRAcceptsAt, BlrPcp.acceptanceProbabilityBLR}
\uses{def:nonzero-field-elements}
\leanok
Given oracle access to $f:\bbF_q^n\to\bbF_q$, the finite-field BLR verifier
$\BLRVerifier$ samples $a,b\in\bbF_q^\times$ and $x,y\in\bbF_q^n$ independently and uniformly at random, and accepts iff
\[
    a f(x)+b f(y)=f(ax+by).
\]
That is,
\[
    \BLRVerifier=1
    \quad\Longleftrightarrow\quad
    a f(x)+b f(y)=f(ax+by).
\]
\end{definition}

\begin{lemma}[Completeness]
\label{lem:blr-completeness}
\lean{BlrPcp.blr_completeness}
\uses{def:finite-field-blr-test,def:linear-functions}
\leanok
If $f\in\lin$, then
\[
    \Pr[\BLRVerifier=1]=1.
\]
\end{lemma}

\begin{proof}\leanok
If $f=\ell_\alpha$ for some $\alpha\in\bbF_q^n$, then for every
$a,b\in\bbF_q^\times$ and $x,y\in\bbF_q^n$,
\[
    f(ax+by)
    =\inner{\alpha}{ax+by}
    =a\inner{\alpha}{x}+b\inner{\alpha}{y}
    =a f(x)+b f(y).
\]
Thus the verifier accepts for every choice of randomness.
\end{proof}

\begin{lemma}[Exact BLR acceptance formula]
\label{lem:exact-blr-acceptance}
\lean{BlrPcp.exact_blr_acceptance_formula}
\uses{def:finite-field-blr-test,def:linear-fourier-score}
\leanok
Let $f:\bbF_q^n\to\bbF_q$.
Then
\[
    \Pr[\BLRVerifier=1]
    = \frac{1}{q}\left(1+\frac{1}{(q-1)^2}
        \sum_{\alpha\in\bbF_q^n} S_\alpha(f)^3
    \right).
\]
\end{lemma}

\begin{proof}\uses{lem:character-sum-indicator,lem:fourier-inversion,lem:character-orthonormal}\leanok
By \cref{lem:character-sum-indicator}, for fixed $a,b,x,y$,
\[
    \ind[a f(x)+b f(y)=f(ax+by)]
    =\frac{1}{q}\sum_{c\in\bbF_q}
        \omega_p^{\Tr(c(a f(x)+b f(y)-f(ax+by)))}.
\]
Taking expectation gives
\[
    \Pr[\BLRVerifier=1]
    =\frac{1}{q}+\frac{1}{q}\bbE_{a,b,x,y}
        \left[\sum_{c\in\bbF_q^\times}
        \fexp{\ff}{ca}(x)\fexp{\ff}{cb}(y)
        \fexp{\ff}{-c}(ax+by)
        \right].
\]
Expand the three complex-valued functions into their Fourier expansions:
\begin{align*}
&\bbE_{a,b,x,y}
        \left[\sum_{c\in\bbF_q^\times}
        \fexp{\ff}{ca}(x)\fexp{\ff}{cb}(y)
        \fexp{\ff}{-c}(ax+by)
        \right] \\
&=\bbE_{a,b}\left[\sum_{c\in\bbF_q^\times}
    \sum_{\alpha,\beta,\gamma\in\bbF_q^n}
    \widehat{\fexp{\ff}{ca}}(\alpha)
    \widehat{\fexp{\ff}{cb}}(\beta)
    \widehat{\fexp{\ff}{-c}}(\gamma)
    \bbE_{x,y}\left[
        \chi_\alpha(x)\chi_\beta(y)\chi_\gamma(ax+by)
    \right]
\right].
\end{align*}
Since
\[
    \chi_\gamma(ax+by)=\chi_{a\gamma}(x)\chi_{b\gamma}(y),
\]
\cref{lem:character-orthonormal} implies that the inner expectation is $1$ exactly when
\[
    \alpha+a\gamma=0
    \quad\text{and}\quad
    \beta+b\gamma=0,
\]
and is $0$ otherwise. Therefore the last display equals
\[
    \bbE_{a,b}\left[\sum_{c\in\bbF_q^\times}
    \sum_{\alpha\in\bbF_q^n}
    \widehat{\fexp{\ff}{ca}}(\alpha)
    \widehat{\fexp{\ff}{cb}}(a^{-1}b\alpha)
    \widehat{\fexp{\ff}{-c}}(-a^{-1}\alpha)
    \right].
\]
Substitute $\alpha=ca\eta$. Then this becomes
\[
    \bbE_{a,b}\left[\sum_{c\in\bbF_q^\times}
    \sum_{\eta\in\bbF_q^n}
    \widehat{\fexp{\ff}{ca}}(ca\eta)
    \widehat{\fexp{\ff}{cb}}(cb\eta)
    \widehat{\fexp{\ff}{-c}}(-c\eta)
    \right].
\]
As $a,b,c$ range over $\bbF_q^\times$, so do $ca$, $cb$, and $-c$. Hence this equals
\[
    \frac{1}{(q-1)^2}
    \sum_{\eta\in\bbF_q^n}
    \left(\sum_{d\in\bbF_q^\times}
        \widehat{\fexp{\ff}{d}}(d\eta)
    \right)^3
    =\frac{1}{(q-1)^2}\sum_{\eta\in\bbF_q^n}S_\eta(f)^3.
\]
Substituting into the expression for $\Pr[\BLRVerifier=1]$ proves the claim.
\end{proof}

\begin{lemma}[Second moment bound]
\label{lem:second-moment-bound}
\lean{BlrPcp.sum_linearFourierScore_sq_upperbound}
\uses{def:linear-fourier-score}
\leanok
For every $f:\bbF_q^n\to\bbF_q$, one has
\[
    \sum_{\alpha\in\bbF_q^n} S_\alpha(f)^2 \le (q-1)^2.
\]
\end{lemma}

\begin{proof}\uses{lem:parseval-identity}\leanok
Expanding the square gives
\[
    \sum_{\alpha\in\bbF_q^n} S_\alpha(f)^2
    =\sum_{a,b\in\bbF_q^\times}
        \sum_{\alpha\in\bbF_q^n}
        \widehat{\fexp{\ff}{a}}(a\alpha)
        \widehat{\fexp{\ff}{b}}(b\alpha).
\]
For fixed $a,b\in\bbF_q^\times$, Cauchy-Schwarz and
\cref{lem:parseval-identity} give
\begin{align*}
\abs{\sum_{\alpha\in\bbF_q^n}
        \widehat{\fexp{\ff}{a}}(a\alpha)
        \widehat{\fexp{\ff}{b}}(b\alpha)}
&\le
\left(\sum_{\alpha\in\bbF_q^n}\abs{\widehat{\fexp{\ff}{a}}(a\alpha)}^2\right)^{1/2}
\left(\sum_{\alpha\in\bbF_q^n}\abs{\widehat{\fexp{\ff}{b}}(b\alpha)}^2\right)^{1/2} \\
&=1.
\end{align*}
The equality uses that multiplication by $a$ and by $b$ permutes $\bbF_q^n$, and that
Parseval plus $\abs{\fexp{\ff}{a}(x)}=\abs{\fexp{\ff}{b}(x)}=1$ for every $x$
make both squared $L_2$ norms equal to \(1\).
Summing over the $(q-1)^2$ choices of $(a,b)$ gives the result.
\end{proof}

\begin{theorem}[Soundness of the finite-field BLR test]
\label{thm:finite-field-blr-soundness}
\lean{BlrPcp.blr_soundness}
\uses{def:finite-field-blr-test,def:distance-to-linearity}
\leanok
For every function $f:\bbF_q^n\to\bbF_q$, one has
\[
    \Pr[\BLRVerifier=1] \le 1-\delta(f,\lin).
\]
\end{theorem}

\begin{proof}\uses{lem:distance-to-linearity-fourier,lem:exact-blr-acceptance,lem:second-moment-bound}\leanok
Let $S_\alpha(f)$ be as in \cref{def:linear-fourier-score}, and let
\[
    M \coloneqq \max_{\alpha\in\bbF_q^n}S_\alpha(f).
\]
By \cref{lem:distance-to-linearity-fourier},
\[
    1-\delta(f,\lin)=\frac{1}{q}(1+M).
\]
By \cref{lem:distance-to-linearity-fourier}, each $S_\alpha(f)$ is real and
$S_\alpha(f)\le M$.
Therefore $S_\alpha(f)^3\le M S_\alpha(f)^2$ for every $\alpha$, because
$S_\alpha(f)^2\ge 0$.
Using \cref{lem:exact-blr-acceptance} and \cref{lem:second-moment-bound},
\begin{align*}
    \Pr[\BLRVerifier=1]
    &=\frac{1}{q}\left(1+\frac{1}{(q-1)^2}
        \sum_{\alpha\in\bbF_q^n}S_\alpha(f)^3\right) \\
    &\le \frac{1}{q}\left(1+\frac{M}{(q-1)^2}
        \sum_{\alpha\in\bbF_q^n}S_\alpha(f)^2\right) \\
    &\le \frac{1}{q}(1+M) \\
    &=1-\delta(f,\lin).
\end{align*}
\end{proof}

% Below is the GH-route to the same soundness results.

%%FILE blueprint/src/subsections/blr.tex
\section{Preliminaries}
\label{sec:prelim}

Let $q=p^s$ for some prime $p$ and $s\in \bbN$. Let $\f:\bbF^n\rightarrow\bbF_q$ be a function over the finite field $\bbF_q$. Let $\omega_p=e^{2\pi i/p}$ be the root of unity modolu $p$. 

\begin{definition}[Trace function in finite fields]
The trace function $\Tr:\bbF_q\rightarrow\bbF_p$ is defined as follows:
$$\Tr(x) = \sum_{i=0}^{s-1}x^{p^i}$$
\end{definition}

% \begin{remark}
% Observe that for $s=1$, meaning that $q=p$, the trace function is simply the identity function. 
% \end{remark}

\begin{definition}
The character group of $\bbF_q$, denoted \characters, contains all homomorphisms $\chi:\bbF^n\rightarrow \bbC$.
\end{definition}

\begin{fact}
\characters consists of all the functions:
$ \chi_{\alpha}(x) = \omega_p^{\Tr(\langle \alpha, x\rangle)}
 $, for all $\alpha\in\bbF^n$. 
\end{fact}

\begin{fact}
$\{\chi_\alpha\}_{\alpha\in \bbF^n}$ form an orthonormal basis for all functions $\f : \bbF^n_q \rightarrow \bbC$ under the inner-product $\langle \f,\g\rangle = \bbE_{x\in \bbF^n}[f(x)\overline{g(x)}].$
\end{fact}


\begin{definition}
Every function  $h:\bbF^n\rightarrow\bbC$ has a Fourier expansion $\hat{h}:\bbF^n\rightarrow\bbC$ given by
$$
h(x)=\sum_{\alpha\in\bbF_q}\hat{h}(\alpha)\chi_{\alpha}(x)
$$
$$
\hat{h}(\alpha)=q^{-n}\sum_{x\in\bbF_q}{h}(x)\overline{\chi_{\alpha}(x)}
$$
\end{definition}

\begin{definition}
The norm of a function  $h:\bbF^n\rightarrow\bbC$ is equal to:
\[
\norm{h}^2=\inner{h}{h}=\bbE_{x\in \bbF^n}[h(x)\overline{h(x)}]=q^{-n}\sum_{x\in\bbF^n}h(x)^2
\]
\end{definition}

\begin{corollary}
Let $h:\bbF^n\rightarrow\bbC$. Then: $\norm{h}\le1$. 
\end{corollary}
\begin{proof}
$$
\norm{h}^2= q^{-n}\sum_{x\in\bbF^n}h(x)^2\le q^{-n}\sum_{x\in\bbF^n}|h(x)|^2\le 1
$$
\end{proof}

\begin{theorem}Let $h_1,h_2:\bbF^n\rightarrow\bbC$.
\begin{itemize}
    \item Parseval's Identity: $\norm{h_1}^2 = \sum_{\alpha\in\bbF^n}\widehat{h_1}(\alpha)^2$.
    \item Plancheral's Identity: $\inner{h_1}{h_2} = \sum_{\alpha\in\bbF^n}\widehat{h_1}(\alpha)\widehat{h_2}(\alpha)$.
\end{itemize}
\end{theorem}

As there is no isomorphism between $\bbF_q$ and $\bbC$, we instead consider $q-1$ functions $\ff_c:\bbF^n\rightarrow \bbC$ for all $c\in \bbF_q^{*}$:
\begin{align}
\ff_c(x) := \omega_p^{\Tr(c\cdot f(x))} = \sum_{\alpha\in \bbF^n} \widehat{\left(\ff_c\right)}(\alpha)\chi_{\alpha}(x)
\end{align}
\begin{definition}
For any $\f:\bbF^n\rightarrow \bbF_q$ we defined the $\bbF_p$-expansion of $f$ as the set:
$
\left\{\ff_c\right\}_{c\in\bbF_q^{*}}
$.
\end{definition}

\begin{definition}
Let $\f,\g:\bbF^n\rightarrow\bbF_q$. We defined the distance between $f,g$ as:
\[
\delta(\f,\g) =\Pr_{x\in\bbF^n}\left[f(x)\ne g(x)\right] 
\]
\end{definition}

Now, we express the distance of $\f,\g$ in terms of the Fourier coefficients $\bbF_p$-expansion sets. 
\begin{claim} \label{claim: delta(f,g)}Let $f,g:\bbF^n\rightarrow\bbF_q$ and their $\bbF_p$-expansion sets be $\{\fexp{\ff}{c}\}_{c\in\bbF_q^{*}},\{\fexp{\fg}{c}\}_{c\in\bbF_q^{*}}$ respectively.
$$\delta(f,g) = 1-\frac{1}{q}\left(1+\sum_{c\in\bbF_q^{*}}\inner{\fexp{\ff}{c}}{\fexp{\fg}{c}}\right) = 1-\frac{1}{q}\left(1+\sum_{c\in\bbF_q^{*}}\sum_{\alpha\in\bbF^n}\widehat{\fexp{\ff}{c}}(\alpha)\overline{\widehat{\fexp{\fg}{c}}(\alpha)} \right)$$    
\end{claim}
\begin{definition}
Let $\lin$ be the set of linear functions form $\bbF^n$ to $\bbF_q$. For $\alpha\in\bbF^n$, we denote $\ell_{\alpha}(x) = \inner{x}{\alpha}$. Then:
\[
\lin = \left\{\ell_{\alpha}\right\}_{\alpha\in\bbF^n}
\]
\end{definition}


\begin{claim}
Let $\f:\bbF^n\rightarrow\bbF_q$ with $\bbF_p$-expansion sets be $\{\fexp{\ff}{c}\}_{c\in\bbF_q^{*}}$ and let $\ell_{\alpha}\in\lin$ for some $\alpha\in\bbF^n$, and let it's $\bbF_p$-expansion set be $\{\fexp{\lambda}{c}\}_{c\in\bbF_q^{*}}$. Then:
$$
\delta(\f,\ell_{\alpha}) = 1-\left(\frac{1}{q}\left(1+\sum_{c\in\bbF_q^{*}}\widehat{\fexp{\ff}{c}}(c\alpha)\right)\right)  
$$
\end{claim}
\begin{proof} For all $\beta\in\bbF^n$, the following is 0 if $\beta \ne c\alpha$ and 1 otherwise:
$$\widehat{\fexp{\lambda}{c}}(\beta)=q^{-n}\sum_{x\in\bbF^n}\fexp{\lambda}{c}(x)\overline{\chi_{\beta}(x) }= q^{-n}\sum_{x\in\bbF^n}\omega_p^{\Tr(c\cdot \ell_\alpha(x)) - \Tr(\langle \beta, x\rangle) }=q^{-n}\sum_{x\in\bbF^n}\omega_p^{\Tr(\langle c\alpha, x\rangle) - \Tr(\langle \beta, x\rangle) }$$ 
\end{proof}

\begin{corollary}\label{bound for sum of expansion q}For all $\alpha\in \bbF^n$, 
$-1\le \sum_{c\in\bbF_q^{*}}\widehat{\fexp{\ff}{c}}(c\alpha)\le q-1$ is a real number.
\end{corollary}
\begin{proof}
This follows because $\delta(f,\ell_{\alpha})\in[0,1]$.
\end{proof}

\begin{proof}
    Equivalent to $0\le\delta\le 1-1/q$
\end{proof}

\begin{definition}
The distance of $\f:\bbF^n\rightarrow\bbF_q$ from \lin is defined as:  $\delta(f,\lin) :=\min_{\alpha\in\bbF^n}\delta(f,\ell_\alpha)$.
\end{definition}

\begin{corollary}\label{distance from linear}
\[
\delta(f,\lin) = \min_{\alpha\in\bbF^n}\left(1-\left(\frac{1}{q}\left(1+\sum_{c\in\bbF_q^{*}}\widehat{\fexp{\ff}{c}}(c\alpha)\right)\right) \right)
\]
Equivalently,
$$\delta(f, \mathcal{LIN}) = 1- \frac{1}{q}\left(1+ \max_{\alpha\in\bbF^n}\left(\sum_{c\in\bbF_q^{*}}\widehat{\fexp{\ff}{c}}(c\alpha)\right)\right)$$
\end{corollary}



%%%%%%%%%%%%%%%%%%%%%%%%%%%%%%%%%%%%%%%%%%%%%%%%%%%%%%%%%%%%%%%%%%%%%%%%%%%%%%%%
%%%%%%%%%%%%%%%%%%%%%%%%%%%%%%%%%%%%%%%%%%%%%%%%%%%%%%%%%%%%%%%%%%%%%%%%%%%%%%%%
%%%%%%%%%%%%%%%%%%%%%%%%%%%%%%%%%%%%%%%%%%%%%%%%%%%%%%%%%%%%%%%%%%%%%%%%%%%%%%%%
\section{BLR Test over Finite Fields}
\label{sec:blr over finite fields}

\paragraph{The BLR Test}Given oracle access to a function $\f:\bbF^n\rightarrow\bbF_q$, the verifier performs the following:
\begin{itemize}
    \item Sample random $a,b\in\bbF^n\setminus\{0\}$ and sample $x,y\in\bbF^n$. 
    \item $\BLRVerifier=1 \ \iff \ a\cdot f(x)+b\cdot f(y) = f(a\cdot x+b\cdot y)$. 
\end{itemize}

\begin{theorem}\label{FieldBLR}
Let $\f:\bbF^n\rightarrow\bbF_q$ with $\bbF_p$-expansion sets be $\{\fexp{\ff}{c}\}_{c\in\bbF_q^{*}}$. 

\[
\Pr_{a,b,x,y}\left[\BLRVerifier=1 \right] = \frac{1}{q}\left(1+\frac{1}{(q-1)^2}\sum_{\alpha\in\bbF^n}\left(\sum_{a\in\bbF_q^{*}}\widehat{\fexp{\ff}{a}}(a\alpha)\right)^3\right)
\]
\end{theorem}


\begin{proof}
Let the $\bbF_p$-expansion set of $f$ be $\{\fexp{\ff}{c}\}_{c\in\bbF_q^{*}}$.
Let also $a,b\in\bbF^n\setminus\{0\}$ and $x,y\in\bbF^n$. We observe that:
\[
\BLRVerifier=1  \ \iff \ \frac{1}{q}\left(\sum_{c\in\bbF_q}\omega_p^{\Tr(c(af(x)+bf(y)-f(ax+by)}\right) = 1
\]
As the above sum is zero unless $af(x)+bf(y)-f(ax+by) = 0$. Hence:
\begin{align*}
\Pr_{a,b,x,y}\left[\BLRVerifier=1\right]&=\bbE_{a,b,x,y}\left[\frac{1}{q}\left(\sum_{c\in\bbF_q}\omega_p^{\Tr(c(af(x)+bf(y)-f(ax+by))}\right)\right]\\
&= \frac{1}{q}+\bbE_{a,b,x,y}\left[\sum_{c\in\bbF_q^{*}}\omega_p^{\Tr(c(af(x)+bf(y)-f(ax+by))}\right]
\end{align*}

Now, we consider the last term:
\begin{align*}
&\bbE_{a,b,x,y}\left[\sum_{c\in\bbF_q^{*}}\omega_p^{\Tr(c(af(x)+bf(y)-f(ax+by))}\right] \\&= \bbE_{a,b,x,y}\left[\sum_{c\in\bbF_q^{*}}\fexp{\ff}{ca}(x)\fexp{\ff}{cb}(y)\fexp{\ff}{-c}(a\cdot x+b\cdot y)\right]\\
&=\bbE_{a,b,x,y}\left[\sum_{c\in\bbF_q^{*}}\sum_{\alpha,\beta,\gamma\in\bbF^n}\widehat{\fexp{\ff}{ca}}(\alpha)\chi_{\alpha}(x)\widehat{\fexp{\ff}{cb}}(\beta)\chi_{\beta}(y)\widehat{\fexp{\ff}{-c}}(\gamma)\chi_{\gamma}(ax+by)\right]\\
&= \bbE_{a,b,x,y}\left[\sum_{c\in\bbF_q^{*}}\sum_{\alpha,\beta,\gamma\in\bbF^n}\widehat{\fexp{\ff}{ca}}(\alpha)\widehat{\fexp{\ff}{cb}}(\beta)\widehat{\fexp{\ff}{-c}}(\gamma)\chi_{\alpha}(x)\chi_{\beta}(y)\chi_{\gamma}(ax+by)\right]\\
&= \bbE_{a,b,x,y}\left[\sum_{c\in\bbF_q^{*}}\sum_{\alpha,\beta,\gamma\in\bbF^n}\widehat{\fexp{\ff}{ca}}(\alpha)\widehat{\fexp{\ff}{cb}}(\beta)\widehat{\fexp{\ff}{-c}}(\gamma)\chi_{\alpha}(x)\chi_{\beta}(y)\chi_{a\gamma}(x)\chi_{b\gamma}(y)\right]\\
&= \bbE_{a,b}\left[\sum_{c\in\bbF_q^{*}}\sum_{\alpha,\beta,\gamma\in\bbF^n}\widehat{\fexp{\ff}{ca}}(\alpha)\widehat{\fexp{\ff}{cb}}(\beta)\widehat{\fexp{\ff}{-c}}(\gamma)\bbE_{x,y}\left[\chi_{\alpha}(x)\chi_{\beta}(y)\chi_{a\gamma}(x)\chi_{b\gamma}(y)\right]\right]
\end{align*}

Now, we consider the following expression:
\begin{align*}
    \bbE_{x,y}\left[\chi_{\alpha}(x)\chi_{\beta}(y)\chi_{a\gamma}(x)\chi_{b\gamma}(y)\right] = \bbE_{x,y}\left[\omega_p^{\Tr(\inner{\alpha}{x}+\inner{\beta}{y}+ \inner{\gamma}{ax}+\inner{\gamma}{by})}\right]
\end{align*}
Observe that the above is zero, unless $\gamma = -a^{-1}\alpha$ and $\beta = a^{-1}b\alpha$. Thus:
\begingroup
\allowdisplaybreaks
\begin{align*}
&\bbE_{a,b,x,y}\left[\sum_{c\in\bbF_q^{*}}\omega_p^{\Tr(c(af(x)+bf(y)-f(ax+by))}\right]\\
&=
\bbE_{a,b}\left[\sum_{c\in\bbF_q^{*}}\sum_{\alpha,\beta,\gamma\in\bbF^n}\widehat{\fexp{\ff}{ca}}(\alpha)\widehat{\fexp{\ff}{cb}}(\beta)\widehat{\fexp{\ff}{-c}}(\gamma)\bbE_{x,y}\left[\chi_{\alpha}(x)\chi_{\beta}(y)\chi_{a\gamma}(x)\chi_{b\gamma}(y)\right]\right] \\
&= 
\bbE_{a,b}\left[\sum_{c\in\bbF_q^{*}}\sum_{\alpha\in\bbF^n}\widehat{\fexp{\ff}{ca}}(\alpha)\widehat{\fexp{\ff}{cb}}(a^{-1}b\alpha)\widehat{\fexp{\ff}{-c}}(-a^{-1}\alpha)\right]\\
&=
\bbE_{a,b}\left[\sum_{c\in\bbF_q^{*}}\sum_{\alpha\in\bbF^n}\widehat{\fexp{\ff}{ca}}(ca\alpha)\widehat{\fexp{\ff}{cb}}(cb\alpha)\widehat{\fexp{\ff}{-c}}(-c\alpha)\right]\\
&= \bbE_{a,b}\left[\sum_{c\in\bbF_q^{*}}\sum_{\alpha\in\bbF^n}\widehat{\fexp{\ff}{a}}(a\alpha)\widehat{\fexp{\ff}{b}}(b\alpha)\widehat{\fexp{\ff}{c}}(c\alpha)\right]\\
&= \frac{1}{(q-1)^2}\sum_{a,b,c\in\bbF_q^{*}}\sum_{\alpha\in\bbF^n}\widehat{\fexp{\ff}{a}}(a\alpha)\widehat{\fexp{\ff}{b}}(b\alpha)\widehat{\fexp{\ff}{c}}(c\alpha)\\
&= \frac{1}{(q-1)^2}\sum_{\alpha\in\bbF^n}\left(\sum_{a\in\bbF_q^{*}}\widehat{\fexp{\ff}{a}}(a\alpha)\right)^3
\end{align*}
\endgroup

\end{proof}


\begin{theorem}
Let $\f:\bbF^n\rightarrow\bbF_q$. 
\[
(1-\delta(f,\lin))^3 \le \Pr_{a,b,x,y}\left[\BLRVerifier=1\right] \le 1-\delta(f,\lin)
\]
\end{theorem}
\begin{proof}
For the left hand-side:
\begin{align*}
\Pr\left[\BLRVerifier=1 \right] &= \frac{1}{q}\left(1+\frac{1}{(q-1)^2}\sum_{\alpha\in\bbF^n}\left(\sum_{a\in\bbF_q^{*}}\widehat{\fexp{\ff}{a}}(a\alpha)\right)^3\right)\\
&\le  \frac{1}{q}\left(1+\frac{1}{(q-1)^2}\max_{\alpha\in\bbF^n}\left(\sum_{a\in\bbF_q^{*}}\widehat{\fexp{\ff}{a}}(a\alpha)\right)\sum_{\alpha\in\bbF^n}\left(\sum_{a\in\bbF_q^{*}}\widehat{\fexp{\ff}{a}}(a\alpha)\right)^2\right)\\
&=  \frac{1}{q}\left(1+\frac{1}{(q-1)^2}\max_{\alpha\in\bbF^n}\left(\sum_{a\in\bbF_q^{*}}\widehat{\fexp{\ff}{a}}(a\alpha)\right)\sum_{\alpha\in\bbF^n}\left(\sum_{a,b\in\bbF_q^{*}}\widehat{\fexp{\ff}{a}}(a\alpha)\widehat{\fexp{\ff}{b}}(b\alpha)\right)\right)\\
&=  \frac{1}{q}\left(1+\frac{1}{(q-1)^2}\max_{\alpha\in\bbF^n}\left(\sum_{a\in\bbF_q^{*}}\widehat{\fexp{\ff}{a}}(a\alpha)\right)\sum_{a,b\in\bbF_q^{*}}\left(\sum_{\alpha\in\bbF^n}\widehat{\fexp{\ff}{a}}(a\alpha)\widehat{\fexp{\ff}{b}}(b\alpha)\right)\right)
\end{align*}
Consider the following, due to Plancherel's identity and the fact that $\fexp{\ff}{a}$ map to the roots of unity modulo $p$:
\begin{align*}
\sum_{\alpha\in\bbF^n}\widehat{\fexp{\ff}{a}}(a\alpha)\widehat{\fexp{\ff}{b}}(b\alpha) 
= \bbE_{\alpha}\left[{\fexp{\ff}{a}(a\alpha)\fexp{\fg}{b}(b\alpha)}\right] 
\le \bbE_{\alpha}\left[{\abs{\fexp{\ff}{a}(a\alpha)\fexp{\fg}{b}(b\alpha)}}\right]\le 1
\end{align*}

Hence, it follows that:
\[
 \Pr_{a,b,x,y}\left[\BLRVerifier=1\right] \le   \frac{1}{q}\left(1+\max_{\alpha\in\bbF^n}\left(\sum_{a\in\bbF_q^{*}}\widehat{\fexp{\ff}{a}}(a\alpha)\right)\right)
= 1-\delta(f,\lin)
\]

For the left hand side of the claim:
\begin{align*}
\Pr\left[\BLRVerifier=1 \right] &= \frac{1}{q}\left(1+\frac{1}{(q-1)^2}\sum_{\alpha\in\bbF^n}\left(\sum_{a\in\bbF_q^{*}}\widehat{\fexp{\ff}{a}}(a\alpha)\right)^3\right)\\
&\ge  \frac{1}{q}\left(1+\frac{1}{(q-1)^2}\max_{\alpha\in\bbF^n}\left(\sum_{a\in\bbF_q^{*}}\widehat{\fexp{\ff}{a}}(a\alpha)\right)^3\right)\\
&\ge  \left(\frac{1}{q}\left(1+\max_{\alpha\in\bbF^n}\left(\sum_{a\in\bbF_q^{*}}\widehat{\fexp{\ff}{a}}(a\alpha)\right)\right)\right)^3\\
&=  \left(1-\delta(f,\lin)\right)^3
\end{align*}
where the last inequality follows from the following: for every $q\geq 2, A\ge \frac{1-|\bbF|^2}{2q-1}$:
\[
\left(\frac{1}{q}\left(1+A\right)\right)^3\le
\frac{1}{q}\left(1+\frac{1}{(q-1)^2}A^3\right)
\]
(we plug $A=\max_{\alpha\in\bbF^n}\left(\sum_{a\in\bbF_q^{*}}\widehat{\fexp{\ff}{a}}(a\alpha)\right)$ in the expression above. By \cref{bound for sum of expansion q}, we observe that the above holds for any $q$).

\end{proof}




% \inputleanmodule{BlrPcp.Basic}

\section{previous chatgpt generated content}
\begin{definition}
    \label{def:ff-linear}
    A function $f : F^n \to F$ is linear if there exists
    $a \in F^n$ such that
    \[
        f(x) = \sum_i a_i x_i
    \]
    for every $x \in F^n$.
\end{definition}

\begin{definition}
    \label{def:ff-blr-test}
    The finite-field BLR test samples $x,y \in F^n$ and
  $a,b \in F \setminus \{0\}$, and accepts iff
  \[
    a f(x) + b f(y) = f(ax + by).
  \]
\end{definition}

\begin{definition}
    \label{def:ff-blr-reject-prob}
    The rejection probability of the finite-field BLR test is
  \[
    \Pr[V_{BLR}^f = 0] = 1 - \Pr[V_{BLR}^f = 1].
  \]
\end{definition}

\begin{definition}
    \label{def:ff-additive-characters}
    Write $F=\mathbb{F}_q$ where $q = p^m$ for a prime number $p$ and a positive integer $m$. For each $\alpha \in \mathbb{F}_q^n$, define
  the additive character $\chi_\alpha : \mathbb{F}_q \rightarrow \mathbb{C}$ as
  \[
    \chi_\alpha(x) := \omega_p^{\mathrm{Tr}(\langle \alpha, x\rangle)},
  \]
  where $ \omega_p = \exp(2π i/p)$ and  $\mathrm{Tr} : \mathbb{F}_q \to \mathbb{F}_p $ is the field trace $\mathrm{Tr}(a) = \sum_{j=0}^{m-1} a^{p^j}$.
\end{definition}

\begin{lemma}
    \label{lem:ff-characters-orthonormal}
The family $\{\chi_\alpha\}_{\alpha \in \mathbb{F}_q^n}$ is an
  orthonormal basis for the complex-valued functions
  $g : \mathbb{F}_q^n \to \mathbb{C}$ with respect to the inner product
  \[
    \langle g,h\rangle := \mathbb{E}_{x \in \mathbb{F}_q^n}[g(x)\overline{h(x)}].
  \]
\end{lemma}

\begin{lemma}
    \label{lem:ff-fourier-expansion}
Every function $g : \mathbb{F}_q^n \to \mathbb{C}$ has a unique
  Fourier expansion
  \[
    g(x) = \sum_{\alpha \in \mathbb{F}_q^n} \widehat g(\alpha)\chi_\alpha(x).
  \]
\end{lemma}

\begin{lemma}
    \label{lem:ff-parseval-plancherel}
    Parseval and Plancherel hold:
  \[
    \langle g,g\rangle = \sum_\alpha |\widehat g(\alpha)|^2,
    \qquad
    \langle g,h\rangle = \sum_\alpha \widehat g(\alpha)\overline{\widehat h(\alpha)}.
  \]
\end{lemma}

\begin{definition}
    \label{def:ff-fourier-set}
    For $f : \mathbb{F}_q^n \to \mathbb{F}_q$, define its Fourier set
  \[
    \Phi(f) := \{\varphi_c\}_{c \in \mathbb{F}_q^\times},
    \qquad
    \varphi_c(x) := \omega_p^{\mathrm{Tr}(c f(x))}.
  \]
\end{definition}

\begin{lemma}
    \label{lem:ff-distance-formula}
If $f,g : \mathbb{F}_q^n \to \mathbb{F}_q$ have Fourier sets
  $\{\varphi_c\}_{c \in \mathbb{F}_q^\times}$ and
  $\{\psi_c\}_{c \in \mathbb{F}_q^\times}$, then
  \[
    \Delta(f,g)
    = 1 - \frac{1}{q}\left(1 + \sum_{c \in \mathbb{F}_q^\times}
      \langle \varphi_c,\psi_c\rangle\right).
  \]
\end{lemma}

\begin{lemma}
    \label{lem:ff-distance-to-linearity}
    If $f$ has Fourier set $\{\varphi_c\}_{c \in \mathbb{F}_q^\times}$,
  then its distance to the linear functions is
  \[
    \Delta(f,\mathrm{LIN})
    = 1 - \frac{1}{q}\left(1 + \max_{\alpha \in \mathbb{F}_q^n}
      \sum_{c \in \mathbb{F}_q^\times} \widehat{\varphi_c}(c\alpha)\right).
  \]
\end{lemma}

\begin{lemma}
    \label{lem:ff-spectral-bound}
A Fourier-analytic calculation shows that the rejection probability
  of the finite-field BLR test dominates the Fourier-set expression for
  $\Delta(f,\mathrm{LIN})$
\end{lemma}

\begin{lemma}
    \label{thm:ff-blr-improved-analysis}
For every function $f : F^n \to F$, the finite-field BLR test
  satisfies
  \[
    \Pr[V_{BLR}^f = 0] \ge \Delta(f,\mathrm{LIN}).
  \]
  This is the improved finite-field analysis proved via Fourier analysis.
\end{lemma}

\begin{lemma}
    \label{cor:ff-high-acceptance-close-to-linear}
If the finite-field BLR test accepts $f$ with probability at least
  $1 - \varepsilon$, then
  \[
    \Delta(f,\mathrm{LIN}) \le \varepsilon.
  \]
\end{lemma}

\begin{lemma}
    \label{rem:q2-special-case}
    When $q = 2$, this reduces to the Boolean Fourier-analytic proof
  of the BLR theorem.
\end{lemma}

%%FILE blueprint/src/subsections/ghZ2.tex
\section{Gowers-Hatami for BLR linearity test over $\FF_2$}

In classical BLR, the mapping is from a field $\mathbb{F}_2^n$. 
First let us relate the BLR test to the notion of approximate representation.

\begin{proposition}[BLR test induces an approximate representation]\label{lem:blr-approx-rep}
Let $f:\mathbb{Z}_2^n\to \mathbb{Z}_2$ satisfy
\[
\Pr_{x,y}[f(x)+f(y)=f(x+y)] \ge 1-\epsilon.
\]
Define $F:\mathbb{Z}_2^n\to \{\pm1\}$ by $F(x)=(-1)^{f(x)}$. Then
\[
\mathbb{E}_{x,y}\big|F(x)F(y)-F(x+y)\big|^2 \le 4\epsilon,
\]
so $F$ is an $(2\epsilon,1)$-approximate representation.
\end{proposition}

\begin{proof}
    First observe that for all $x,y\in G$,
    \[
    F(x)F(y) = (-1)^{f(x)+f(y)}.
    \]
    Hence
    \[
    F(x)F(y)=F(x+y) \quad \Longleftrightarrow \quad f(x)+f(y)=f(x+y)\ (\mathrm{mod}\ 2).
    \]

    Define the indicator of BLR success:
    \[
    \mathbf{1}_{\mathrm{good}}(x,y) := \mathbf{1}[f(x)+f(y)=f(x+y)].
    \]
    Then
    \[
    \Pr_{x,y}[\mathbf{1}_{\mathrm{good}}=1] \ge 1-\epsilon.
    \]

    Now compute the squared error. Since $F(x)\in\{\pm1\}$,
    \[
    |F(x)F(y)-F(x+y)|^2 =
    \begin{cases}
    0 & \text{if } F(x)F(y)=F(x+y),\\
    4 & \text{otherwise}.
    \end{cases}
    \]

    Taking expectation over $(x,y)$,
    \begin{align*}
    \mathbb{E}_{x,y}|F(x)F(y)-F(x+y)|^2
    &= 4 \Pr_{x,y}[F(x)F(y)\neq F(x+y)] \\
    &= 4 \Pr_{x,y}[f(x)+f(y)\neq f(x+y)] \\
    &\le 4\epsilon.
    \end{align*}
\end{proof}

The Gowers-Hatami theorem can be applied to show that if a function $f: \mathbb{Z}_2^n \to \mathbb{F}_2$ passes the BLR linearity test with high probability, 
then there exists a linear function $\ell: \mathbb{F}_2^n \to \mathbb{F}_2$ that is close to $f$ in terms of Hamming distance $\delta(f, \ell)$ defined in Def.~\ref{def:relative-hamming-distance}.
To do this, we first need to apply the Gowers-Hatami theorem to the approximate representation $F(x)=(-1)^{f(x)}$ given by Proposition~\ref{lem:blr-approx-rep}.

\begin{corollary}[Gowers-Hatami Correlation for BLR]\label{cor:gh-Z2}
    Let $F: \mathbb{Z}_2^n \to U(1)$ be an $(\epsilon, \sigma)$ approximate representation. 
    Then there exists a unit vector $v$ and an exact representation $G$ such that 
    \be
    \mathbb{E}_{x \in \mathbb{Z}_2^n} \Re \langle F(x), \langle v, G(x) v\rangle \rangle \geq 1 - \epsilon.
    \ee
    
\end{corollary}

Next, we would like to conclude from the above corollary that there exists a linear function $\ell: \mathbb{Z}_2^n \to  \mathbb{Z}_2$ such that 
$\Pr_x[f(x) \neq \ell(x)] \leq \epsilon/2$. To do this, we need to use a fact that about the regular representation of $\mathbb{Z}_2^n$.


% need to prove this fact
\begin{fact}[Regular representation of $\mathbb{Z}_2^n$]
    Let $\rho(a)$ act on $L=\{f:\mathbb{Z}_2^n\to\mathbb{C}\}$ by $(\rho(a)f)(x)=f(x+a)$. 
    Then in the Fourier basis $\{\ket{S}\}$, we have
    \[
    \rho(a) = \sum_{S\subseteq[n]} \chi_S(a)\ket{S}\bra{S}.
    \]
\end{fact}

\begin{lemma}[Gowers-Hatami implies near linearity]\label{lem:gh-linearity}
    Let $F: \mathbb{Z}_2^n \to \mathbb{C}$ be an $(\epsilon, 1)$ approximate representation of $\mathbb{Z}_2^n$ from the BLR test 
    such that $F(x) = (-1)^{f(x)}$ for some function $f: \mathbb{Z}_2^n \to \mathbb{Z}_2$.
    Then there exists a linear function $\ell: \mathbb{F}_2^n \to \mathbb{F}_2$ such that the distance between $f$ and $\ell$ is bounded
    \be
    \delta(f, \ell) \leq \epsilon/2.
    \ee
\end{lemma}


\begin{proof}
    By Corollary~\ref{cor:gh-Z2}, there exists a representation $G$ and a unit vector $v$ such that
    \[
        \mathbb{E}_{x} \Re \langle F(x), \langle v, G(x)v\rangle \rangle \ge 1 - \epsilon.
    \]
    Using the diagonal form of $G(x)$ in the character basis,
    \[
        G(x) = \sum_S \chi_S(x) \ket{S}\bra{S}.
    \]

    The inner product condition can be rewritten as
    \[
       \sum_{S} p_S  \mathbb{E}_{x} \Re \langle F(x), \chi_S(x) \rangle \ge 1 - \epsilon,
    \]
        where $p_S := |\bra{v}S\rangle|^2$ and $\sum_S p_S = 1$.
    Therefore,
    \[
        \sum_S p_S \widehat{F}(S) \ge 1 - \epsilon,
    \]
    where $\widehat{F}(S):=\mathbb{E}_x[F(x)\chi_S(x)]$ is the real-valued Fourier coefficient of $F$ at $S$.

    Since $(p_S)_S$ is a probability distribution, there exists some $S^*$ such that
    \[
        \widehat{F}(S^*) \ge 1 - \epsilon.
    \]
    This means that there is a large Fourier coefficient at $S^*$.
    Define $\ell(x)=x\cdot S^*$. Then $\chi_{S^*}(x)=(-1)^{\ell(x)}$, and
    \[
        \mathbb{E}_x[F(x)\chi_{S^*}(x)] = 1 - 2\Pr_x[f(x)\neq \ell(x)].
    \]
    Hence
    \[
        1 - 2\delta(f,\ell) \ge 1 - \epsilon,
    \]
    which implies
    \[
        \delta(f,\ell) \le \epsilon/2.
    \]
\end{proof}

\begin{proposition}[Equivalence for BLR test]\label{prop:equiv_blr_z2}
    Let $ \Pr[V_{\text{BLR}}^f=0]$ denote the probability that the BLR verifier rejects a function $f: \mathbb{F}_2^n \to \mathbb{F}_2$.
    Let $\delta(f, \mathcal{LIN})$ denote the distance between $f$ and the closest linear function in $\mathcal{LIN}$.
    Then the following statements are equivalent:
    \begin{itemize}
        \item If $\Pr[V_{\text{BLR}}^f=0] \leq \epsilon$, then $\delta(f, \mathcal{LIN}) \leq \epsilon$.
        \item $\Pr[V_{\text{BLR}}^f=0] \geq \delta(f, \mathcal{LIN})$.
    \end{itemize}
\end{proposition}    


Now we provide the proof of Soundness of BLR test using the Gowers-Hatami theorem.

\begin{theorem}[Soundness of the $Z_2$ BLR]
\label{thm:z2-blr-soundness}
For every function $f:\mathbb{F}_2^n\to\mathbb{F}_2$,
\[
    \Pr[V_{\text{BLR}}^f=0] \ge \delta(f,\mathcal{LIN}).
\]
Equivalently,
\[
    \Pr[V_{\text{BLR}}^f=1] \le 1-\delta(f,\mathcal{LIN}).
\]
\end{theorem}

\begin{proof}
    Let $F: \mathbb{Z}_2^n \to \{-1,1\}$ be defined by $F(x)=(-1)^{f(x)}$, as given by Proposition~.\ref{lem:blr-approx-rep}.
    $F$ is an $(2\epsilon,1)$ approximate representation of $\mathbb{Z}_2^n$.
    Applying Gowers-Hatami theorem by Lemma.~\ref{lem:gh-linearity},
    then there exists a linear function $\ell$ such that $\delta(f, \ell) \leq \epsilon$ so that $\delta(f,\mathcal{LIN})\le \epsilon$.

    This proves the implication 
    $\Pr[V_{\text{BLR}}^f=0]\le \epsilon \Rightarrow \delta(f,\mathcal{LIN})\le \epsilon$,
    which is equivalent to
    \[
        \Pr[V_{\text{BLR}}^f=0]\ge \delta(f,\mathcal{LIN}),
    \]
     by Proposition~\ref{prop:equiv_blr_z2}.
\end{proof}

%%FILE blueprint/src/subsections/gh_blr.tex
\section{BLR linearity test over $\FF_q$ via Gowers-Hatami}\label{sec:blr-finite-field}

Assume $q = p^s$ be a prime with power $s\in \mathbb{Z}^+$, $n \ge 1$, and we set
\[
X := \F_q^n, \qquad A := \F_q^\times, \qquad \omega := e^{2\pi i/p}.
\]

% \begin{definition}[BLR verifier and distance]\label{def:blr}\lean{BlrPcp.BLRAcceptsAt, BlrPcp.distanceToLinear, BlrPcp.acceptanceProbabilityBLR}\leanok
% Given $f : X \to \F_q$, the verifier $\BLR^f$ samples $a, b \in A$ and
% $x, y \in X$ uniformly and accepts iff $af(x) + bf(y) = f(ax + by)$. Let the test rejection probability be bounded by
% \[
% \Pr_{a,b,x,y}\!\bigl[af(x) + bf(y) \ne f(ax + by)\bigr] \leq \varepsilon.
% \]
% For $\alpha \in X$, write $\ell_\alpha(x) := \inner{\alpha}{x}$, $\delta_\alpha := \Pr_x[f(x) \ne \ell_\alpha(x)]$,
% and $\delta(f, \lin) := \min_{\alpha} \delta_\alpha$.
% \end{definition}
% At the end of the section, we prove the BLR soundness using Gowers-Hatami theorem.

First notice that linearity implies two properties: 
\begin{enumerate}
    \item Additivity $f(x + y) = f(x) + f(y), \quad \forall x, y \in X$. 
    
    Note that for $q$ prime additivity is actually equivalent to linearity. However this is not the case for prime powers $q=p^s$ for $s > 1$.
    
    \item Homogeneity $f(a x) = a f(x), \quad \forall x \in X, \, a \in A$.

    For $q=2$, the two properties are redundant. 
    For $q > 2$, turns out testing homogeneity is crucial for a test to have a high soundness.
\end{enumerate}

% In the rest of the section, we discuss an encoding of both properties through a group construction by ChatGPT.


\subsection{The twisted group \texorpdfstring{$G$}{G}}\label{sec:group}

Define a group $G$ by the set $X \times A$ with the binary operation
\begin{equation}\label{eq:star}
(x, a) \star (y, b) := \bigl(b^{-1}x + a^{-1}y,\; ab\bigr).
\end{equation}

\begin{lemma}[$G$ is an abelian group]\label{lem:group}\lean{BlrPcp.Twist.instCommGroup}\leanok
The operation $\star$ in~\equref{eq:star} makes $G$ into an abelian group with identity
$e = (0, 1)$ and inverses $(x, a)^{-1} = (-a^2 x, a^{-1})$. 
\end{lemma}

\begin{proof}\leanok
\emph{Associativity.} For $(x, a), (y, b), (z, c) \in G$,
\begin{align*}
\bigl((x, a) \star (y, b)\bigr) \star (z, c)
&= (b^{-1}x + a^{-1}y,\; ab) \star (z, c) = \bigl(b^{-1}c^{-1}x + a^{-1}c^{-1}y + a^{-1}b^{-1}z,\; abc\bigr),
% &= \bigl(c^{-1}(b^{-1}x + a^{-1}y) + (ab)^{-1}z,\; abc\bigr) \\
% &= \bigl(b^{-1}c^{-1}x + a^{-1}c^{-1}y + a^{-1}b^{-1}z,\; abc\bigr),
\end{align*}
and an identical computation gives the same expression for $(x, a) \star \bigl((y, b) \star (z, c)\bigr)$.

\emph{Commutativity.}
$(y, b) \star (x, a) = (a^{-1}y + b^{-1}x,\; ba) = (b^{-1}x + a^{-1}y,\; ab) = (x, a) \star (y, b).$

\emph{Identity.} For any $(y, b) \in G$,
$(0, 1) \star (y, b) = (b^{-1}\cdot 0 + 1^{-1}\cdot y,\; 1\cdot b) = (y, b)$.

\emph{Inverse.} For any $(x, a) \in G$,
\[
(x, a) \star (-a^2 x, a^{-1}) = \bigl((a^{-1})^{-1}\cdot x + a^{-1}\cdot(-a^2 x),\; a\cdot a^{-1}\bigr)
= (ax - ax,\; 1) = (0, 1).
\]

\end{proof}

It is at first surprising to come up with $\star$ operation, but $G$ is isomorphic to the direct product group. 
\begin{remark}[Group isomorphism to product group]
The map
\be
\Psi : G \to (X, +) \times (A, \cdot), \qquad
\Psi(x, a) := (ax, a),
\ee
is a group isomorphism onto the direct product of $(X, +)$ and $(A, \cdot)$ with operation
$\star_\Psi:(x,\; a) \star_\Psi (y,\; b) = (x+y,\; ab)$. Indeed,
\[
\Psi\bigl((x, a) \star (y, b)\bigr) = \bigl(ab\cdot(b^{-1}x + a^{-1}y),\; ab\bigr) = (ax + by,\; ab)
= \Psi(x, a) \star_\Psi \Psi(y, b),
\]
and $\Psi$ is bijective with inverse $\Psi^{-1}(y, b) = (b^{-1}y, b)$.

This isomorphism is used only to \emph{motivate} the characters below. The Lean development does not
formalize $\Psi$: it equips $G$ with its group structure (\lemref{lem:group}) and defines the characters
$\chi_{\alpha,k}$ (\propref{prop:group_characaters}) directly on the twisted group, proving the homomorphism
property by direct computation rather than by pulling back through $\Psi$. Accordingly it is recorded as a
remark, with no \texttt{\textbackslash lean} link and no node in the dependency graph.
\end{remark}

The isomorphism gives us a easy way to compute its irreducible representations which forms the Pontryagin dual group (the set of irrep characters for an abelian group).
\begin{remark}[Pontryagin dual over direct product]
For finite abelian groups $G_1, G_2$ with direct product $H = G_1 \times G_2$, the Pontryagin dual (the set of irreducible characters) factors as $\hat{H} = \hat{G_1} \times \hat{G_2}$. Like the isomorphism above, this is used only to motivate the character formula and is not part of the formalized graph.
\end{remark}
\begin{proposition}[Group characters]\label{prop:group_characaters}\lean{BlrPcp.gChar, BlrPcp.gChar_mul, BlrPcp.card_gChar_index, BlrPcp.gChar_injective}\leanok
The group $(G, \star)$ has the Pontryagin dual character group 
\be
\hat{G} = \{\chi_{\alpha, k}\}_{\alpha \in X, k \in [0, \cdots, q-2]}, \quad \chi_{\alpha, k}(x,a) = \omega^{\Tr \inner{\alpha}{a x}} \psi_k(a),
\ee
where $\psi_k(a) = e^{\frac{i 2 \pi m k}{q-1}}$ and $a=g_0^m$ for a fixed multiplicative generator $g_0\in A$.
\end{proposition}
\begin{proof}\leanok
By the isomorphism $\Psi(x, a) = (ax, a)$ in the remark above, $(G, \star) \cong (X, +) \times (A, \cdot)$. First we consider the characters of the product group. Using the Pontryagin-dual factorization in the remark above, we have
\[
\widehat{(X, +) \times (A, \cdot)} = \hat{X} \times \hat{A}.
\]
The character groups of $X$ and $A$ are
\begin{align}
    \hat{X} &= \{\phi_\alpha(y) = \omega^{\Tr \inner{\alpha}{y}}\}_{\alpha\in X}, \\
     \hat{A} &= \{\psi_k(b) =e^{\frac{i 2 \pi m k}{q-1}}\}_{k \in [0, \cdots, q-2]},
\end{align}
where $b=g_0^m$ for a fixed multiplicative generator $g_0\in A$. 
Here the field trace $\Tr: \F_q \to \F_p$ is defined by $\Tr(x) = x + x^p + x^{p^2} + \cdots + x^{p^{s-1}}$ for $x \in \F_q$ (see Definition~1.1 in Chapter 1).%(see Definition~\ref{def:field-trace} in \secref{sec:blr-general-finite-fields}).
Therefore, $\widehat{(X, +) \times (A, \cdot)} = \{\Theta_{\alpha, k}(y, b) = \phi_\alpha(y) \psi_k(b)\}_{\alpha, k}$.
The character $\chi \in \hat{G}$ is given by the pull-back through the isomorphism
\be
\chi_{\alpha, k}(x, a) = \Theta_{\alpha, k}(\Psi(x, a)) = \Theta_{\alpha, k}(ax, a).
\ee
Therefore, we have $\chi_{\alpha, k}(x, a) = \phi_\alpha(ax) \psi_k(a) =\omega^{\Tr \inner{\alpha}{a x}} \psi_k(a)$.
\end{proof}
\subsection{Approximate representation and matrix lift of $f$}\label{sec:rep}

\begin{definition}[Matrix lift $F_f$]\label{def:Flift}\lean{BlrPcp.liftMatrix}\leanok
Define a unitary matrix valued lift $F_f : G \to U(\cc^{q-1})$ by the diagonal matrix
\begin{equation}\label{eq:Flift-def}
F_f(x, a) = D(a f(x)) := \sum_{c \in A} \omega^{\Tr\, c a f(x)}\ket{c}\bra{c},
\end{equation}
where $\{\ket{c}\}$ denotes the orthonormal basis vectors of a Hilbert space $\mathcal{H} \cong \cc^{q-1}$ and the inner product between operators defined by $\inner{A}{B}_\sigma=\Tr(A^* B \sigma)$ and choose $\sigma = \I/(q-1)$.
\end{definition}

% \begin{lemma}[$F_f$ is unitary-valued]\label{lem:Ff-unitary}\uses{def:Flift}\lean{BlrPcp.liftMatrix_mem_unitaryGroup}\leanok
% For every $(x, a) \in G$, the lift $F_f(x, a)$ is a unitary matrix, i.e.\ $F_f : G \to U(\cc^{q-1})$.
% \end{lemma}
We first show that the BLR condition indeed gives an approximate represenatation.

\begin{lemma}[Lift operation from BLR condition]\label{lem:rep-iff-blr}\uses{lem:group, def:Flift}\lean{BlrPcp.liftFin_isApproxRepresentation}\leanok
Let $f : X \to \F_q$ and $(x, a), (y, b) \in G$ be rejected by the BLR test with probability $\varepsilon$,
then the unitary matrix lift $F_f$ in ~\equref{eq:Flift-def} is an $(\frac{q}{q-1}\varepsilon, \frac{\I}{q-1})$ approximate representation of $(G, \star)$.
\end{lemma}

\begin{proof}\leanok
The BLR test condition gives
\[
\Pr_{c,d,x,y}\!\bigl[c f(x) + d f(y) \neq f(cx + dy)\bigr] \leq \varepsilon.
\]
First we observe that
\begin{equation}\label{eq:rep-eq-blr}
\Pr_{g, h \in G}\!\bigl[F_f(g \star h) = F_f(g)F_f(h)\bigr]
= \Pr_{c,d,x,y}\!\bigl[c f(x) + d f(y) = f(cx + dy)\bigr].
\end{equation}
Using~\equref{eq:star} and~\equref{eq:Flift-def},
\begin{align*}
F_f\bigl((x,a)\star(y,b)\bigr) &= \sum_{c\in A}\omega^{\Tr\, cab\,f(b^{-1}x + a^{-1}y)} \ket{c}\bra{c}, \\
F_f(x, a)\,F_f(y, b) &= \sum_{c\in A} \omega^{\Tr\, c(af(x) + bf(y))} \ket{c}\bra{c}.
\end{align*}
These coincide iff their exponents agree modulo $p$ for $c \in A$: 
\[
ab\, f(b^{-1}x + a^{-1}y) = af(x) + bf(y) \mod p, \quad \forall c \in A.
\]
Therefore, we have equality over $\F_p$ 
\[
ab\, f(b^{-1}x + a^{-1}y) = af(x) + bf(y).
\]
Since $a, b$ have a multiplicative inverse, substitute $c := b^{-1}$, $d := a^{-1}$, we have
\[
f(cx + d y) = c f(x) + d f(y),
\]
which proves ~\equref{eq:rep-eq-blr}.

To relate the test condition to the approximate representation, first note the equivalent definition of approximate representation
\be
\mathbb{E}_{g, h \in G} \Re \left\langle F_f(g)^{*} F_f(h), F_f\left(g^{-1} h\right)\right\rangle_{\sigma} = \mathbb{E}_{g, h \in G} \Re\left\langle F_f(g) F_f(h), F_f\left(g h\right)\right\rangle_{\sigma},
\ee
because $F_f$ is unitary so $F_f(g)^* = F_f(g^{-1}), \forall g \in G$. Then evaluate the right hand side
\begin{align}
    \mathbb{E}_{g, h \in G} \Re \left\langle F_f(g) F_f(h), F_f\left(g h\right)\right\rangle_{\sigma} &= \mathbb{E}_{g, h \in G} \Re \Tr\left[F_f^*(h) F_f^*(g) F_f\left(g h\right)\sigma\right], \nonumber \\
    & = \mathbb{E}_{(x,a), (y,b) \in G} \frac{1}{q-1}\Re \Tr\left[ \sum_{c\in A} \omega^{-\Tr\, c(af(x) + bf(y))} \omega^{\Tr\, cab\,f(b^{-1}x + a^{-1}y)} \ket{c}\bra{c}\right], \nonumber \\
    & = \mathbb{E}_{(x,a), (y,b) \in G} \frac{1}{q-1}\Re \sum_{c\in A} \omega^{\Tr\, c z},
\end{align}
where we define $z=c (ab\,f(b^{-1}x + a^{-1}y) - (af(x) + bf(y)))$. To compute the expectation, we can split to two cases
\begin{align}
\sum_{c\in A} \omega^{\Tr\, c z} = 
    \begin{cases}
        q-1, \quad z=0 \quad \text{when $f$ passes BLR}, \\
        -1, \quad z \neq 0 \quad \text{when $f$ fails BLR}.
    \end{cases}
\end{align}
The second case is because $\sum_{c\in A} \omega^{\Tr\, c z} + 1= \sum_{c\in \FF_q} \omega^{\Tr\, c z} =0$. Therefore, we have the approximate representation
\be
\mathbb{E}_{g, h \in G} \Re \left\langle F_f(g) F_f(h), F_f\left(g h\right)\right\rangle_{\sigma} \geq (1- \varepsilon) - \frac{1}{q-1} \varepsilon = 1 - \frac{q}{q-1}\varepsilon. 
\ee
Hence $F_f$ is an approximate representation with  $\epsilon= \frac{q}{q-1}\varepsilon$ as claimed.

\end{proof}
The approximate representation ``appears'' slightly worse than what we would have hoped $\epsilon > \varepsilon$, but we will see this is not an issue for achieving a high soundness. Remarkably, the twisted group operation let us recover the linearity test condition exactly.

% \begin{fact}
%     $F_f$ is an exact representation of $(G, \star)$ if and only if
% $f \in \lin$.
% \end{fact}
% \begin{proof}
%     If $F_f$ is an exact representation of $G$, then~\equref{eq:rep-eq-blr} forces
% $\varepsilon = 0$. Setting $c = d = 1$ in the BLR condition gives $f(x + y) = f(x) + f(y)$
% for all $x, y$, so $f$ is additive. Set $y = 0$ and $c = 1 \in A$: $f(x) = f(x) + d\, f(0)$ for any
% $d \in A$. Set $y = 0$ and $d=1$, then $f(0)=0$ and $f(c x) = c f(x)$. Combined with additivity, $f$ is $\F_q$-linear, so $f \in \lin$.

% Conversely, if $f = \ell_\alpha$ for some $\alpha \in X$, then so $F_f$ is a group homomorphism and an exact representation of $G$.
% \end{proof}

\subsection{Gowers-Hatami implies BLR soundness}
In this section, we first prove the key lemma relating the implication of GH theorem to the hamming distance. 
The proof outlined here correspond to the formalized version in lean. The most convenient proof still uses properties of the Fourier transform of $F_f$. However, we remark that there are other alternative proofs, which we leave as comment in the current source file.

% Finally, we provide a proof of \thmref{thm:main}.

% We provide two alternative proofs of the key lemma:
% \begin{enumerate}
%     \item The first proof in \lemref{lem:gh_linearity_1} should only use the existence of an isometry $V$ and exact representation $R$. Unfortunately, we ended up with an extra constraint about the range of $\epsilon$. 
%     \item In the second proof in \lemref{lem:gh_linearity_2}, we use an explicit isometry $V$ constructed via $F_f$ in the proof of Gowers-Hatami theorem. This proof might require a extra care in formalization of lean. However, this hybrid Fourier-GH approach gives us an unconditional high soundness of the linear BLR test.   
% \end{enumerate}

% \begin{lemma}[Gowers-Hatami to linearity (Convex character proof)] \label{lem:gh_linearity_1}
% Suppose $F_f: G \to U(\cc^{q-1})$ defined in ~\equref{eq:Flift-def} via $f : \FF_q^n \to \FF_q$ is an $(\epsilon, \sigma)$ approximate unitary representation of $(G, \star)$ defined in ~\equref{eq:star}. 
% % Furthermore, suppose $1 - \epsilon > \frac{\sqrt{q}}{q-1}$.
% Then, there exists a linear function $\ell: \FF_q^n \to \FF_q$ such that the relative Hamming distance $\delta(f, \ell)$ is bounded by:
% \be\label{eq:GH-linear-bound}
%     \delta(f, \ell) \le \frac{q-1}{q} \epsilon.
% \ee
% \end{lemma}

% \begin{proof}
% % \yan{In this specific proof below without any Fourier analysis, we also need the condition $1 - \epsilon > \frac{\sqrt{q}}{q-1}$, hence the additional assumption.}

%     By Gowers-Hatami theorem, we have an isometry $V: \mathcal{H} \to \mathcal{H}^\prime$ and an exact representation $R: G \to \End(\mathcal{H}^\prime)$, where the Hilbert space $\mathcal{H}\cong \cc^{q-1}$ and the space $R$ acts on $\mathcal{H}^\prime = L(G, \mathcal{H})$ 
%     \begin{equation}
%     \EE_{g \in G} \Re \langle F_f(g), V^* R(g) V \rangle_\sigma \ge 1 - \epsilon. \label{eq:inner1}
% \end{equation}
%     Let $\{ |c\rangle \}_{c \in \FF_q^\times}$ be the orthonormal basis for $\mathcal{H}$, and the set of characters $\{\chi\}_{\chi\in\hat{G}}$ be an orthonormal basis for $L(G, \cc)$. 
    
%     Because $F_f(g)$ is diagonal in $\{\ket{c}\}$ with entries $\omega^{\Tr\, caf(x)}$, we can rewrite the inner product trace:
%     \[
%         \langle F_f(g), V^* R(g) V \rangle_\sigma = \frac{1}{q-1} \sum_{c \in \FF_q^\times} \langle c | F_f^*(g) V^* R(g) V | c \rangle = \frac{1}{q-1} \sum_{c \in \FF_q^\times} \omega^{-\Tr\, caf(x)} \langle v_c, R(g) v_c \rangle_{L(G, \mathcal{H})},
%     \]
%     where $\ket{v_c} = V\ket{c} \in \mathcal{H}^\prime$.
%     Taking the expectation over $g = (x, a)$ and taking the real part, condition ~\equref{eq:inner1} becomes:
%     \begin{equation}
%         \frac{1}{q-1} \sum_{c \in \FF_q^\times} \EE_{x, a} \Re \left[ \omega^{-\Tr\, caf(x)} \langle v_c, R(x, a) v_c \rangle \right] \ge 1 - \epsilon. \label{eq:expanded}
%     \end{equation}

%     A basis for $L(G, \mathcal{H})$ is given by $\{\chi \ket{c}\}_{\chi\in\hat{G}, c\in \FF_q^\times}$. By \propref{prop:group_characaters}, the characters are given by
%     \[
%     \hat{G} = \{\chi_{\alpha, k}\}_{\alpha \in X, k \in [0, \cdots, q-2]}, \chi_{\alpha, k}(x,a) = \omega^{\Tr \inner{\alpha}{a x}} \psi_k(a).
%     \]
%     The regular representation $R(g)$ acts on $L(G, \mathcal{H})$ so it can be decomposed as 
%     \be\label{eq:reg_repr_decomp}
%     R(g) = \sum_{\alpha, k} \chi_{\alpha, k}(g) P_{\alpha, k},
%     \ee
%     where $P_{\alpha, k} = \ket{\chi_{\alpha, k}}\bra{\chi_{\alpha, k}} \otimes \I_{\mathcal{H}}$ projects onto the character function subspace (this leaves $\mathcal{H}$ invariant). 
    
%     Substituting this into the matrix element for $v_c$, we obtain:
%     \[
%         \langle v_c, R(g) v_c \rangle = \sum_{\alpha, k} \chi_{\alpha, k}(g) \langle v_c, P_{\alpha, k} v_c \rangle.
%     \]
%     Let $p_{c, \alpha, k} := \langle v_c, P_{\alpha, k} v_c \rangle$. As $P_{\alpha, k}$ is a projector, then 
%     \be
%     p_{c, \alpha, k} = \langle v_c \ket{\chi_{\alpha, k}}\bra{\chi_{\alpha, k}} \otimes \I_{\mathcal{H}}\ket{v_c} \geq 0.
%     \ee
    
%     Since $V$ is an isometry $V^* V = \I_{\mathcal{H}}$, we have $\bra{c} V^* V \ket{c} = \inner{c}{c}_{\mathcal{H}} = 1$. 
%     Equivalently, we have $\inner{v_c}{v_c}_{L(G, \mathcal{H})} = 1$. 
%     Since $\sum P_{\alpha, k} = I_{L(G,H)}$, it follows that $p_{c, \alpha, k}$ is a valid probability distribution over the characters,
%     \be
%     \sum_{\alpha, k} p_{c, \alpha, k} = 1, \quad \forall c \in \FF_q^\times.
%     \ee

%     Substituting the spectral decomposition ~\equref{eq:reg_repr_decomp} back into ~\equref{eq:expanded}, we get:
%     \[
%         \sum_{c \in \FF_q^\times} \left( \frac{1}{q-1} \right) \sum_{\alpha, k} p_{c, \alpha, k} \Re \left( \EE_{x, a} \left[ \omega^{-\Tr\, caf(x)} \chi_{\alpha, k}(x, a) \right] \right) \ge 1 - \epsilon.
%     \]
%     This is a convex combination (a weighted average) of real values less than $1$. Because the average is at least $1 - \epsilon$, there must exist at least one specific configuration $(c^*, \alpha^*, k^*)$ with $c^* \neq 0$ such that its individual value satisfies:
%     \begin{equation}
%         \Re \left( \EE_{x, a} \left[ \omega^{-\Tr\, c^*af(x)} \omega^{\Tr \langle \alpha^*, ax \rangle} \psi_{k^*}(a) \right] \right) \ge 1 - \epsilon. \label{eq:dominant}
%     \end{equation}

%     Because $c^* \in \FF_q^\times$ is invertible, we define a specific linear function $\ell(x) := \langle (c^*)^{-1} \alpha^*, x \rangle$. Notice that $\langle \alpha^*, ax \rangle = c^* a \ell(x)$. We can rewrite the argument of the expectation in ~\equref{eq:dominant} as:
%     \[
%         E^* := \EE_{x} \left[ \EE_{a} \left[ \omega^{\Tr\, c^* a (\ell(x) - f(x))} \psi_{k^*}(a) \right] \right].
%     \]

%     For the rest of this proof, without using specific form of $p_{c, \alpha, k}$, we consider two cases: $k^*=0$ and $k^* \neq 0$. Define $z(x) = c^*(\ell(x) - f(x))$.

%     \begin{enumerate}
%         \item Consider $k^*=0$ so that $\psi_0(a) = 1$, the inner expectation over $a \in \FF_q^\times$ for a fixed $x$:
%     \be
%     \EE_{a \in \in \FF_q^\times}\omega^{\Tr\, c^* a (\ell(x) - f(x))} = 
%     \begin{cases}
%         1, \quad z(x) = 0, \\
%         - \frac{1}{q-1}, \quad z(x) \neq 0,
%     \end{cases}
%     \ee
%     where in the second case we have again used average of roots of unity is $\sum_{b \in \in \FF_q^\times} \omega^{\Tr b} +1 = 0$.

%     Evaluating the outer expectation over $x$ and recall the definition of hamming distance 
%     \be
%     \delta_\alpha := \Pr_x[f(x) \ne \ell_\alpha(x)],
%     \ee
%     we find:
%     \[
%         \Re E^* = E^* = 1 \cdot (1 - \delta(f, \ell)) + \left( \frac{-1}{q-1} \right) \cdot \delta(f, \ell) = 1 - \delta(f, \ell) \left( 1 + \frac{1}{q-1} \right) = 1 - \delta(f, \ell) \left( \frac{q}{q-1} \right).
%     \]
%     Applying our bound from ~\equref{eq:dominant}, we have:
%     \[
%         1 - \delta(f, \ell) \left( \frac{q}{q-1} \right) \ge 1 - \epsilon.
%     \]
%     Rearranging yields the final result:
%     \[
%         \delta(f, \ell) \le \frac{q-1}{q} \epsilon.
%     \]

%     \item Consider $k^* \neq 0$, then by \propref{prop:group_characaters} we have $\psi_k(a) = e^{\frac{i 2 \pi m k}{q-1}}$ and $a=g_0^m$ for a fixed multiplicative generator $g_0\in \FF_q^\times$. Note that averaging over $\psi_k(a)$ we have
%     \be\label{eq: averaging_psi}
%     \EE_{a \in \FF_q^\times} \psi_k(a) = \frac{1}{q-1} \sum_{m=0}^{q-2} e^{\frac{i 2 \pi m k}{q-1}} = 0.
%     \ee
%     The average expectation over $a \in \FF_q^\times$ for a fixed $x$ gives:
%     \be
%     \EE_{a \in \in \FF_q^\times}\omega^{\Tr \, c^* a (\ell(x) - f(x))}  \psi_{k^*}(a) = 
%     \begin{cases}
%         0, \quad z(x) = 0, \\
%         \frac{1}{q-1} \sum_{a \in \in \FF_q^\times} \omega^{\Tr a z(x)}\psi_{k^*}(a), \quad z(x) \neq 0.
%     \end{cases}
%     \ee
%     In the second case, let $b = a z(x)$ so that
%     \begin{align}
%         \frac{1}{q-1} \sum_{a \in \in \FF_q^\times} \omega^{\Tr\, a z(x)}\psi_{k^*}(a) &= \frac{1}{q-1} \sum_{b \in \in \FF_q^\times} \omega^{\Tr b}\psi_{k^*}(z^{-1}b), \nonumber \\
%         & =  \frac{\psi_{k^*}(z^{-1})}{q-1} \sum_{b \in \in \FF_q^\times} \omega^{\Tr b}\psi_{k^*}(b).
%     \end{align}
%     It turns out the summation is a known quantitiy called \textit{the Gauss sum}
%     \be
%     \mathcal{G}(\psi_{k^*}) = \sum_{b \in \in \FF_q^\times} \omega^{\Tr b}\psi_{k^*}(b).
%     \ee
%     Is it believed that there is no analytic expression for $\mathcal{G}$ but it is known that
%     \be\label{eq:gauss_sum_mag}
%     \abs{\mathcal{G}(\psi_{k^*})} = \sqrt{q}
%     \ee

%     Evaluating the outer expectation over $x$ and recall the definition of hamming distance gives
%     \[
%     \Re E^* = 0 \cdot (1 - \delta(f, \ell)) + \Re \left( \frac{\psi_{k^*}(z^{-1})}{q-1} \mathcal{G}(\psi_{k^*})\right) \cdot \delta(f, \ell) .
%     \]
%     Therefore we have
%     \[
%     \Re E^* = \Re \left( \frac{\psi_{k^*}(z^{-1})}{q-1} \mathcal{G}(\psi_{k^*})\right) \delta(f, \ell) \leq \frac{\abs{\mathcal{G}(\psi_{k^*})}}{q-1} = \frac{\sqrt{q}}{q-1}.
%     \]
%     On the other hand, from ~\equref{eq:dominant} we have
%     \[
%     \Re E^* \geq 1- \epsilon.
%     \]

%     This is a contradiction to the existence of $k^*$ if 
%     \be\label{eq:epsilon_bound}
%     \frac{\sqrt{q}}{q-1} < 1- \epsilon.
%     \ee
%     Then if we have a good test outcome such that ~\equref{eq:epsilon_bound} is satisfied, then we only have a contribution from $k^*=0$, where we have already achieved the bound.
%     \end{enumerate}
    
% \end{proof}

% The above bound is tight assuming we pass the test with high probability such that $\epsilon < 1 - \frac{\sqrt{q}}{q-1} $. Next, we try to give a proof that is unconditional, using Fourier coefficients. 

First let us write down statements about Fourier coefficients of $F_f$, which we will use in the proof.

\begin{proposition}[Fourier coefficient of $F_f$]\label{prop:Ff_fourier}\uses{def:Flift, prop:group_characaters}\lean{BlrPcp.gFourierCoeff_factor}\leanok
    Let the matrix lift $F_f : G \to U(\cc^{q-1})$ be defined by the diagonal matrix
\begin{equation*}
F_f(x, a) = D(a f(x)) := \sum_{c \in A} \omega^{\Tr\, c a f(x)}\ket{c}\bra{c},
\end{equation*}
then its Fourier coefficients can be written as
\be\label{eq:factored-coeff}
[\hat{F}_f(\chi_{c\beta, k})]_{cc} = \psi_k(c) \nu_k(\beta),
\ee
where 
\be\label{def:nu-k}
\nu_k(\beta) := \EE_{x,b} \left[ \omega^{\Tr b(f(x) - \langle \beta, x \rangle)} \overline{\psi_k(b)} \right],
\ee

\end{proposition}

\begin{proof}\uses{prop:group_characaters}\leanok
    By \propref{prop:group_characaters}, the characters are given by
    \[
    \hat{G} = \{\chi_{\alpha, k}\}_{\alpha \in X, k \in [0, \cdots, q-2]}, \chi_{\alpha, k}(x,a) = \omega^{\Tr \inner{\alpha}{a x}} \psi_k(a).
    \]
    Therefore the Fourier coefficients of $F_f$ are
    \begin{align}
        \hat{F}_f(\chi_{\alpha, k}) 
        & = \sum_{c \in \FF_q^\times } \EE_{x, a} \omega^{\Tr\, c a f(x)} \bar{\chi}_{\alpha, k}(x, a) , \nonumber \\
        & = \sum_{c \in \FF_q^\times }\EE_{x,a} \left[ \omega^{\Tr\, a(cf(x) - \langle \alpha, x \rangle)} \overline{\psi_k(a)} \right] \ket{c}\bra{c}.
    \end{align}
    

Since $c \in \FF_q^\times$ is invertible, we perform a change of variables. Let $\alpha = c\beta$ for a fresh $\beta \in \FF_q^n$. Because $c$ is fixed, this is a bijection on $\FF_q^n$. 
We also substitute $b = ca$, which is a bijection on $\FF_q^\times$. 
This means $a = bc^{-1}$, and by the multiplicity of Dirichlet characters, $\overline{\psi_k(a)} = \overline{\psi_k(bc^{-1})} = \overline{\psi_k(b)}\psi_k(c)$.

Substituting these into the expectation isolates the $c$-dependence: 
\begin{align} 
    [\hat{F}_f(\chi_{c\beta, k})]_{cc} &= \EE_{x,b} \left[ \omega^{\Tr b(f(x) - \langle \beta, x \rangle)} \overline{\psi_k(b)}\psi_k(c) \right] \nonumber \\
    &= \psi_k(c) \nu_k(\beta), 
\end{align}
where $\nu_k(\beta) := \EE_{x,b} \left[ \omega^{\Tr b(f(x) - \langle \beta, x \rangle)} \overline{\psi_k(b)} \right]$ is a $c$-independent coefficient.    
\end{proof}


% \begin{lemma}[Gowers-Hatami to linearity (Convex character Fourier hybrid)] \label{lem:gh_linearity_2}
% Suppose $F_f: G \to U(\cc^{q-1})$ defined in ~\equref{eq:Flift-def} via $f : \FF_q^n \to \FF_q$ is an $(\epsilon, \sigma)$ approximate unitary representation of $(G, \star)$ defined in ~\equref{eq:star}. 
% Then, there exists a linear function $\ell: \FF_q^n \to \FF_q$ such that the relative Hamming distance $\delta(f, \ell)$ is bounded by:
% \be\label{eq:GH-linear-bound}
%     \delta(f, \ell) \le \frac{q-1}{q} \epsilon.
% \ee
% \end{lemma}

% \begin{proof}

%     By Gowers-Hatami theorem, we have an isometry $V: \mathcal{H} \to \mathcal{H}^\prime$ and an exact representation $R: G \to \End(\mathcal{H}^\prime)$, where the Hilbert space $\mathcal{H}\cong \cc^{q-1}$ and the space $R$ acts on $\mathcal{H}^\prime = L(G, \mathcal{H})$
%     \begin{equation}
%     \EE_{g \in G} \Re \langle F_f(g), V^* R(g) V \rangle_\sigma \ge 1 - \epsilon. \label{eq:inner1}
% \end{equation}
%     Let $\{ |c\rangle \}_{c \in \FF_q^\times}$ be the orthonormal basis for $\mathcal{H}$, and the set of characters $\{\chi\}_{\chi\in\hat{G}}$ be an orthonormal basis for $L(G, \cc)$. 
    
%     % Since $V$ is an isometry, the vectors $\{v_c := V|c\rangle\}$ form an orthonormal set of $q-1$ basis in $L(G, \mathcal{H})$. 
    
%     Because $F_f(g)$ is diagonal in $\{\ket{c}\}$ with entries $\omega^{\Tr\, caf(x)}$, we can rewrite the inner product trace:
%     \[
%         \langle F_f(g), V^* R(g) V \rangle_\sigma = \frac{1}{q-1} \sum_{c \in \FF_q^\times} \langle c | F_f^*(g) V^* R(g) V | c \rangle = \frac{1}{q-1} \sum_{c \in \FF_q^\times} \omega^{-\Tr\, caf(x)} \langle v_c, R(g) v_c \rangle_{L(G, \mathcal{H})},
%     \]
%     where $\ket{v_c} = V\ket{c} \in \mathcal{H}^\prime$.
%     Taking the expectation over $g = (x, a)$ and taking the real part, condition ~\equref{eq:inner1} becomes:
%     \begin{equation}
%         \frac{1}{q-1} \sum_{c \in \FF_q^\times} \EE_{x, a} \Re \left[ \omega^{-\Tr\, caf(x)} \langle v_c, R(x, a) v_c \rangle \right] \ge 1 - \epsilon. \label{eq:expanded}
%     \end{equation}

%     A basis for $L(G, \mathcal{H})$ is given by $\{\chi \ket{c}\}_{\chi\in\hat{G}, c\in \FF_q^\times}$. By \propref{prop:group_characaters}, the characters are given by
%     \[
%     \hat{G} = \{\chi_{\alpha, k}\}_{\alpha \in X, k \in [0, \cdots, q-2]}, \chi_{\alpha, k}(x,a) = \omega^{\Tr \inner{\alpha}{a x}} \psi_k(a).
%     \]

%     Moreover, the Fourier coefficients of $F_f$ are
%     \begin{align}
%         \hat{F}_f(\chi_{\alpha, k}) 
%         & = \sum_{c \in \FF_q^\times } \EE_{x, a} \omega^{\Tr\, caf(x)} \bar{\chi}_{\alpha, k}(x, a) , \nonumber \\
%         & = \sum_{c \in \FF_q^\times }\EE_{x,a} \left[ \omega^{\Tr\, a(cf(x) - \langle \alpha, x \rangle)} \overline{\psi_k(a)} \right] \ket{c}\bra{c}.
%     \end{align}
    
%     The regular representation $R(g)$ acts on $L(G, \mathcal{H})$ so it can be decomposed as 
%     \be\label{eq:reg_repr_decomp}
%     R(g) = \sum_{\alpha, k} \chi_{\alpha, k}(g) P_{\alpha, k},
%     \ee
%     where $P_{\alpha, k} = \ket{\chi_{\alpha, k}}\bra{\chi_{\alpha, k}} \otimes \I_{\mathcal{H}}$ projects onto the character function subspace (this leaves $\mathcal{H}$ invariant). 
    
%     Substituting this into the matrix element for $v_c$, we obtain:
%     \[
%         \langle v_c, R(g) v_c \rangle = \sum_{\alpha, k} \chi_{\alpha, k}(g) \langle v_c, P_{\alpha, k} v_c \rangle.
%     \]
%     Let $p_{c, \alpha, k} := \langle v_c, P_{\alpha, k} v_c \rangle$. As $P_{\alpha, k}$ is a projector, then 
%     \be
%     p_{c, \alpha, k} = \langle v_c \ket{\chi_{\alpha, k}}\bra{\chi_{\alpha, k}} \otimes \I_{\mathcal{H}}\ket{v_c} \geq 0.
%     \ee
%     Let the isometry be given by $V\ket{c} = [F_f(g)]_{c}$, then the matrix element is given by the Fourier coefficients
%     \be
%     p_{c, \alpha,k} = \abs{[\hat{F}_f(\chi_{\alpha, k})]_{cc}}^2.
%     \ee
%     By \propref{prop:Ff_fourier}, let $\alpha = c \beta$, we can show
%     \be
%     p_{c, c\beta,k} = \abs{\psi_k(c) \nu_k(\beta)}^2 = \abs{\nu_k(\beta)}^2,
%     \ee
%     where 
%     \[
%     \nu_k(\beta) := \EE_{x,b} \left[ \omega^{\Tr\, b(f(x) - \langle \beta, x \rangle)} \overline{\psi_k(b)} \right],
%     \]
%     is defined in ~\equref{def:nu-k}.   
    
%     Since $V$ is an isometry $V^* V = \I_{\mathcal{H}}$, we have $\bra{c} V^* V \ket{c} = \inner{c}{c}_{\mathcal{H}} = 1$. 
%     Equivalently, we have $\inner{v_c}{v_c}_{L(G, \mathcal{H})} = 1$. 
%     Since $\sum P_{\alpha, k} = I_{L(G,H)}$, it follows that $p_{c, \alpha, k}$ is a valid probability distribution over the characters,
%     \be
%     \sum_{\alpha, k} p_{c, \alpha, k} = 1, \quad \forall c \in \FF_q^\times.
%     \ee

%     Equivalently, we have the relabeled probability
%     \be
%     \sum_{\beta, k} p_{c, c \beta, k} = 1, \quad \forall c \in \FF_q^\times.
%     \ee
%     Therefore $|\nu_k(\beta)|^2$ is a probability distribution indexed by $(k, \beta)$. Because the $k=0$ probability is a subset, we have
%     \be\label{eq:nu_less_1}
%     \sum_{\beta \in \FF_q^n} |\nu_0(\beta)|^2 \leq \sum_{\beta \in \FF_q^n} \sum_k |\nu_k(\beta)|^2 = 1.
%     \ee
    
%     Substituting the spectral decomposition ~\equref{eq:reg_repr_decomp} back into ~\equref{eq:expanded}, we define:
%     \begin{align*}
%             E &:= \frac{1}{q-1} \sum_{c \in \FF_q^\times} \EE_{x, a} \Re \left[ \omega^{-\Tr\, caf(x)} \langle v_c, R(x, a) v_c \rangle \right], \\
%             % & = \sum_{c \in \FF_q^\times} \left( \frac{1}{q-1} \right) \sum_{\alpha, k} p_{c, \alpha, k} \Re \left( \EE_{x, a} \left[ \omega^{-\Tr\, caf(x)} \chi_{\alpha, k}(x, a) \right] \right), \\
%             & = \frac{1}{q-1} \sum_{c \in \FF_q^\times}\sum_{\alpha, k} p_{c, \alpha, k} \Re \left( \EE_{x, a} \left[ \omega^{-\Tr\, caf(x)} \omega^{\Tr\, \inner{\alpha}{ax}} \psi_k(a) \right] \right).
%     \end{align*}

%     Relabel the sum by $\alpha = c \beta$ and $b=ca$, we have
%     \[
%     E =  \frac{1}{q-1} \sum_{c \in \FF_q^\times}\sum_{\beta, k} \abs{\nu_k(\beta)}^2 \Re \left( \EE_{x, b} \left[ \omega^{\Tr\, b(\inner{\beta}{x} - f(x))} \psi_k(c^{-1})\psi_k(b) \right] \right).
%     \]
%     Use the indicator function of $\FF_q^\times$:
%     \[
%         \sum_{c \in \FF_q^\times} {\psi_k(c^{-1})} = (q-1) \I(k=0)
%     \] 
%     we have     
%     \[
%     E = \sum_{\beta} \abs{\nu_0(\beta)}^2 \Re \left( \EE_{x, b} \left[ \omega^{\Tr\, b(\inner{\beta}{x} - f(x))} \right] \right).
%     \]
    
%     This is a convex combination (a weighted average) of real values less than $1$ with the weight sum to less than 1 by ~\equref{eq:nu_less_1}. Because the average is at least $E \geq 1 - \epsilon$, there must exist at least one specific configuration $\beta^*$ such that its individual value satisfies:
%     \begin{equation}
%         E^* := \Re \left( \EE_{x, b} \left[ \omega^{\Tr\, b(\inner{\beta^*}{x} - f(x))} \right] \right) \ge 1 - \epsilon. \label{eq:dominant}
%     \end{equation}
    
%     Let $\ell_\beta(x):=\inner{\beta^*}{x}$, the inner expectation over $b \in \FF_q^\times$ for a fixed $x$:
%         \be
%         \EE_{b \in \in \FF_q^\times}\omega^{\Tr\, b z(x)} = 
%         \begin{cases}
%             1, \quad z(x) = 0, \\
%             - \frac{1}{q-1}, \quad z(x) \neq 0,
%         \end{cases}
%         \ee
%         where in the second case we have again used average of roots of unity is $\sum_{b \in \in \FF_q^\times} \omega^{\Tr\, b} +1 = 0$.
    
%         Evaluating the outer expectation over $x$ and recall the definition of hamming distance 
%         \be
%         \delta(f, \ell_\alpha) := \Pr_x[f(x) \ne \ell_\alpha(x)],
%         \ee
%         we find:
%         \[
%             E^* = 1 \cdot (1 - \delta(f, \ell_\beta)) + \left( \frac{-1}{q-1} \right) \cdot \delta(f, \ell_\beta) = 1 - \delta(f, \ell_\beta) \left( 1 + \frac{1}{q-1} \right) = 1 - \delta(f, \ell_\beta) \left( \frac{q}{q-1} \right).
%         \]
%         Applying our bound from ~\equref{eq:dominant}, we have:
%         \[
%             1 - \delta(f, \ell_\beta) \left( \frac{q}{q-1} \right) \ge 1 - \epsilon.
%         \]
%         Rearranging yields the final result: $\exists \ell_\beta = \inner{\beta^*}{x}$ such that 
%         \[
%             \delta(f, \ell_\beta) \le \frac{q-1}{q} \epsilon.
%         \]
    
% \end{proof}


\begin{lemma}[Gowers-Hatami to linearity (formalized proof)] \label{lem:gh_linearity_formalized}\uses{lem:group, def:Flift, approx-representation, def:linear-functions, def:relative-hamming-distance}\lean{BlrPcp.gh_linearity}\leanok
Suppose $F_f: G \to U(\cc^{q-1})$ defined in ~\equref{eq:Flift-def} via $f : \FF_q^n \to \FF_q$ is an $(\epsilon, \sigma)$ approximate unitary representation of $(G, \star)$ defined in ~\equref{eq:star}. 
Implicitly, $\epsilon \leq 1$.

Then, there exists a linear function $\ell: \FF_q^n \to \FF_q$ such that the relative Hamming distance $\delta(f, \ell)$ is bounded by:

\be\label{eq:GH-linear-bound}
\delta(f, \ell) \le \frac{q-1}{q} \epsilon.
\ee
\end{lemma}

\begin{proof}\uses{gowers-hatami-correlation, prop:group_characaters, prop:Ff_fourier}\leanok
Write the diagonal lift as $F_f(g) = \operatorname{diag}_{c \in \FF_q^\times}(\varphi_c(g))$, where the scalar functions are $\varphi_c(x,a) = \omega^{\Tr(caf(x))}$. 

Because $F_f$ is diagonal and $\sigma = \I/(q-1)$, the trace equates to the average over the diagonal elements. 
By an intermediate statement in the proof of Gowers-Hatami theorem using Lemma.~\ref{gowers-hatami-correlation}, the compression from an exact representation can be rewritten as a pair-correlation:
\begin{equation}
    \mathcal{C}(f) := \EE_{g,h} \Re \inner{F_f(g) F_f(h)}{F_f(gh)}_\sigma \ge 1 - \epsilon. \label{eq:gh-pair-correlation}
\end{equation}


To connect the operator pair-correlation directly to the cubic trace of the Fourier coefficients, we express the lift matrices using the inverse Fourier transform, $F_f(x) = \sum_{\chi \in \widehat{G}} \hat{F}_f(\chi) \chi(x)$. Substituting this into the inner product definition of the pair-correlation and using the character homomorphism $\chi(gh) = \chi(g)\chi(h)$, we expand the trace:
\begin{align}
    \mathcal{C}(f) &= \EE_{g,h \in G} \Re \Tr\left( F_f(g)^* F_f(h)^* F_f(gh) \sigma \right) \nonumber \\
    &= \Re \Tr \left( \sum_{\chi_1, \chi_2, \chi_3} \hat{F}_f^*(\chi_1) \hat{F}_f^*(\chi_2) \hat{F}_f(\chi_3) \sigma \cdot \EE_g\left[\overline{\chi_1(g)}\chi_3(g)\right] \EE_h\left[\overline{\chi_2(h)}\chi_3(h)\right] \right).
\end{align}
By the orthogonality of characters, the expectations over $g$ and $h$ independently vanish unless $\chi_1 = \chi_3$ and $\chi_2 = \chi_3$. This completely collapses the triple sum to a single sum over $\chi$. Because the Fourier coefficients of the diagonal lift $F_f$ are themselves diagonal matrices (and thus commute), we can rearrange the surviving terms to recover the exact cubic trace expression:
\begin{equation}
    \mathcal{C}(f) = \Re \sum_{\chi \in \widehat{G}} \Tr\left( \hat{F}_f^*(\chi) \hat{F}_f(\chi) \hat{F}_f^*(\chi) \sigma \right).
\end{equation}

    Because the Fourier coefficients are all fully diagonal since $F_f$ is diagonal, we expand the product
    \[
    \sum_{\chi_{\alpha, k} \in \widehat{G}} \Tr(\hat{F}_f^*(\chi)\hat{F}_f(\chi)\hat{F}_f^*(\chi) \sigma) = \frac{1}{q-1}\sum_{\alpha \in \FF_q^n} \sum_{k=0}^{q-2} \sum_{c \in \FF_q^\times} \left| [\hat{F}_f(\chi_{\alpha, k})]_{cc} \right|^2 \overline{[\hat{F}_f(\chi_{\alpha, k})]_{cc}}.
    \]

    Using the substitution $\alpha = c\beta$, summing over all $\alpha$ is equivalent to summing over all $\beta$. 
Substituting our factored coefficient from ~\eqref{eq:factored-coeff} in \propref{prop:Ff_fourier}: 
    \begin{align*}
        \sum_{\chi \in \widehat{G}} \Tr(\hat{F}_f^*(\chi)\hat{F}_f(\chi)\hat{F}_f^*(\chi) \sigma) &= \frac{1}{q-1}\sum_{c \in \FF_q^\times} \sum_{\beta \in \FF_q^n} \sum_{k=0}^{q-2} |\psi_k(c)\nu_k(\beta)|^2 \overline{\psi_k(c)\nu_k(\beta)} \\
        &= \frac{1}{q-1}\sum_{\beta, k} |\nu_k(\beta)|^2 \overline{\nu_k(\beta)} \sum_{c \in \FF_q^\times} \overline{\psi_k(c)},
    \end{align*}
    where 
    \[
    \nu_k(\beta) := \EE_{x,b} \left[ \omega^{b(f(x) - \langle \beta, x \rangle)} \overline{\psi_k(b)} \right],
    \]
    is defined in \equref{def:nu-k}.

    Use the indicator function of $\FF_q^\times$:
    \[
        \sum_{c \in \FF_q^\times} \overline{\psi_k(c)} = (q-1) \I(k=0)
    \] 
we have
    \be\label{eq:triple-F-nu0}
        \sum_{\chi \in \widehat{G}} \Tr(\hat{F}_f^*(\chi)\hat{F}_f(\chi)\hat{F}_f^*(\chi) \sigma) = \sum_{\beta \in \FF_q^n} |\nu_0(\beta)|^2 \overline{\nu_0(\beta)}.
    \ee
    Note that $\nu_0(\beta) \in \mathbb{R}$ is related to the hamming distance $\delta(f, \ell_\beta)$ for liner function $\ell_\beta = \inner{f}{\ell_\beta}$ via
    \be\label{eq:nu0-hamming}
    \nu_0(\beta) = (1 - \delta(f, \ell_\beta)) - \frac{1}{q-1}\delta(f, \ell_\beta) = 1 - \frac{q}{q-1}\delta(f, \ell_\beta).
    \ee
    Moreover, because the Fourier coefficients sum to identity by unitarity of $F_f$, we have
    \begin{align}
        1& = \frac{1}{q-1}\sum_{\chi\in\hat{G}}\Tr\left[\hat{F}^*_f(\chi)\hat{F}_f(\chi)\right] , \nonumber \\
        &= \frac{1}{q-1}\sum_{c \in \FF_q^\times} \sum_{\beta \in \FF_q^n} \sum_k |\psi_k(c)\nu_k(\beta)|^2, \nonumber \\
        &= \sum_{\beta \in \FF_q^n} \sum_k |\nu_k(\beta)|^2.
    \end{align}

    Therefore $|\nu_k(\beta)|^2$ is a probability distribution indexed by $(k, \beta)$. Because the $k=0$ probability is a subset, we have
    \be
    \sum_{\beta \in \FF_q^n} |\nu_0(\beta)|^2 \leq \sum_{\beta \in \FF_q^n} \sum_k |\nu_k(\beta)|^2 = 1.
    \ee

    Let $\beta^*$ be the closest linear function such that
    \be
    \nu_0(\beta^*) := \max_{\beta \in \FF_q^n} \nu_0(\beta).
    \ee
    Therefore from Gowers-Hatami implication in \equref{eq:gh-pair-correlation} and the simplification of Fourier coefficeints in \equref{eq:triple-F-nu0}, we have the bound
    \begin{align}
        1 - \epsilon \leq \Re \sum_{\chi\in\widehat G} \Tr(\hat{F}_f^*(\chi)\hat{F}_f(\chi)\hat{F}_f^*(\chi) \sigma) = \sum_{\beta \in \FF_q^n} |\nu_0(\beta)|^2 \overline{\nu_0(\beta)} \leq  \nu_0(\beta^*) \sum_{\beta \in \FF_q^n} |\nu_0(\beta)|^2 \leq  \nu_0(\beta^*).
    \end{align}

    Unpack the relation to the hamming distance, we have
    \be
    1 - \frac{q}{q-1} \delta(f, \ell_{\beta^*}) \ge 1 - \epsilon \implies \delta(f, \ell_{\beta^*}) \le \frac{q-1}{q}\epsilon.
    \ee

\end{proof}

% Finally, we are ready to prove the main theorem of BLR linearity test soundness. 


% \setcounter{theorem}{0}
% \begin{theorem}[Soundness of the BLR test]\label{thm:main}\uses{lem:rep-iff-blr, lem:gh_linearity_formalized}\lean{BlrPcp.blr_soundness}\leanok
% Let $q$ be prime. For every $f : \F_q^n \to \F_q$, $\delta(f, \lin) \le \Pr_{a,b,x,y}\!\bigl[af(x) + bf(y) \ne f(ax + by)\bigr]$.
% \end{theorem} 

\begin{lemma}
[Soundness of the BLR test]\label{BLR_via_GH}
\lean{BlrPcp.gh_blr_soundness}
\uses{def:finite-field-blr-test,def:distance-to-linearity}
\leanok
Let $q=p^s$ be a prime power. For every $f : \F_q^n \to \F_q$, $\delta(f, \lin) \le \Pr_{a,b,x,y}\!\bigl[af(x) + bf(y) \ne f(ax + by)\bigr]$.
\end{lemma}

\begin{proof}[Alternative proof via Gowers–Hatami]\uses{lem:rep-iff-blr, lem:gh_linearity_formalized}\leanok
First denote
\[
X := \F_q^n, \qquad A := \F_q^\times, \qquad \omega := e^{2\pi i/p}.
\]
Define a group $G$ on the set $X \times A$ with the binary operation
\[
(x, a) \star (y, b) := \bigl(b^{-1}x + a^{-1}y,\; ab\bigr).
\]
Let $F_f$ be the unitary matrix valued lift $F_f : G \to U(\cc^{q-1})$ given by the diagonal matrix
\[
F_f(x, a) = D(a f(x)) := \sum_{c \in A} \omega^{\Tr\, c a f(x)}\ket{c}\bra{c},
\]
where $\{\ket{c}\}$ denotes the orthonormal basis vectors of a Hilbert space $\mathcal{H} \cong \cc^{q-1}$, the
inner product between operators is $\inner{A}{B}_\sigma=\Tr(A^* B \sigma)$, and $\sigma = \I/(q-1)$. By
\lemref{lem:rep-iff-blr}, since $f$ is rejected with probability $\varepsilon$, the lift $F_f$ is a
$(\frac{q}{q-1}\varepsilon,\frac{\I}{q-1})$ approximate representation of $(G, \star)$.

Then by \lemref{lem:gh_linearity_formalized}, there exists a linear function $\ell: \FF_q^n \to \FF_q$ whose
relative Hamming distance to $f$ is bounded by
\[
    \delta(f, \ell) \le \frac{q-1}{q}\cdot\frac{q}{q-1}\varepsilon = \varepsilon.
\]
This proves the claim
\[
\delta(f, \lin) := \min_{\alpha} \delta(f, \ell_\alpha) \leq \delta(f, \ell) \leq \varepsilon.
\]
Miraculously, the loose factor $\frac{q}{q-1}$ in the approximate-representation parameter cancels the tighter
factor $\frac{q-1}{q}$ of \lemref{lem:gh_linearity_formalized}.

In Lean this route is formalized as \texttt{gh\_blr\_soundness\_acceptance}, which proves
\texttt{acceptanceProbabilityBLR f $\le$ 1 - distanceToLinear f} (the statement of \texttt{blr\_soundness}),
valid for $\frac{q}{q-1}\varepsilon\le1$; the unconditional statement is the direct proof above, since for
$\frac{q}{q-1}\varepsilon>1$ the third-moment extraction's $\max\ge0$ step fails.
\end{proof}

% \begin{proof}
% First denote 
% \[
% X := \F_q^n, \qquad A := \F_q^\times, \qquad \omega := e^{2\pi i/p}.
% \]
% Define a group $G$ by the set $X \times A$ with the binary operation
% \begin{equation}
% (x, a) \star (y, b) := \bigl(b^{-1}x + a^{-1}y,\; ab\bigr).
% \end{equation}

%     Let $F_f$ be a unitary matrix valued lift $F_f : G \to U(\cc^{q-1})$ by the diagonal matrix
% \begin{equation}
% F_f(x, a) = D(a f(x)) := \sum_{c \in A} \omega^{\Tr\, c a f(x)}\ket{c}\bra{c},
% \end{equation}
% where $\{\ket{c}\}$ denotes the orthonormal basis vectors of a Hilbert space $\mathcal{H} \cong \cc^{q-1}$ and the inner product between operators defined by $\inner{A}{B}_\sigma=\Tr(A^* B \sigma)$ and choose $\sigma = \I/(q-1)$. By \lemref{lem:rep-iff-blr}, since $f$ is rejected with probability $\varepsilon$, then $F_f$ is an $(\frac{q}{q-1}\varepsilon,\frac{\I}{q-1})$ approximate representation of $(G, \star)$.

% Then by \lemref{lem:gh_linearity_3}, there exists a linear function $\ell: \FF_q^n \to \FF_q$ such that the relative Hamming distance $\delta(f, \ell)$ is bounded by:
% \be
%     \delta(f, \ell) \le \frac{q-1}{q} \frac{q}{q-1}\varepsilon = \varepsilon.
% \ee

% This proves the claim
% \[
% \delta(f, \lin) := \min_{\alpha} \delta(f, \ell_\alpha) \leq \delta(f, \ell) \leq \varepsilon.
% \]
% \end{proof}

% Miraculously, the loose factor in $(\epsilon^\prime =\frac{q}{q-1}\varepsilon,\frac{\I}{q-1})$ approximate representation $F_f$ cancels with the tighter factor in 
% \[
% \delta(f, \ell) \le \frac{q-1}{q} \epsilon^\prime.
% \]


%%FILE blueprint/src/subsections/gowers-hatami.tex
%%%%%%%%%%%%%%%%%%%%%%
% This file is a self-contained blueprint for the Gowers-Hatami theorem, following Thomas' lecture notes.
% Many definitions have been formalized in representation theory in mathlib:
% https://leanprover-community.github.io/mathlib_docs/representation_theory/basic.html
% However for completeness we include the definitions as comments.
%%%%%%%%%%%%%%%%%%%%%%

\section{Preliminaries}
\begin{definition}[$\sigma$ inner product.]\label{sigmaInner}\lean{sigmaInner}\leanok
    The $\sigma$ inner product between matrices $A$ and $B$ is defined as
    \be
    \inner{A}{B}_{\sigma} = \Tr(A^* B \sigma),
    \ee
    where $\sigma \geq 0$ and $A^*$ denotes conjugate-transpose.
\end{definition}

\begin{definition}[$(\epsilon, \sigma)$-representation.]\label{approx-representation}\uses{sigmaInner}\lean{IsApproxRepresentation}\leanok
Given a finite group \(G\), an integer \(d \geq 1, \varepsilon \geq 0\), and a \(d\)-dimensional positive semidefinite matrix \(\sigma\) with trace 1 , an \((\varepsilon, \sigma)\)-representation of \(G\) is a map \(f: G \rightarrow U_{d}(\mathbb{C})\), to the  \(d \times d\) unitary matrices, such that
\be
\mathbb{E}_{x, y \in G} \Re\left(\left\langle f(x)^{*} f(y), f\left(x^{-1} y\right)\right\rangle_{\sigma}\right) \geq 1-\varepsilon,
\ee
where $\Re$ denotes taking real part of the inner product. 
\end{definition}

\begin{proposition}[$\sigma$-norm Square]\label{sigmaNormSqSub}\uses{sigmaInner}\lean{sigmaNormSq_sub_eq}\leanok
Let $\sigma$ be positive semidefinite and let $A, B$ be matrices of the same dimension. Then
\[
\|A - B\|_\sigma^2 = \|A\|_\sigma^2 + \|B\|_\sigma^2 - 2\,\Re\inner{A}{B}_\sigma.
\]
\end{proposition}
\begin{proof}\leanok
    Expanding by linearity of the inner product,
    \[
    \|A-B\|_\sigma^2 = \Re\inner{A-B}{A-B}_\sigma
    = \|A\|_\sigma^2 + \|B\|_\sigma^2 - \Re\inner{A}{B}_\sigma - \Re\inner{B}{A}_\sigma.
    \]
    Since $\sigma$ is Hermitian, $\inner{B}{A}_\sigma = \overline{\inner{A}{B}_\sigma}$, so the two cross terms have the same real part, yielding $-2\,\Re\inner{A}{B}_\sigma$.
\end{proof}


\section{Representation theory}

% Representation theory exists on mathlib but for completeness we include the definitions as comments

% \begin{definition}[Representation]
%     Let $G$ be a group and $V$ a finite-dimensional complex vector space. A \emph{representation} of $G$ is a homomorphism
%     \be
%         \rho : G \to \mathrm{GL}(V),
%     \ee
%     such that for all $x,y \in G$,
%     \be
%         \rho(xy) = \rho(x)\rho(y).
%     \ee
% \end{definition}

In general, a representation of a group \(G\) on a vector space \(V\) over a
field \(k\) is a homomorphism from \(G\) to the linear endomorphisms of \(V\).
In mathlib this is the abbreviation \(\texttt{Representation}\), defined as
\(G \to^* (V \to_{\!k} V)\), i.e. a monoid homomorphism into the \(k\)-linear
endomorphisms of \(V\).

We call a representation \emph{unitary} when its image lies in the unitary
group; in the finite-dimensional matrix setting used below this means a
homomorphism \(\rho : G \to U(V)\). In the Lean development this is represented
directly as a monoid homomorphism into a matrix unitary group, for example
\(G \to^* \texttt{Matrix.unitaryGroup}\).

% \begin{definition}[Irreducible representation]
%     A representation $\rho : G \to \mathrm{GL}(V)$ is called \emph{irreducible} if the only subspaces $W \subseteq V$ satisfying
%     \be
%         \rho(x) W \subseteq W \quad \text{for all } x \in G
%     \ee
%     are $W = \{0\}$ and $W = V$.
% \end{definition}

% For the rest of the notes, we will write ``irrep'' as a shorthand to denote irreducible representation.

% \begin{definition}[Completely reducible]
%     A representation $\rho : G \to \mathrm{GL}(V)$ is \emph{completely reducible} if $V$ can be decomposed as a direct sum of invariant subspaces,
%     \be
%         V = \bigoplus_i V_i,
%     \ee
%     where each $V_i$ carries an irreducible representation of $G$.
% \end{definition}

% \begin{definition}[Inequivalent representations]
%     Two representations $\rho_1 : G \to \mathrm{GL}(V)$ and $\rho_2 : G \to \mathrm{GL}(W)$ are \emph{equivalent} if there exists an invertible linear map $T : V \to W$ such that
%     \be
%         T \rho_1(x)T^{-1} = \rho_2(x) \quad \text{for all } x \in G.
%     \ee
%     They are \emph{inequivalent} if no such map exists.
% \end{definition}

% \begin{definition}[Character]
%     Let $\rho : G \to \mathrm{GL}(V)$ be a representation. The \emph{character} of $\rho$ is the function $\chi_\rho : G \to \mathbb{C}$ defined by
%     \be
%         \chi_\rho(x) = \mathrm{Tr}[\rho(x)].
%     \ee
% \end{definition}

\begin{definition}[Regular representation]\label{RegularRep}
    Let $R: G \rightarrow \mathbb{C}[G]$ be a regular representation where $\{\ket{g}, g\in G\}$ forms a basis of \(\mathbb{C}[G]\). The representation is given by
    \be
    R(x) \ket{g} = \ket{xg}.
    \ee
    
\end{definition}

It is possible to decompose any representation of a finite group into direct sums over irreps. 

\begin{theorem}[Maschke's theorem]\label{Maschke}
    Let $G$ be a finite group and let $\rho : G \to \mathrm{GL}(V)$ be a finite-dimensional representation over $\mathbb{C}$. Then $\rho$ is completely reducible. In particular, there exists a decomposition
    \be
        \rho \cong \bigoplus_{m} r_m \, \rho_m,
    \ee
    where $\{\rho_m\}$ is a set of inequivalent irreducible representations of $G$, and $r_m \in \mathbb{Z}^+$ are the multiplicities.
\end{theorem}

\begin{sproof}
    Maschke's theorem is partially formalized in $\texttt{Mathlib.RepresentationTheory.Maschke}$ but incomplete.
    Filippo Berta, Jonathan Conrad, and Younous Malliefer are formalizing this.
\end{sproof}

% mathlib: `FDRep.char_orthonormal`
% \begin{theorem}[Petit orthogonality theorem]\label{PetitOrthogonality}
%      Let $R_a$ and $R_b$ be two non-equivalent irreps of $G$. The characters satisfy the orthogonality relation
% \be
% \sum_{g\in G} \chi^*_a(g) \chi_b(g) = \abs{G} \delta_{a,b}.
% \ee    
% \end{theorem}

\begin{proposition}[Decomposition of the regular representation]\label{RegDecompIrrep}\uses{Maschke,RegularRep}
    Let $R : G \to \mathrm{GL}(\mathbb{C}[G])$ be the regular representation of a finite group $G$. Then
    \be
        R \cong \bigoplus_{\rho \in \widehat{G}} d_\rho \, \rho,
    \ee
    where $\widehat{G}$ is a complete set of inequivalent irreducible representations of $G$, and $d_\rho = \dim(V_\rho)$ is the dimension of $\rho$.
\end{proposition}

\begin{sproof}
    This can be proven using Maschke's theorem in \ref{Maschke} to decompose into irreps with multiplicity. 
    Then use character orthogonality theorem in mathlib $\texttt{FDRep.char\_orthonormal}$ to show the multiplicity equals irrep dimension. 
    A more formal statement is given by the Peter-Weyl theorem. 
\end{sproof}


\begin{proposition}[Character of regular representation.]\label{regular_repr_character}\uses{RegDecompIrrep}
Let $G$ be a finite group and $\sum_\rho$ denotes summing over the complete set of inequivalent irreps $\hat{G}$
    \be\label{eq:regular-repr-characterr}
    \sum_{\rho \in \hat{G}} d_{\rho} \operatorname{Tr}(\rho(x))=|G| \delta_{x, e}.
    \ee
\end{proposition}
\begin{sproof}
    The proof uses decomposition of regular representation and definition of character. 
\end{sproof}

% \begin{fact}
%     Let $\hat{G}$ be the complete set of inequivalent irreps of $G$ and $d_\rho = \dim(\rho)$ the dimension, then
%      \be\label{eq:dim_rho_group}
%      \sum_{\rho} d_{\rho}^{2}=|G|.
%      \ee
% \end{fact}
% \begin{sproof}
%     This is a corollary from Fact.~\ref{regular_repr_character} for $x=e$.
% \end{sproof}

\begin{definition}[Group fourier transform]\label{GroupFourier}
    Let $f: G \rightarrow U_d(\cc)$ be a map and \(\rho: G \rightarrow U_{d_{\rho}}(\mathbb{C})\) be a unitary irrep. The fourier transform of $f$ at the representation $\rho$ is denoted by $\hat{f}$:
    \be
\hat{f}(\rho)=\mathbb{E}_{x \in G} f(x) \otimes \overline{\rho(x)},
    \ee
    where $\overline{\rho(x)}$ denotes complex conjugate of $\rho(x)$. 
\end{definition}

\section{Gowers Hatami Theorem}

\begin{theorem}[Gowers-Hatami]\label{gowers-hatami}
  \uses{sigmaNormSqSub, approx-representation,GroupFourier, regular_repr_character}
Let \(G\) be a finite group, \(\varepsilon \geq 0\), and \(f: G \rightarrow U_{d}(\mathbb{C})\) an \((\varepsilon, \sigma)\)-representation of \(G\). Then there exists a \(d^{\prime} \geq d\), an isometry \(V: \mathbb{C}^{d} \rightarrow \mathbb{C}^{d^{\prime}}\), and a representation \(g: G \rightarrow U_{d^{\prime}}(\mathbb{C})\) such that
\be
\mathbb{E}_{x \in G}\left\|f(x)-V^{*} g(x) V\right\|_{\sigma}^{2} \leq 2 \varepsilon.
\ee 
\end{theorem}

\begin{sproof}
    The theorem is due to Gowers and Hatami~\cite{GowHat17}; the constructive proof below follows Vidick's lecture notes, Theorem 2.3~\cite{VidickQMATH22}.
    In particular, the proof is constructive. The isometry is defined as 
    \(V: \mathbb{C}^{d} \rightarrow \mathbb{C}^{d} \otimes\left(\oplus_{\rho} \mathbb{C}^{d_{\rho}} \otimes \mathbb{C}^{d_{\rho}}\right)\) by
    \be\label{eq:isometry-V}
    V u = \bigoplus_{\rho} d_{\rho}^{1 / 2} \sum_{i=1}^{d_{\rho}}\left(\hat{f}(\rho)\left(u \otimes e_{i}\right)\right) \otimes e_{i}.
    \ee
    And the exact representation is defined as 
    \be
    g(x)=\bigoplus_{\rho}\left(I_{d} \otimes I_{d_{\rho}} \otimes \rho(x)\right).
    \ee

    The proof requires the following two ingredients:
    \begin{enumerate}
        \item $V$ is an isometry:
        \be
        V^{*} V  = \I_d
        \ee 
        \item The isometry ``averages'' over group multiplication. For any \(x \in G\),
        \be
        V^{*} g(x) V = \mathbb{E}_{z} f(z)^{*} f(z x).
        \ee
    \end{enumerate}
    Both ingredients follow from the key identity over the regular representation in \equref{eq:regular-repr-characterr}.
\end{sproof}

Next, we will give a different proof that does not directly decompose into irreps but instead uses the action of the regular representation.
\begin{lemma}[Gowers-Hatami Pair-Correlation]\label{gowers-hatami-correlation}\lean{gh2_average_correlation, gh2_compression_avg_eq_pair, gh2_pair_correlation}\leanok\uses{approx-representation}
Let \(G\) be a finite group, \(\varepsilon \geq 0\), and \(\rho: G \rightarrow U(\mathcal{H})\) an \((\varepsilon, \sigma)\)-representation of \(G\) on space $\mathcal{H}$. 
Define the pair-correlation of $\rho$ as:
\be
    \mathcal{C}(\rho) := \EE_{x,y \in G} \Re \langle \rho(y)\rho(x), \rho(yx) \rangle_\sigma.
\ee
Then there exists an isometry \(V: \mathcal{H} \rightarrow \mathcal{H}^\prime\) to a different space $\mathcal{H}^\prime$, 
and an exact representation \(\rho^\prime: G \rightarrow U(\mathcal{H}^\prime)\) such that the average correlation between $\rho$ and $\rho^\prime$ is exactly the pair-correlation, which is bounded by $1-\varepsilon$:
\be
    \EE_{x \in G} \Re \langle \rho(x), V^{*} \rho^\prime(x) V\rangle_{\sigma} = \mathcal{C}(\rho) \geq 1-\varepsilon.
\ee    
\end{lemma}

\begin{proof}\leanok
    Here $\mathcal{H}^\prime = L(G, \mathcal{H})$ is the space of functions $f: G \to \mathcal{H}$.
    
    Let the isometry be defined by the action of the approximate representation on a state $u\in \mathcal{H}$ as
    \bea
    &V:& \mathcal{H} \to L(G, \mathcal{H}), \\
    & V(u):& G \to \mathcal{H}, \quad x \mapsto \rho(x)(u), \quad \forall x \in G.
    \eea

    The adjoint map $V^*$ is defined as the unique map such that for all $u \in \mathcal{H}$ and $f \in L(G, \mathcal{H})$,
    \be
    \langle V(u), f \rangle_{L(G, \mathcal{H})} = \langle u, V^*(f) \rangle_{\mathcal{H}}.
    \ee

    In particular, $V^*$ is explicitly given by the group average:
    \bea
    &V^*:& L(G, \mathcal{H}) \to \mathcal{H}, \\
    &V^*(f)& = \EE_{x \in G} \left[ \rho^*(x) f(x) \right].
    \eea

    The exact representation is given by the right regular representation on $L(G, \mathcal{H})$:
    \bea
    &\rho^\prime:& G \to U(L(G, \mathcal{H})), \\
    &(\rho^\prime(x)f)(y)& = f(yx), \quad \forall x,y \in G.
    \eea

    $V$ is an isometry because $\rho(x)$ is unitary, yielding:
    \be
    V^* V(u) = \EE_{x \in G} \left[ \rho^*(x) \rho(x)(u) \right] = \EE_{x \in G} [u] = u.
    \ee

    The compression of the exact representation ``averages'' over group multiplication. For any $u \in \mathcal{H}$:
    \bea\label{eq:isometry-average} 
    V^* \rho^\prime(x) V(u) & = & \EE_{y \in G} \left[ \rho^*(y) (\rho^\prime(x)\rho(y))(u) \right] \nonumber \\
    & = & \EE_{y \in G} \left[ \rho^*(y) \rho(yx)(u) \right].
    \eea

    Substituting the average property from ~\equref{eq:isometry-average} into the inner product, and using the unitarity of $\rho(y)$ to rewrite the adjoints ($\langle A, B^* C \rangle_\sigma = \langle B A, C \rangle_\sigma$), we obtain:
    \bea \label{eq:inner-product-bound}
    \EE_{x \in G} \Re \langle \rho(x), V^* \rho^\prime(x) V\rangle_\sigma 
    & = & \EE_{x,y \in G} \Re \langle \rho(x), \rho^*(y) \rho(yx) \rangle_\sigma \nonumber \\ 
    & = & \EE_{x,y \in G} \Re \langle \rho(y)\rho(x), \rho(yx) \rangle_\sigma \nonumber \\ 
    & = & \mathcal{C}(\rho).
    \eea    
    By the definition of an $(\varepsilon, \sigma)$-approximate representation, the pair-correlation inherently satisfies $\mathcal{C}(\rho) \geq 1 - \varepsilon$, which concludes the proof.
\end{proof}
% Below is older lemma that was refactored above
% \begin{lemma}[Gowers-Hatami-Correlation]\label{gowers-hatami-correlation}\lean{gh2_average_correlation}\leanok\uses{approx-representation}
% Let \(G\) be a finite group, \(\varepsilon \geq 0\), and \(\rho: G \rightarrow U(\mathcal{H})\) an \((\varepsilon, \sigma)\)-representation of \(G\) on space $\mathcal{H}$. 
% Then there exists an isometry \(V: \mathcal{H} \rightarrow \mathcal{H}^\prime\) to a different space $\mathcal{H}^\prime$, 
% and a representation \(\rho^\prime: G \rightarrow \mathcal{H}^\prime\) such that the average correlation between $\rho$ and $\rho^\prime$ is:
% \be
% \mathbb{E}_{x \in G} \Re \langle f(x), V^{*} \rho^\prime(x) V\rangle_{\sigma} \geq 1-\varepsilon,
% \ee    
% \end{lemma}

% \begin{proof}
%     Here $\mathcal{H}^\prime = L(G, \mathcal{H})$ is the space of functions $f: G \to \mathcal{H}$.
    
%     Let the isometry be defined by the action of the approximate representation on a state $u\in \mathcal{H}$ as
%     \bea
%     &V:& \mathcal{H} \to L(G, \mathcal{H}), \\
%     & V(u):& G \to \mathcal{H}, \quad x \mapsto \rho(x)(u), \forall x \in G.
%     \eea

%     The adjoint map $V^*$ is defined as the unique map such that for all $u \in \mathcal{H}$ and $g \in L(G, \mathcal{H})$,
%     \be
%     \langle V(u), g \rangle_{L(G, \mathcal{H})} = \langle u, V^*(g) \rangle_{\mathcal{H}}.
%     \ee

%     In particular, $V^*$ is explicitly given by the group average
%     \bea
%     &V^*:& L(G, \mathcal{H}) \to \mathcal{H}, \\
%     &V^*(f)& = \EE_{x \in G} \left[ \rho^*(x) f(x) \right].
%     \eea

%     The exact representation is given by the right regular representation on $L(G, \mathcal{H})$:
%     \bea
%     &\rho^\prime:& G \to \End(L(G, \mathcal{H})), \\
%     &(\rho^\prime(x)f)(y)& = f(yx), \quad \forall x,y \in G.
%     \eea

%     $V$ is an isometry because
%     \be
%     V^* V(u) = \EE_{x \in G} \left[ \rho^*(x) \rho(x)(u) \right] = u.
%     \ee

%     The isometry ``averages'' over group multiplication because
%     \bea\label{eq:isometry-average} 
%     V^* \rho^\prime(x) V(u) & = & \EE_{y \in G} \left[ \rho^*(y) (\rho^\prime(x)\rho(y))(u) \right] \\
%     & = & \EE_{y \in G} \left[ \rho^*(y) (\rho(yx)(u)) \right].\nonumber
%     \eea

%     Finally, using the isometry average property in ~\equref{eq:isometry-average} and the definition of $(\varepsilon, \sigma)$-representation, we have
%     \bea \label{eq:inner-product-bound}
%     && \EE_{x \in G} \Re \langle \rho(x), V^* \rho^\prime(x) V\rangle_\sigma, \\ \nonumber
%     & = & \EE_{x,y \in G} \Re \langle \rho(x), \rho^*(y) \rho(yx) \rangle_\sigma, \\ \nonumber
%     &\ge & 1 - \varepsilon. 
%     \eea    
% \end{proof}

\begin{theorem}[Gowers-Hatami-2]\label{gowers-hatami-2}\lean{gowers_hatami_GH2_product}\leanok
  \uses{sigmaNormSqSub, approx-representation}
Let \(G\) be a finite group, \(\varepsilon \geq 0\), and \(\rho: G \rightarrow U(\mathcal{H})\) an \((\varepsilon, \sigma)\)-representation of \(G\) on space $\mathcal{H}=\mathbb{C}^{d}$. 
Then there exists an isometry \(V: \mathcal{H} \rightarrow \mathcal{H}^\prime\) to a different space $\mathcal{H}^\prime$, 
and a representation \(\rho^\prime: G \rightarrow \mathcal{H}^\prime\) such that
\be
\mathbb{E}_{x \in G}\left\|\rho(x)-V^{*} \rho^\prime(x) V\right\|_{\sigma}^{2} \leq 2 \varepsilon.
\ee 
\end{theorem}

\begin{proof}\leanok
    \uses{gowers-hatami-correlation}
    \leanok
    A regular-representation proof of the Gowers-Hatami theorem appears in Morel's slides~\cite{Morel20}, following the original Gowers-Hatami theorem~\cite{GowHat17}.

    The rest of the proof follows by expanding the $\sigma$-norm, bounding the norm of $\EE_{x \in G} V^* \rho^\prime(x) V$ and applying the definition of $(\varepsilon, \sigma)$-representation.
     In particular, we have 
     \bea
        & &\EE_{x \in G}\left\|\rho(x)-V^{*} \rho^\prime(x) V\right\|_{\sigma}^{2} \\ 
        &= & \EE_{x \in G} \| \rho(x) \|_\sigma^2 + \EE_{x \in G} \|V^* \rho^\prime(x) V\|_\sigma^2 - 2 \EE_{x \in G} \Re \langle \rho(x), V^* \rho^\prime(x) V\rangle_\sigma \nonumber
    \eea        

    Because $\rho$ is unitary and $\Tr \sigma = 1$, $\EE_{x \in G} \| \rho(x) \|_\sigma^2 = 1$.
    For the second term, recall from \equref{eq:isometry-average} that the compression is the group average $V^* \rho^\prime(x) V = \EE_{y \in G} \rho^*(y)\rho(yx)$. Each summand $\rho^*(y)\rho(yx)$ is a product of unitaries, hence unitary, so $\|\rho^*(y)\rho(yx)\|_\sigma = 1$. By convexity of the $\sigma$-norm---the triangle inequality applied to the average---we obtain
    \be \label{eq:isometry-unitary-bound}
    \|V^* \rho^\prime(x) V\|_\sigma = \left\| \EE_{y \in G} \rho^*(y)\rho(yx) \right\|_\sigma \le \EE_{y \in G} \|\rho^*(y)\rho(yx)\|_\sigma = 1,
    \ee
    hence $\EE_{x \in G}\|V^* \rho^\prime(x) V\|_\sigma^2 \le 1$.

    Therefore using the bounds in ~\equref{eq:isometry-unitary-bound} and ~\equref{eq:inner-product-bound} by applying Lemma~\ref{gowers-hatami-correlation}, we have
    \be
        \EE_{x \in G}\left\|\rho(x)-V^{*} \rho^\prime(x) V\right\|_{\sigma}^{2} \le  2 - 2(1 - \varepsilon) = 2\varepsilon.
    \ee
\end{proof}


%%FILE blueprint/src/subsections/lpcp-lineq.tex
\begin{definition}\label{LINEQ}\lean{LINEQ}\leanok
For a field \(\F\) and parameters \(m, n \in \N\), we define the language
\[\LINEQ(\F,m,n) := \{ (M, c) \in \F^{m\times n} \times \F^{m} \mid \exists b\in \F^n, \, Mb = c \}.\]
\end{definition}
\begin{definition}\label{LINEQ.verifier}\lean{LINEQ.verifier}\uses{LPCPVerifier,LINEQ}\leanok
Let \(V_{\LINEQ(\F,m,n)}^{\langle b, \cdot \rangle } (M,c)\) be the LPCP verifier defined as follows:
\begin{enumerate}
  \item Sample \(r \leftarrow \F^m\).
  \item Query \(b\in \F^n \) at \(u:= \Transpose{M}r \in \F^n\).
  \item Accept if and only if \(\langle b,u\rangle = \langle c, r\rangle\).
\end{enumerate}
\end{definition}

\begin{lemma}
\label{LINEQ.linear_form_uniform_prob_mul_card_le_one}
\lean{LINEQ.linear_form_uniform_prob_mul_card_le_one}
\leanok
Let \(a\in\F^m\) be nonzero. Then
\[
  |\F|\cdot\Pr_{r\leftarrow\F^m}\!\left[\langle a,r\rangle=0\right]\leq 1.
\]
Equivalently, a nonzero linear form over \(\F^m\) vanishes on a uniformly
random input with probability at most \(1/|\F|\).
\end{lemma}
\begin{proof}\leanok
Choose a coordinate \(i\) with \(a_i\neq 0\). After fixing all other
coordinates of \(r\), there is exactly one value of \(r_i\) making
\(\langle a,r\rangle=0\).
\end{proof}

\begin{theorem}\label{LINEQ_LPCP}\lean{LINEQ_LPCP}\uses{LINEQ, LPCP}\leanok
  \[\LINEQ(\F,m,n) \in \LPCP\left[\varepsilon_c=0, \varepsilon_s= 1 / |\F|, \Sigma = \F, \ell = |x|, q = 1, r = m\right].\]
\end{theorem}
\begin{proof}\uses{LINEQ.verifier,LINEQ.linear_form_uniform_prob_mul_card_le_one}\leanok
Consider the LPCP verifier \(V_{\LINEQ(\F,m,n)}^{\langle b, \cdot \rangle } (M,c)\).\\
\emph{Completeness:} If \(Mb =c\), then for every \(r \in \F^m\), we have \[\langle b, u \rangle = \Transpose{b} (\Transpose{M} r) = \Transpose{(Mb)} r = \langle c , r\rangle.\]
\emph{Soundness:} If \(Mb \neq c\), then 
\begin{gather*}
  \Pr_{r\leftarrow \F^m} \left[\langle b, u \rangle =  \langle c , r\rangle \right] = \Pr_{r\leftarrow \F^m} \left[\langle Mb, r \rangle =  \langle c , r\rangle \right] 
  \\= \Pr_{r\leftarrow \F^m} \left[ \sum_{i=1}^{m} (Mb - c)_{i} \cdot r_{i} =  0\right] \leq 1 / |\F|,
\end{gather*}
where the last inequality follows from \ref{LINEQ.linear_form_uniform_prob_mul_card_le_one}.
\end{proof}


%%FILE blueprint/src/subsections/lpcp-tensor.tex
For a field \(\F\), vectors \(a \in \F^n\) and \(b \in \F^m\), the tensor product \(a \otimes b \in \F^{n \times m}\)
is the matrix defined by
\[
  (a \otimes b)_{i,j} := a_i b_j \quad \forall i \in [n],\, j \in [m].
\]
We write \(\mathrm{flat}(a \otimes b) \in \F^{nm}\) for the row-concatenation of \(a \otimes b\).

For our purposes, we want a test that checks whether $\mathrm{flat}(a \otimes a) = b$

\begin{definition}\label{TENSORQ}\lean{TENSORQ}\leanok
For a field \(\F\) and \(n \in \N\), define the tensor test language
\[
  \mathrm{TENSORQ}(\F, n) :=
    \bigl\{ (a, b) \in \F^n \times \F^{n^2} \mid b = \mathrm{flat}(a \otimes a) \bigr\}.
\]
\end{definition}

\begin{definition}\label{TENSORQ.verifier}\lean{TENSORQ.verifier}\uses{LPCPVerifier,TENSORQ}\leanok
% Lecture 08 presents this as a subcheck inside the QESAT verifier, where the
% queried proof oracles are already named \(a\) and \(b\). The standalone Lean
% language has explicit input \((a,b)\), so its verifier also checks that the
% proof's two linear parts agree with that input before running the tensor
% multiplicativity check.
Given \((a,b) \in \F^n \times \F^{n^2}\), the LPCP verifier
\(V_{\mathrm{TENSORQ}(\F,n)}^{\langle \pi,\cdot\rangle}(a,b)\) uses a proof
\(\pi=\alpha\parallel\beta \in \F^{n+n^2}\), intended to be the linear encoding of the public input
\((a,b)\). Since the proof may be dishonest, the verifier checks consistency with \((a,b)\) on the sampled
queries before checking the tensor relation:
\begin{enumerate}
  % \item Sample \(s,t \leftarrow \F^n\).
  % \item Query the \(a\)-part of \(\pi\) at \(s\) and at \(t\), obtaining \(\langle a, s \rangle\) and \(\langle a, t \rangle\).
  % \item Query the \(b\)-part of \(\pi\) at \(\mathrm{flat}(s\otimes t)\), obtaining \(\langle b, \mathrm{flat}(s \otimes t) \rangle\).
  % \item Accept if and only if \(\langle b, \mathrm{flat}(s \otimes t) \rangle = \langle a, s \rangle \cdot \langle a, t \rangle\).
  \item Sample \(s,t \leftarrow \F^n\).
  \item Query the \(\alpha\)-part of \(\pi\) at \(s\) and at \(t\), obtaining \(y_s\) and \(y_t\).
  \item Query the \(\beta\)-part of \(\pi\) at \(\mathrm{flat}(s\otimes t)\), obtaining \(y_{s,t}\).
  \item Accept if and only if
  \[
      y_s=\langle a,s\rangle,\qquad
      y_t=\langle a,t\rangle,\qquad
      y_{s,t}=\langle b,\mathrm{flat}(s\otimes t)\rangle,\qquad
      y_{s,t}=y_s y_t.
    \]
\end{enumerate}
\end{definition}

\begin{definition}\label{TENSORQ.selfVerifier}\lean{TENSORQ.selfVerifier}\uses{OracleComp,randOracleSpec}\leanok
In settings where \(a\) and \(b\) are supplied only by the proof oracle, the tensor self-check uses the same three
queries but has no public input to compare against. Given oracle access to
\(\pi=\alpha\parallel\beta\in\F^{n+n^2}\), sample \(s,t\leftarrow\F^n\), query
\[
  y_s=\langle\alpha,s\rangle,\qquad
  y_t=\langle\alpha,t\rangle,\qquad
  y_{s,t}=\langle\beta,\mathrm{flat}(s\otimes t)\rangle,
\]
and accept if and only if \(y_{s,t}=y_s y_t\). This is the tensor component used inside the QESAT verifier.
\end{definition}

\begin{lemma}\label{TENSORQ.tensor_check_complete}\lean{TENSORQ.tensor_check_complete}\leanok
For every \(a,s,t\in\F^n\),
\[
  \langle \mathrm{flat}(a\otimes a),\mathrm{flat}(s\otimes t)\rangle
  = \langle a,s\rangle \langle a,t\rangle.
\]
\end{lemma}
\begin{proof}\leanok
\begin{align*}
  \langle \mathrm{flat}(a\otimes a),\mathrm{flat}(s\otimes t)\rangle
    &= \sum_{i,j} a_i a_j s_i t_j
     = \left(\sum_i a_i s_i\right)\left(\sum_j a_j t_j\right)
     = \langle a,s\rangle \langle a,t\rangle.
\end{align*}
\end{proof}

\begin{lemma}\label{TENSORQ.verifier_soundness_after_sampling}\lean{TENSORQ.verifier_soundness_after_sampling}\uses{TENSORQ.verifier}\leanok
Let \(\F\) be a finite field. If \(b \neq \mathrm{flat}(a\otimes a)\), then for every proof
\(\pi\in \F^{n+n^2}\),
\[
  \Pr_{s,t}\!\left[
    V_{\mathrm{TENSORQ}(\F,n)}^{\langle\pi,\cdot\rangle}(a,b)=1
  \right]
  \leq \frac{2|\F|-1}{|\F|^2}.
\]
\end{lemma}
\begin{proof}\uses{TENSORQ.verifier,LINEQ.linear_form_uniform_prob_mul_card_le_one}\leanok
Suppose \(b \neq \mathrm{flat}(a \otimes a)\), i.e., there exists \(i^*, j^* \in [n]\) such that
\(b_{i^* j^*} \neq a_{i^*} \cdot a_{j^*}\).
For each \(i \in [n]\) define the linear polynomial \(p_i(t) := \sum_{j \in [n]} (b_{ij} - a_i a_j)\, t_j\).
For fixed \(s,t\), if the verifier accepts, then its consistency checks force
\(y_s=\langle a,s\rangle\), \(y_t=\langle a,t\rangle\), and
\(y_{s,t}=\langle b,\mathrm{flat}(s\otimes t)\rangle\). Thus acceptance can only occur when
\[
  \langle b, \mathrm{flat}(s \otimes t) \rangle - \langle a,s\rangle\langle a,t\rangle
  = \sum_{i \in [n]}  \sum_{j \in [n]} (b_{ij} - a_i a_j)\, s_i t_j
  = \sum_{i \in [n]} p_i(t)\, s_i.
\]
Immediate application of Polynomial Identity Testing yields a lower bound of $1 - \tfrac{2}{|\F|}$ but we can do better.
Since \(b_{i^* j^*} \neq a_{i^*} a_{j^*}\), the polynomial \(p_{i^*}\) is nonzero, so by the
Schwartz--Zippel lemma applied to \(t\),
\[
  \Pr_{t}\!\left[p_{i^*}(t) \neq 0\right] \geq 1 - \tfrac{1}{|\F|}.
\]
Conditioned on \(p_{i^*}(t) \neq 0\), the expression \(\sum_i p_i(t)\,s_i\) is a nonzero linear polynomial
in \(s\), so by the Schwartz--Zippel lemma applied to \(s\),
\[
  \Pr_{s \leftarrow \F^n}\!\left[\sum_i p_i(t)\,s_i \neq 0 \mid p_{i^*}(t) \neq 0\right]
  \geq 1 - \tfrac{1}{|\F|}.
\]
Hence
\[
  \Pr_{s,t}\!\left[ V_{\mathrm{TENSORQ}(\F,n)}^{\langle\pi,\cdot\rangle}(a,b)=0 \right]
  \geq \left(1 - \tfrac{1}{|\F|}\right)^{\!2},
\]
and therefore
\[
  \Pr_{s,t}\!\left[ V_{\mathrm{TENSORQ}(\F,n)}^{\langle\pi,\cdot\rangle}(a,b)=1 \right]
  \leq 1 - \left(1 - \tfrac{1}{|\F|}\right)^{\!2}
  = \frac{2|\F|-1}{|\F|^2}.
\]
\end{proof}

\begin{lemma}\label{TENSORQ.selfVerifier_soundness_after_sampling}\lean{TENSORQ.selfVerifier_soundness_after_sampling}\uses{TENSORQ.selfVerifier}\leanok
Let \(\F\) be a finite field and let \(\pi=\alpha\parallel\beta\in\F^{n+n^2}\). If
\((\alpha,\beta)\notin\mathrm{TENSORQ}(\F,n)\), then the tensor self-check accepts with probability at most
\[
  \frac{2|\F|-1}{|\F|^2}.
\]
\end{lemma}
\begin{proof}\uses{TENSORQ.verifier_soundness_after_sampling}\leanok
Apply \ref{TENSORQ.verifier_soundness_after_sampling} to the public pair \((\alpha,\beta)\). The self-check is
exactly the multiplicativity part of that verifier, with \((\alpha,\beta)\) read from the proof oracle itself.
\end{proof}


\begin{theorem}\label{TENSORQ_LPCP}\uses{TENSORQ, LPCP}\lean{TENSORQ_LPCP}\leanok
For every finite field \(\F\) and \(n \in \N\),
\[
  \mathrm{TENSORQ}(\F, n) \in
  \LPCP\!\left[
    \varepsilon_c = 0,
    \varepsilon_s = \frac{2|\F|-1}{|\F|^2},
    \Sigma = \F,
    \ell = n + n^2,
    q = 3,
    r = 2n
  \right].
\]
\end{theorem}
\begin{proof}\uses{TENSORQ.verifier,TENSORQ.tensor_check_complete,TENSORQ.verifier_soundness_after_sampling}\leanok
Use the verifier of \ref{TENSORQ.verifier}. If
\((a,b)\in\mathrm{TENSORQ}(\F,n)\), the honest proof \(\pi=(a,b)\) passes the
three input-consistency checks identically, and the multiplicativity check is
\ref{TENSORQ.tensor_check_complete}. Soundness is
\ref{TENSORQ.verifier_soundness_after_sampling}. The verifier makes three proof
queries and uses two random vectors in \(\F^n\), which gives the displayed query
and randomness bounds.
\end{proof}

\begin{corollary}\label{TENSORQ_LPCP_Zmod2}\uses{TENSORQ, LPCP}\leanok
\[
  \mathrm{TENSORQ}(\F_2, n) \in
  \LPCP\!\left[
    \varepsilon_c = 0,
    \varepsilon_s = 3 / 4,
    \Sigma = \F_2,
    \ell = n + n^2,
    q = 3,
    r = 2 n
  \right].
\]
\end{corollary}
\begin{proof}\uses{TENSORQ_LPCP}\lean{TENSORQ_LPCP_Zmod2}\leanok
Specialize \ref{TENSORQ_LPCP} to \(\F_2\), where \(|\F_2|=2\).
\end{proof}


%%FILE blueprint/src/subsections/lpcp-to-pcp.tex
\begin{lemma}[Chernoff bound]
\label{LPCPToPCP.chernoff_bound}
\lean{LPCPToPCP.chernoff_bound}
\leanok
  Let $X_1,\dots, X_N$ be independent random variables over $\{0,1\}$, and let $\mu=\EE\left[\sum_{i = 1}^N X_i\right]$. For any $\Delta > 0$ one has
  \begin{equation*}
    \Pr\left[\sum_{i = 1}^N X_i \ge \mu+\Delta\right] \le \exp(-2\Delta^2/N)
  \end{equation*}
\end{lemma}
\begin{proof}\leanok
  Write $p_i=\EE[X_i]$, so that $\mu=\sum_{i=1}^N p_i$. We first prove the elementary exponential-moment estimate
  \[
    \EE[\exp(\lambda(X_i-p_i))]\le \exp(\lambda^2/8)
  \]
  for every $\lambda>0$ and every $i$. Since $X_i$ is $\{0,1\}$-valued,
  \[
    \EE[\exp(\lambda(X_i-p_i))]
    = (1-p_i)\exp(-\lambda p_i)+p_i\exp(\lambda(1-p_i)).
  \]
  Equivalently, it is enough to show
  \[
    \log(1-p_i+p_i\exp(\lambda))-p_i\lambda \le \lambda^2/8.
  \]
  For $p\in[0,1]$, set $f(\lambda)=\log(1-p+p\exp(\lambda))-p\lambda$. Then $f(0)=f'(0)=0$, and
  \[
    f''(\lambda)
    = \frac{p\exp(\lambda)}{1-p+p\exp(\lambda)}
      \left(1-\frac{p\exp(\lambda)}{1-p+p\exp(\lambda)}\right)
    \le \frac14.
  \]
  Hence, for $\lambda>0$,
  \[
    f(\lambda)
    = \int_0^\lambda\int_0^u f''(v)\,dv\,du
    \le \int_0^\lambda\int_0^u \frac14\,dv\,du
    = \lambda^2/8.
  \]

  Let $S=\sum_{i=1}^N X_i$. By independence and the estimate above,
  \[
    \EE[\exp(\lambda(S-\mu))]
    = \prod_{i=1}^N \EE[\exp(\lambda(X_i-p_i))]
    \le \exp(N\lambda^2/8).
  \]
  Markov's inequality gives, for every $\lambda>0$,
  \[
    \Pr[S\ge \mu+\Delta]
    = \Pr[\exp(\lambda(S-\mu))\ge \exp(\lambda\Delta)]
    \le \exp(-\lambda\Delta)\EE[\exp(\lambda(S-\mu))]
  \]
  Using the moment bound above, this is at most
  \[
    \exp(-\lambda\Delta+N\lambda^2/8).
  \]
  Taking $\lambda=4\Delta/N$ gives
  \[
    \Pr[S\ge \mu+\Delta]\le \exp(-2\Delta^2/N).
  \]
\end{proof}

\begin{definition}
\label{LPCPToPCP.numShifts}
\uses{LPCPVerifier}
\lean{LPCPToPCP.numShifts}
\leanok
For an LPCP verifier making at most \(q\) linear queries, set
\[
  t := 8\bigl(\lfloor \log_2(100q)\rfloor+1\bigr).
\]
% Lecture 08 only needs \(t=O(\log q)\). The Lean proof fixes this base-two
% value because the formalization uses `Nat.log2`; it is stronger than needed
% for \(q\exp(-t/8)\le 1/100\).
This is the concrete logarithmic shift count used in Lean. The converted PCP
verifier uses \(t\) independent random shifts to decode each linear query by
plurality.
\end{definition}

\begin{definition}[Encoded proof]
\label{LPCPToPCP.encodedProof}
% TODO: define LPCP proofs and add it to \uses
\lean{LPCPToPCP.encodedProof}
\leanok
For a linear LPCP proof \(\pi\in\F_2^\ell\), its PCP encoding \(\widehat{\pi}_{\mathrm{lin}} \in \F_2^{\F_2^\ell}\) is the truth
table
\[
  \widehat{\pi}_{\mathrm{lin}}(a) := \langle \pi,a\rangle,
  \qquad a\in\F_2^\ell.
\]
\end{definition}

\begin{definition}
\label{LPCPToPCP.correctedAnswer}
% TODO: define PCP proofs and add it to \uses
\lean{LPCPToPCP.correctedAnswer,LPCPToPCP.decodedLinearQuery,LPCPToPCP.pluralityZMod2}
\leanok
Given a PCP proof \(\widehat{\pi}:\F_2^\ell\to\F_2\), shifts
\(r_1,\dots,r_t\in\F_2^\ell\), and a linear query \(a\in\F_2^\ell\), define
the locally corrected answer
\[
  \mathsf{Corr}_{\widehat{\pi}}(a;r_1,\dots,r_t)
  :=
  \operatorname{plurality}_{i\in[t]}
  \bigl(\widehat{\pi}(a+r_i)-\widehat{\pi}(r_i)\bigr).
\]
\end{definition}

\begin{definition}
\label{LPCPToPCP.verifier}
\lean{LPCPToPCP.verifier}
\uses{PCPVerifier,LPCPVerifier,BLR.basicVerifier,LPCPToPCP.numShifts,LPCPToPCP.correctedAnswer}
\leanok
Let \(V\) be an LPCP verifier over \(\F_2\). The converted PCP verifier
\(\overline V\), with oracle access to
\(\widehat{\pi}:\F_2^\ell\to\F_2\), proceeds as follows on input \(x\):
\begin{enumerate}
  \item If \(V_{\mathrm{BLR}(\F_{2}, n)}^{\widehat{\pi}}=0\), reject. Otherwise, proceed to the next step.
  \item Sample \(t\) shifts \(r_1,\dots,r_t\in\F_2^\ell\), with \(t\) as in
    \ref{LPCPToPCP.numShifts}.
  \item Simulate \(V(x)\). Whenever \(V\) asks a linear query
    \(a\in\F_2^\ell\), answer it by
    \(\mathsf{Corr}_{\widehat{\pi}}(a;r_1,\dots,r_t)\).
\end{enumerate}
% We define a PCP verifier as follows, for any $\bar{pi} \in \mathbb{F}_2^{\mathbb{F}_2^\ell}$ and $x \in \{0,1\}^*$.
% \vspace{1em}
% \noindent \textbf{Verifier} $\bar{V}^{\bar{\pi}} (x)$:
% \begin{enumerate}
%     \item If $V_{\text{BLR}}^{\bar{\pi}}=0$ output $0$ and halt, otherwise proceed with the next step.
%     \item Independently for each $i \in [t]$, sample $r_i \gets \mathbb{F}_2^\ell$.
%     \item Simulate $V$ by answering each linear query $a$ with $\text{plurality}(\bar{\pi}(a+r_i)-\bar{\pi}(r_i))_{i \in [t]}$.
% \end{enumerate}
\end{definition}


\begin{lemma}[Completeness]
\label{LPCPToPCP.verifier_completeness}
\lean{LPCPToPCP.verifier_completeness}
\uses{LPCPToPCP.verifier,LPCPToPCP.encodedProof}
\leanok
If \(\Pr\!\left[V^{\langle \pi, \cdot \rangle}(x)=1\right]\ge 1-\varepsilon_c\), then \(\Pr\!\left[\overline{V}^{\widehat{\pi}_{\mathrm{lin}}}(x)=1\right]\ge 1-\varepsilon_c\).
\end{lemma}
\begin{proof}\uses{LPCPToPCP.verifier,LPCPToPCP.encodedProof,BLR_basic_completeness_ZMod2}\leanok
The encoded proof is linear, so the BLR test accepts with probability \(1\) by
\ref{BLR_basic_completeness_ZMod2}. For every query \(a\) and every shift
\(r_i\),
\[
  \widehat{\pi}_{\mathrm{lin}}(a+r_i)-\widehat{\pi}_{\mathrm{lin}}(r_i)
  = \langle \pi,a\rangle.
\]
Thus every plurality answer equals the corresponding LPCP oracle answer, and
the simulation has exactly the same acceptance probability as \(V\).
\end{proof}


\begin{lemma}[Soundness, far case]
\label{LPCPToPCP.verifier_soundness_far_case_of_isFar}
\lean{LPCPToPCP.verifier_soundness_far_case_of_isFar}
\uses{LPCPToPCP.verifier,def:distance-to-linearity}
\leanok
If \(\distanceToLin{\widehat{\pi}} \geq  1 / 8\), then
\[
  \Pr[\overline V^{\widehat{\pi}}(x)=1]\leq 7/8 .
\]
\end{lemma}
\begin{proof}\uses{LPCPToPCP.verifier,BLR_basic_soundness_ZMod2}\leanok
By \ref{BLR_basic_soundness_ZMod2}, the BLR precheck accepts with probability at most
\(1-\distanceToLin{\widehat{\pi}}\leq 7/8\). Since \(\overline V\) accepts only when this test accepts, its
acceptance probability is at most \(7/8\).
\end{proof}

\begin{lemma}
\label{LPCPToPCP.local_correction}
\lean{LPCPToPCP.localCorrectionBad_single_le_quarter_of_not_far,LPCPToPCP.correctedAnswer_failure_le_exp_of_single,LPCPToPCP.chernoff_half_of_mean_le_quarter,LPCPToPCP.numShifts_eta_bound}
\uses{LPCPToPCP.correctedAnswer,def:relative-hamming-distance}
\leanok
If \(\distance{\widehat{\pi}}{f} < 1 / 8\) for some linear
function \(f\) of the form \(u\mapsto\langle\pi,u\rangle\), then
\[
  \Pr_{r_1,\dots,r_t}\!\left[
    \mathsf{Corr}_{\widehat{\pi}}(a;r_1,\dots,r_t)\neq \langle\pi,a\rangle
  \right]
  \leq \exp(-t/8), \quad \forall a\in\F_2^\ell.
\]
\end{lemma}
\begin{proof}\uses{LPCPToPCP.correctedAnswer,LPCPToPCP.chernoff_bound}\leanok
For one shift \(r\), the corrected sample can be wrong only if
\(\widehat{\pi}(a+r)\) or \(\widehat{\pi}(r)\) disagrees \(f\). A union bound gives error probability at most
\(2\cdot(1/8)=1/4\). If \(X_i\) is the indicator of failure for shift \(r_i\),
then \(\EE[\sum_i X_i]\le t/4\). The plurality answer is wrong only if
\(\sum_i X_i\ge t/2\), so \ref{LPCPToPCP.chernoff_bound} with
\(\Delta=t/4\) gives \(\exp(-t/8)\).
% Lecture 07 states a looser concentration estimate with different constants.
% The blueprint follows the Hoeffding normalization used by the Lean lemma
% `chernoff_half_of_mean_le_quarter`, which is mathematically stronger here.
\end{proof}

\begin{lemma}
\label{LPCPToPCP.corrected_simulation_hybrid}
\lean{LPCPToPCP.corrected_run_accept_le_linear_add,LPCPToPCP.verifier_accept_le_linear_add_of_query_correction}
\uses{LPCPToPCP.verifier}
\leanok
Assume every possible linear query \(a\) is locally corrected incorrectly with
probability at most \(\eta\). If the simulated LPCP computation makes at most
\(q\) proof queries, then its acceptance probability under the corrected PCP
oracle is at most the acceptance probability under the nearby linear proof
plus \(q\eta\).
\end{lemma}
\begin{proof}\leanok
Expose the LPCP computation one oracle query at a time. At a proof query, the
two simulations can diverge only when the local correction for the queried
linear form is wrong; by assumption this has probability at most \(\eta\).
Inducting over the query-bound derivation gives an additive loss of at most
\(q\eta\).
% Lecture 08 presents this step as a union bound over the verifier's \(q\)
% queries. The Lean proof phrases the same estimate as a sequential hybrid,
% which also covers adaptive LPCP verifiers whose later queries may depend on
% earlier corrected answers.
\end{proof}


\begin{lemma}[Soundness, near case]
\label{LPCPToPCP.verifier_soundness_near_case}
\lean{LPCPToPCP.verifier_soundness_near_case_of_single_query_bound}
\uses{LPCPToPCP.verifier,LPCP,def:relative-hamming-distance}
\leanok
If \(x\notin L\), \(V\) has LPCP soundness \(\varepsilon_s\), and
\(\distance{\widehat{\pi}}{f} < 1/8\) for some linear function \(f\) of the form \(u\mapsto\langle\pi,u\rangle\), then
\[
  \Pr[\overline V^{\widehat{\pi}}(x)=1]\leq \varepsilon_s+\frac{1}{100}.
\]
\end{lemma}
\begin{proof}\uses{LPCPToPCP.local_correction,LPCPToPCP.corrected_simulation_hybrid,LPCPToPCP.numShifts,LPCP}\leanok
By \ref{LPCPToPCP.local_correction}, every possible linear query has local
correction error at most \(\eta=\exp(-t/8)\). The hybrid estimate
\ref{LPCPToPCP.corrected_simulation_hybrid} compares the corrected simulation
with the simulation using the nearby linear proof and loses at most
\(q\eta\). The choice of \(t\) in \ref{LPCPToPCP.numShifts} gives
\(q\eta\le 1/100\). The nearby linear proof is a valid LPCP proof, so its
acceptance probability is at most \(\varepsilon_s\) because \(x\notin L\).

  % \paragraph{Case $\delta(\hat{\pi}, \text{LIN}) < 1/8$}
  % Let $S = \{\tilde{\pi} \in \lin: \, \delta(\hat{\pi},\tilde{\pi} \le \delta(\hat{\pi}, \text{LIN}))\}$.
  % Suppose for a contradiction that $|S| \ge 2$, and let $\bar{\pi},\tilde{\pi} \in S$ be distinct: since 
  % $\delta(\hat{\pi},\bar{\pi} \le \delta(\hat{\pi}, \text{LIN})) < 1/8$ and $\delta(\hat{\pi},\tilde{\pi} \le \delta(\hat{\pi}, \text{LIN})) < 1/8$,
  % by triangle inequality we have $\delta(\bar{\pi},\tilde{\pi}) < 1/4$;
  % since $|\bbF_2|=2$, this contradicts the fact that any two distinct $\bar{\pi},\tilde{\pi} \in \lin$ satisfy $\delta(\bar{\pi},\tilde{\pi}) \ge 1/2$.

  % Let then $\tilde{\pi}$ be the unique element of $\lin$ with minimal distance to $\hat{\pi}$, and let $z \in \bbF^\ell$ the unique vector such that $\tilde{\pi}(x)=\langle z, x \rangle$ for all $x$.
  % If for every query $a$ that $V$ asks we have $\text{plurality}(\hat{\pi}(a+r_i)-\hat{\pi}(r_i))_{i \in [t]}=\tilde{\pi}(a)$, then $\bar{V}^{\hat{\pi}}(x)=V^{\langle z, \cdot \rangle}(x)$, so
  % provided that the $r_i$ satisfy this condition we have that $\Pr[\bar{V}^{\hat{\pi}}(x)=1]=\Pr[V^{\langle z, \cdot \rangle}(x)=1] \le \varepsilon_s$.
  % We are only left to bound the probability that for each of the $q$ queries that $V$ asks we have $\text{plurality}(\hat{\pi}(a+r_i)-\hat{\pi}(r_i))_{i \in [t]}=\tilde{\pi}(a)$.
  % Note that for any $a$ and $r_i$, if $\hat{\pi}(a+r_i)=\tilde{\pi}(a+r_i)$ and $\hat{\pi}(r_i)=\tilde{\pi}(r_i)$ then $\hat{\pi}(a+r_i)-\hat{\pi}(r_i)=\tilde{\pi}(a)$.
  % Hence, the proabability over one $r_i$ that $\hat{\pi}(a+r_i)-\hat{\pi}(r_i) \neq \tilde{\pi}(a)$ is at most $2\cdot \delta(\hat{\pi},\tilde{\pi})\le 1/4$ by a union bound.
  % Also note that $\text{plurality}(\hat{\pi}(a+r_i)-\hat{\pi}(r_i))_{i \in [t]}\neq \tilde{\pi}(a)$ means that $\hat{\pi}(a+r_i)-\hat{\pi}(r_i) \neq \tilde{\pi}(a)$ for at least $\lceil t/2 \rceil \ge t/2$ of the $i$'s.
  % Let $X_i = \ind[\hat{\pi}(a+r_i)-\hat{\pi}(r_i) \neq \tilde{\pi}(a)]$. Then $\mu = \EE[\sum_{i \in [t]}X_i] \le t/4$. Applying Chernoff bound with $\Delta = t/4$, the probability over the choice of $r_1,\dots, r_t$ that there are at least $t/2$ $i$'s with $X_i=1$ is at most $\exp(-t/8)$.
  % Union bounding over the $q$ queries that $V$ asks, we have that the local correction fails with probability at most $q\cdot \exp(-t/8) \le 1/100$ by our choice of $t$.

\end{proof}


\begin{lemma}[Soundness]
\label{LPCPToPCP.verifier_soundness_of_cases}
\lean{LPCPToPCP.verifier_soundness_of_cases}
\uses{LPCPToPCP.verifier}
\leanok
For every \(x\notin L\) and every PCP proof \(\widehat{\pi}\),
\[
  \Pr[\overline V^{\widehat{\pi}}(x)=1]
  \leq \max\{7/8,\varepsilon_s+1/100\}.
\]
\end{lemma}
\begin{proof}\uses{LPCPToPCP.verifier_soundness_far_case_of_isFar,LPCPToPCP.verifier_soundness_near_case}\leanok
If \(\distanceToLin{\widehat{\pi}} \geq 1 / 8\), apply \ref{LPCPToPCP.verifier_soundness_far_case_of_isFar}. Otherwise, apply
\ref{LPCPToPCP.verifier_soundness_near_case}.
\end{proof}

\begin{lemma}[LPCPs admit exponential-size PCPs]
\uses{LPCP,PCP,LPCPToPCP.numShifts}
\label{LPCP_to_PCP_ZMod2}
\lean{LPCP_to_PCP_ZMod2}
\leanok
\begin{gather*}
\LPCP[\varepsilon_c,\varepsilon_s,\Sigma=\F_{2},\ell,q,r] \subseteq\\
\PCP\left[
\varepsilon_c,\
\varepsilon_s'=\max\{7/8,\varepsilon_s+1/100\},
\Sigma'=\F_{2},\
\ell'=2^{\ell},\
q'=3+2tq,\
r'=r+(2+t)\ell
\right].
\end{gather*}
\end{lemma}
\begin{proof}\uses{LPCPToPCP.verifier,LPCPToPCP.verifier_completeness,LPCPToPCP.verifier_soundness_of_cases}\leanok
Let $L \in \LPCP[\varepsilon_c,\varepsilon_s,\Sigma=\F_{2},\ell,q,r]$ and let $V$ be the LPCP verifier for $L$. Use the converted PCP verifier $\overline V$ defined in \ref{LPCPToPCP.verifier}.
Completeness is
\ref{LPCPToPCP.verifier_completeness}, and soundness is
\ref{LPCPToPCP.verifier_soundness_of_cases}. The BLR test makes three proof
queries. Each simulated LPCP query uses two PCP queries for each of the \(t\)
shifts, so the query bound is \(3+2tq\). The randomness consists of the
original \(r\) random bits, \(2\ell\) bits for the BLR test, and \(t\ell\) bits
for the shifts.
\end{proof}


%%FILE blueprint/src/subsections/pcp.tex
\nocite{Chiesa2024-fopp-course}
\begin{definition}
\label{QESAT}
\lean{QESAT}
\leanok
Let \(m,n >0\) be integers, \(\F\) be a field, \(R = \F[x_{1},\dots,x_{n}]\) be its polynomial ring and \(R_{d}\) denote the subspace of polynomials of degree \(\leq d\).

\[
  \QESAT(\F,n) :=
  \{ (p_1, \ldots, p_m) \in R^m_{2} \mid
    \exists a \in \F^n, \,
    \forall i \in [m], \,
    p_i(a_1, \ldots, a_n) = 0 \}.
\]
\end{definition}

For example, $(x + 1, xy + z) \in \QESAT(\F_2,3)$.
The goal of this chapter is to construct an exponential-length PCP for
\(\QESAT(\F_2,n)\) (Theorem \ref{QESAT_exp_PCP}). Since \(\QESAT(\F_2,n)\) is NP-complete, this will imply the existence of exponential-length PCPs for all languages in NP.

\section{Oracle computations and verifiers}
Instead of oracle Turing machines, we will define verifiers using Lean monads, specifically the \texttt{OracleComp} free monad from VCV-io.

\begin{definition}[Oracle specification]
\label{OracleSpec}
\lean{OracleSpec}
\leanok
An \emph{oracle specification} defines the type of the queries and outputs that an oracle can handle.
\end{definition}

\begin{definition}
\label{randOracleSpec}
\lean{randOracleSpec}
\uses{OracleSpec}
\leanok
For this section, the \emph{randomness oracle specification} will be of a monadic function that takes
nothing and returns an element from a field \(\F\).
\end{definition}

\begin{definition}[Query implementation]
\label{QueryImpl}
\lean{QueryImpl}
\leanok
A \emph{query implementation} is a monadic function whose behavior conforms to a given oracle specification.
\end{definition}

\begin{definition}
\label{uniformSample}
\uses{QueryImpl, randOracleSpec}
\lean{uniformSample}
\leanok
The \emph{uniform sampling query implementation} for a field \(\F\) is a (mathematical) 
query implementation for the randomness oracle specification that samples uniformly from \(\F\).
\end{definition}
This is the the oracle implementation that will be used to prove theorems about our verifiers,
since it is compatible with the more classical notion of probability distributions used in the previous chapters.
\begin{remark}[Implementing verifiers]
If we want to actually run the verifiers, we can also use concrete implementations based on the IO monad
and system randomness for pseudorandom sampling.
\end{remark}
\begin{definition}[Oracle context]
\label{OracleContext}
\lean{OracleContext}
\uses{OracleSpec, QueryImpl}
\leanok
An \emph{oracle context} is a bundle consisting of an oracle specification and a query implementation for that specification.
\end{definition}
\begin{definition}[Randomness oracle]
\label{randOracle}
\uses{OracleContext, randOracleSpec, uniformSample}
\lean{randOracle} \leanok
As \emph{randomness oracle}, we will use the {uniform sampling query implementation} for the randomness oracle specification.
\end{definition}
\begin{definition}[Oracle computation]
\label{OracleComp}
\uses{OracleSpec}
\lean{OracleComp}\leanok
An \emph{oracle computation} is a monadic computation in the free monad generated by an oracle specification,
i.e. a description (a ``script'') of a program consisting of a chain of oracle queries (compatible with any query implementation of the specification)
and deterministic Lean computations.
\end{definition}

A \emph{verifier} is a function that takes an input \(x \in \{0,1\}^*\) and produces an oracle computation
(with query access to a randomness oracle) that outputs a boolean. This definition can be extended to include
other oracles (typically a proof oracle). In Lean this convention is represented directly by function types returning
\texttt{OracleComp} computations, rather than by a standalone verifier declaration.

We also need to define verifiers that do not take any inputs (for example the BLR verifier).
Since they use the same randomness oracle specification as the normal verifiers, they can
compose nicely with them.
A \emph{unit verifier} (or just \emph{verifier} when it is clear from the context) is an oracle computation with query access to a randomness oracle that outputs a boolean.
This definition can be extended to include
other oracles (typically a proof oracle).

\begin{remark}[Time complexity]
The price to pay for using these definitions for verifiers instead of Turing machines is that we have no simple way to enforce
the time complexity of the verifier. For that, a possible approach would be to restrict the allowed operations to a strict set
of operations instead of general Lean computations, and track the time cost using query complexity.
\end{remark}
\section{PCP and LPCP definitions}

\begin{definition}
\label{PCPVerifier}
\lean{PCPVerifier}
\uses{OracleComp,randOracleSpec}
\leanok
A \emph{PCP verifier} is a verifier with query access to a function $\pi : [\ell (|x|)] \to \Sigma$.
\end{definition}

\begin{definition}
\uses{PCPVerifier, randOracle}
\label{PCP}
\lean{PCP}
\leanok
A language $L \subseteq \{0,1\}^*$ is in
$\PCP[\varepsilon_c, \varepsilon_s, \Sigma, \ell, q, r]$
if there exists a polynomial-time PCP verifier $V$
such that for every $x \in \{0,1\}^*$,
$V$ makes at most $q(|x|)$ queries to the proof oracle and
at most $r(|x|)$ randomness queries, and the following holds:
\begin{itemize}
  \item Completeness: If $x \in L$, then
  $\exists \pi \in \Sigma^{\ell(|x|)}, \,
  \Pr\left[V^\pi(x)=1 \right] \geq 1 - \varepsilon_c$.
  \item Soundness: If $x \notin L$, then
  $\forall \widetilde{\pi} \in \Sigma^{\ell(|x|)},\,
  \Pr\left[V^{\widetilde{\pi}}(x)=1 \right] \leq \varepsilon_s$.
\end{itemize}
\end{definition}

\begin{definition}
\label{LPCPVerifier}
\lean{LPCPVerifier}
\uses{OracleComp,randOracleSpec}
\leanok
A \emph{LPCP verifier} is a verifier with query access to a linear function
$\langle \pi,\cdot \rangle : \Sigma^{\ell(|x|)} \to \Sigma$.
\end{definition}

\begin{definition}
\uses{LPCPVerifier, randOracle}
\label{LPCP}
\lean{LPCP}
\leanok
A language $L \subseteq \{0,1\}^*$ is in
$\LPCP[\varepsilon_c, \varepsilon_s, \Sigma, \ell, q, r]$
if there exists a polynomial-time LPCP verifier $V$
such that for every $x \in \{0,1\}^*$,
$V$ makes at most $q(|x|)$ queries to the linear proof oracle and
at most $r(|x|)$ randomness queries, and the following holds:
\begin{itemize}
  \item Completeness: If $x \in L$, then
  $\exists \pi \in \Sigma^{\ell(|x|)}, \,
  \Pr\left[V^{\langle \pi, \cdot \rangle}(x)=1 \right]
  \geq 1 - \varepsilon_c$.
  \item Soundness: If $x \notin L$, then
  $\forall \widetilde{\pi} \in \Sigma^{\ell(|x|)},\,
  \Pr\left[V^{\langle \widetilde{\pi}, \cdot \rangle}(x)=1 \right]
  \leq \varepsilon_s$.
\end{itemize}
\end{definition}


%%FILE blueprint/src/subsections/qesat-exp-pcp.tex
% \begin{theorem}\label{old-QESAT_poly_LPCP}\lean{QESAT_poly_LPCP}\uses{QESAT, LPCP}\leanok
% \[
%   \mathrm{QESAT}(\F_2,n) \in
%   \mathrm{LPCP}[
%     \varepsilon_c = 0,
%     \varepsilon_s = 3/4,
%     \Sigma = \F_2,
%     \ell = n + n^2,
%     q = 4,
%     r = |x| + 2*n
%   ].
% \]
% \end{theorem}
% \begin{proof}%\uses{} tensor stuff

% \noindent Let $(p_1, \dots, p_m)$ be an instance of $\text{QESAT}(\mathbb{F})$ with $n$ variables. \\
% The LPCP verifier expects a proof $\pi = (a, b) \in \mathbb{F}^{n+n^2}$ and works as follows:

% \vspace{1em}
% \noindent \textbf{Verifier} $V^{(a,b)}(p_1, \dots, p_m)$:
% \begin{enumerate}
%     \item Sample $r \gets \mathbb{F}^m$ and $s, t \gets \mathbb{F}^n$.
%     \item Define 
%     \[
%     M := \begin{bmatrix} \text{coeff}(p_1) \\ \vdots \\ \text{coeff}(p_m) \end{bmatrix} \in \mathbb{F}^{m \times (n+n^2)} 
%     \quad \text{and} \quad 
%     c := \begin{bmatrix} -\text{const}(p_1) \\ \vdots \\ -\text{const}(p_m) \end{bmatrix} \in \mathbb{F}^m.
%     \]
%     \textit{(Note: $\text{coeff}(p_i) :=$ non-constant coeffs of $p_i$, and $\text{const}(p_i) :=$ constant coeff of $p_i$)}
    
%     \item \textit{(linear equation test)} $\rightarrow$ Query $(a,b)$ at $M^T r$ and check that $\langle a \parallel b, M^T r \rangle = \langle c, r \rangle$.
    
%     \item \textit{(tensor test)} $\rightarrow$ Query $b$ at $\text{flat}(s \otimes t)$, $a$ at $s$ and $t$, and check that 
%     \[
%     \langle b, \text{flat}(s \otimes t) \rangle = \langle a, s \rangle \cdot \langle a, t \rangle.
%     \]
% \end{enumerate}

% \vspace{1em}
% \noindent \textbf{Completeness:} 
% Suppose $p_1(a) = \dots = p_m(a) = 0$ and set $b := \text{flat}(a \otimes a)$. \\
% Then 
% \[
% \Pr_{s,t} \left[ \langle b, \text{flat}(s \otimes t) \rangle = \langle a, s \rangle \cdot \langle a, t \rangle \right] = 1
% \]
% and 
% \[
% \Pr_r \left[ \langle a \parallel b, M^T r \rangle = \langle c, r \rangle \right] = 1 \quad \left(\text{since } M \begin{bmatrix} a \\ \text{flat}(a \otimes a) \end{bmatrix} = c \right).
% \]

% \vspace{1em}
% \noindent \textbf{Soundness:}
% Suppose $\forall a \in \mathbb{F}^n \ \exists i \in [m] \text{ s.t. } p_i(a) \neq 0$. Fix any $\pi = (a, b)$. \\
% Either:
% \begin{enumerate}
%     \item[(i)] $b \neq \text{flat}(a \otimes a) \implies$ tensor test passes w.p. $\leq \frac{2|\mathbb{F}|-1}{|\mathbb{F}|^2}$
%     \item[OR (ii)] $b = \text{flat}(a \otimes a)$ and $M \begin{bmatrix} a \\ b \end{bmatrix} \neq c \implies$ linear equation test passes w.p. $\leq \frac{1}{|\mathbb{F}|}$
% \end{enumerate}

% \end{proof}


Let \(x=(p_1,\dots,p_m)\) be a \(\QESAT(\F_2,n)\) instance. We use the following
linearization. For each quadratic
polynomial \(p_i\), form a row \(M_i\in\F_2^{n+n^2}\) by replacing every
linear monomial \(X_j\) by the variable \(a_j\) and every quadratic monomial
\(X_jX_k\) by the tensor variable \(b_{jk}\). Let
    \[
    M_{x} := \begin{bmatrix} M_1 \\ \vdots \\ M_m \end{bmatrix} \in \F_2^{m \times (n+n^2)} 
    \quad \text{and} \quad 
    c_{x} := \begin{bmatrix} -\operatorname{const}(p_1) \\ \vdots \\ -\operatorname{const}(p_m) \end{bmatrix} \in \F_2^m,
    \]
where \(\operatorname{const}(p_i)\) is the constant term of \(p_i\).
Then \(M_{x}(a\parallel b)=c_{x}\) expresses the original equations after
replacing \(a_ja_k\) by \(b_{jk}\).

\begin{definition}
\label{QESAT.verifier}
\lean{QESAT.verifier}
\uses{LPCPVerifier,QESAT,LINEQ.verifier,TENSORQ.selfVerifier}
\leanok
Let \(x=(p_1,\dots,p_m) \in R_{2}^{m}\) be a QESAT instance as defined in \ref{QESAT}. Let \(V_{\QESAT(\F_2,n)}^{\pi}(x)\) be the LPCP verifier 
with oracle access to \(\pi=a\parallel b\in\F_2^{n+n^2}\), defined as follows:
\begin{enumerate}
  \item Sample $r \gets \F_{2}^m$ and $s, t \gets \F_{2}^n$.
  \item Construct \((M_{x},c_{x}) \in \F_2^{m \times (n+n^2)} \times \F_2^m\) by the linearization above.
  \item Check that \(V_{\LINEQ(\F_2,m,n+n^2)}^{\langle \pi, \cdot \rangle } (M_{x},c_{x}) = 1\).
  \item Query \(\pi\) at
  \((s,0)\), \((t,0)\), and \((0,\mathrm{flat}(s\otimes t))\) and check that
  \[
    \langle b,\mathrm{flat}(s\otimes t)\rangle
    =\langle a,s\rangle\langle a,t\rangle .
  \]
  Equivalently, this last step is the tensor self-check of \ref{TENSORQ.selfVerifier} specialized to \(\F_2\).
\end{enumerate}
\end{definition}

\begin{lemma}[Completeness]
\label{QESAT.verifier_completeness}
\lean{QESAT.verifier_correct}
\uses{QESAT,QESAT.verifier}
\leanok
If \(x\in \QESAT(\F_2,n)\), then exists a proof
\(\pi\in\F_2^{n+n^2}\) such that \[\Pr\!\left[V_{\QESAT(\F_2,n)}^{\pi}(x)=1\right]=1.\]
\end{lemma}
\begin{proof}\uses{QESAT.verifier,LINEQ.verifier,TENSORQ.tensor_check_complete}\leanok
Let \(a\in\F_2^n\) satisfy all equations and set
\(b=\mathrm{flat}(a\otimes a)\). By construction of the linearized system,
\(M_x(a\parallel b)=c_x\), so the LINEQ check
accepts for every \(r\). The tensor self-check accepts for every \(s,t\) by
\ref{TENSORQ.tensor_check_complete}.
\end{proof}

\begin{lemma}[Soundness]
\label{QESAT.verifier_soundness}
\lean{QESAT.verifier_correct}
\uses{QESAT,QESAT.verifier}
\leanok
If \(x\notin \QESAT(\F_2,n)\), then for every proof
\(\tilde\pi\in\F_2^{n+n^2}\),
\[\Pr\!\left[V_{\QESAT(\F_2,n)}^{\tilde\pi}(x)=1\right]\leq 3 / 4.\]
\end{lemma}
\begin{proof}\uses{QESAT.verifier,LINEQ.linear_form_uniform_prob_mul_card_le_one,TENSORQ.selfVerifier_soundness_after_sampling}\leanok
If some \(p_i\) has total degree \(>2\), the verifier rejects immediately, so
the acceptance probability is \(0\). Assume now that every \(p_i\) is quadratic
and write \(\pi=a\parallel b\). If \(b\neq\mathrm{flat}(a\otimes a)\), then
the tensor self-check accepts with probability at most \(3/4\) by
\ref{TENSORQ.selfVerifier_soundness_after_sampling} specialized to \(\F_2\), and
the full verifier accepts only if that check accepts.

Otherwise \(b=\mathrm{flat}(a\otimes a)\). Since \(x\notin\QESAT(\F_2,n)\),
the assignment \(a\) does not satisfy all equations. By construction of the
linearized system, \(M_x(a\parallel b)\neq c_x\). The LINEQ check
therefore accepts with probability at most \(1/|\F_2|=1/2\) by
\ref{LINEQ.linear_form_uniform_prob_mul_card_le_one}. Again the full verifier
accepts only if this check accepts, so the acceptance probability is at most
\(1/2\leq 3/4\).
\end{proof}

\begin{theorem}
\label{QESAT_poly_LPCP}
\lean{QESAT_poly_LPCP}
\uses{QESAT,LPCP}
\leanok
\[
  \mathrm{QESAT}(\F_2,n) \in
  \mathrm{LPCP}[
    \varepsilon_c = 0,
    \varepsilon_s = 3/4,
    \Sigma = \F_2,
    \ell = n + n^2,
    q = 4,
    r = |x| + 2n
  ].
\]
\end{theorem}
\begin{proof}\uses{QESAT.verifier,QESAT.verifier_completeness,QESAT.verifier_soundness}\leanok
% Lecture 08 gives the exact LINEQ randomness as \(m\log|\F|\). In Lean,
% \(|x|\) is the encoded input size and bounds the number of equations \(m\);
% after specializing to \(\F_2\), this yields the displayed \(|x|+2n\) bound.
Use the verifier \ref{QESAT.verifier}. It makes one query for the LINEQ check
and three queries for the tensor self-check, for a total of four. Its
randomness is the LINEQ randomness plus two vectors in \(\F_2^n\). Completeness
and soundness are \ref{QESAT.verifier_completeness} and
\ref{QESAT.verifier_soundness}.
\end{proof}

\begin{theorem}
\label{QESAT_exp_PCP_before_repetition}
\lean{QESAT_exp_PCP_before_repetition}
\uses{QESAT,PCP}
\leanok
\[
  \mathrm{QESAT}(\F_2,n) \in
  \mathrm{PCP}[
    \varepsilon_c = 0,
    \varepsilon_s = 7/8,
    \Sigma = \F_2,
    \ell = 2^{|x|},
    q = O(1),
    r = |x|^{O(1)}
  ].
\]
\end{theorem}
\begin{proof}\uses{QESAT_poly_LPCP,LPCP_to_PCP_ZMod2}\leanok
% The lecture's final no-abstraction slide shares local-correction work in a
% direct QESAT verifier and therefore displays different constants. This
% formalization deliberately routes through the modular LPCP-to-PCP theorem,
% so only the asymptotic constant-query statement is retained here.
Apply \ref{LPCP_to_PCP_ZMod2} to \ref{QESAT_poly_LPCP}. Since the LPCP
soundness is \(3/4\), the converted soundness is
\[
  \max\{7/8,\,3/4+1/100\}=7/8.
\]
The query bound remains constant because \(q=4\), and the randomness is
polynomial in the input size. The Lean formalization pads the proof oracle from
length \(2^{n+n^2}\) to length \(2^{|x|}\).
\end{proof}

\begin{lemma}
\label{QESAT.constant_repetition}
\lean{QESAT_exp_PCP}
\uses{PCP}
\leanok
\[\PCP[0,7/8,\F_2,2^{|x|},O(1),|x|^{O(1)}] \subseteq \PCP[0,1/2,\F_2,2^{|x|},O(1),|x|^{O(1)}].\]
\end{lemma}
\begin{proof}\uses{QESAT_exp_PCP_before_repetition}\leanok
% This repetition is folded into the Lean theorem `QESAT_exp_PCP` rather than
% exposed as a standalone declaration.
Invoke \ref{QESAT_exp_PCP_before_repetition} to obtain the \(7/8\)-sound
verifier. Run it independently six times and accept only if all six runs
accept. Completeness error remains \(0\). For \(x\notin L\), the acceptance
probability is at most \((7/8)^6\leq 1/2\). Query and randomness bounds are
multiplied by the constant \(6\).
\end{proof}

\begin{theorem}
\uses{QESAT,PCP}
\label{QESAT_exp_PCP}
\lean{QESAT_exp_PCP}
\leanok
\[
  \mathrm{QESAT}(\F_2,n) \in
  \mathrm{PCP}[
    \varepsilon_c = 0,
    \varepsilon_s = 1/2,
    \Sigma = \F_2,
    \ell = 2^{|x|},
    q = O(1),
    r = |x|^{O(1)}
  ].
\]
\end{theorem}
\begin{proof}\uses{QESAT_exp_PCP_before_repetition,QESAT.constant_repetition}\leanok
Combine \ref{QESAT_exp_PCP_before_repetition} with the constant repetition
step \ref{QESAT.constant_repetition}.
\end{proof}

